\documentclass[12pt,spanish,fleqn,openany,letterpaper,pagesize]{scrbook}

% ==============================================================================
% CODIFICACIÓN Y LENGUAJE
% ==============================================================================
\usepackage[utf8]{inputenc}
\usepackage[T1]{fontenc}
\usepackage[spanish]{babel}

% ==============================================================================
% GEOMETRÍA DE PÁGINA
% ==============================================================================
\usepackage[
  letterpaper,
  top=4cm,
  bottom=4cm,
  inner=3.5cm,
  outer=2.5cm,
  headheight=15pt
]{geometry}

% ==============================================================================
% PAQUETES MATEMÁTICOS Y SÍMBOLOS
% ==============================================================================
\usepackage{amsmath}
\usepackage{amsthm}
\usepackage{amssymb}
\usepackage{amsfonts}
\usepackage{mathtools}
\usepackage{amscd}
\usepackage{bm}
\usepackage{dsfont}
\usepackage{pifont}
\usepackage{tipa}
\usepackage{comment}
\usepackage{mathtools}

% ==============================================================================
% TABLAS Y FIGURAS
% ==============================================================================
\usepackage{graphicx}
\usepackage{epsfig}
\usepackage{epic}
\usepackage{eepic}
\usepackage{here}
\usepackage{lscape}
\usepackage{wrapfig}
\usepackage{subfigure}
\usepackage{threeparttable}
\usepackage{tabularx}
\usepackage{longtable} % Se carga solo una vez aquí
\usepackage{caption}   % Se carga después de longtable
\usepackage{cellspace}
\setlength\cellspacetoplimit{8pt}    % Margen superior dentro de la celda
\setlength\cellspacebottomlimit{8pt} % Margen inferior dentro de la celda

% ==============================================================================
% TIKZ Y GRÁFICOS
% ==============================================================================
\usepackage{tikz-cd}
\usepackage{tikz}
\usetikzlibrary{decorations.pathmorphing, decorations.markings, fit}
\usetikzlibrary{arrows.meta}
\usetikzlibrary{babel}
\usetikzlibrary{shapes.geometric}
\usetikzlibrary{backgrounds}

% ==============================================================================
% OTROS PAQUETES DE FORMATO Y TEXTO
% ==============================================================================
\usepackage{fancyhdr}
\usepackage{lipsum}
\usepackage{bold-extra}
\usepackage{needspace}
\usepackage[numbers]{natbib}

% ==============================================================================
% ENTORNO DE TEOREMAS Y CONTADORES
% ==============================================================================
\counterwithin{equation}{section}
\theoremstyle{definition}
\newtheorem{Tma}[equation]{Teorema}
\newtheorem{Prop}[equation]{Proposición}
\newtheorem{Def}[equation]{Definición}
\newtheorem{Lema}[equation]{Lema}
\newtheorem{Cor}[equation]{Corolario}
\newtheorem{Ejc}[equation]{Ejercicio}
\newtheorem{Ejm}[equation]{Ejemplo}
\newtheorem{Not}[equation]{Notación}
\newtheorem{Obs}[equation]{Observación}

\counterwithin{figure}{section}
%\counterwithin{table}{section}

\renewcommand*{\proofname}{Demostración}

% ==============================================================================
% FORMATO DE CAPÍTULO PERSONALIZADO (KOMA-SCRIPT / scrbook)
% ==============================================================================
\addtokomafont{chapter}{\rmfamily\bfseries}
\addtokomafont{chapterprefix}{\rmfamily\mdseries\Large}
\addtokomafont{section}{\rmfamily\bfseries}
\addtokomafont{subsection}{\rmfamily\bfseries}
\addtokomafont{subsubsection}{\rmfamily\bfseries}
\addtokomafont{paragraph}{\rmfamily\bfseries}
\addtokomafont{subparagraph}{\rmfamily\bfseries}

\KOMAoptions{chapterprefix=true}

\renewcommand*{\chapterformat}{%
  \raggedleft\MakeUppercase{\chaptername}\ \Huge\thechapter\par\vspace{1ex}%
}

\renewcommand*{\chapterlineswithprefixformat}[3]{%
  \raggedleft
  #2% Imprime "CAPÍTULO X"
  \rule{\textwidth}{0.8pt}\par\nobreak
  \vspace{0.8ex}%
  {\raggedleft #3\par}% Imprime el Título alineado a la derecha
  \vspace{0.5ex}%
  \rule{\textwidth}{0.8pt}%
}

% Nombres en español
\renewcaptionname{spanish}{\tablename}{Tabla}
\renewcaptionname{spanish}{\figurename}{Figura}
\renewcommand{\listtablename}{Lista de Tablas}
\renewcommand{\listfigurename}{Lista de Figuras}
\renewcommand{\contentsname}{Contenido}

% ==============================================================================
% ENCABEZADOS Y PIE DE PÁGINA
% ==============================================================================
\pagestyle{fancyplain}

\renewcommand{\chaptermark}[1]{\markboth{\footnotesize\MakeUppercase{\chaptername\ \thechapter.\ #1}}{}}
\renewcommand{\sectionmark}[1]{\markright{\footnotesize\MakeUppercase{\thesection\quad #1}}}

\lhead[\fancyplain{}{\normalsize\thepage}]{\fancyplain{}{\rightmark}}
\rhead[\fancyplain{}{\leftmark}]{\fancyplain{}{\normalsize\thepage}}

\fancyfoot{}
\fancyfoot[C]{\normalsize\thepage}
\thispagestyle{fancy}

\unitlength1mm
\mathindent0cm
\marginparwidth0cm
\parindent0cm

\newcommand{\PreserveBackslash}[1]{\let\temp=\\#1\let\\=\temp}
\let\PBS=\PreserveBackslash

\renewcommand{\baselinestretch}{1.1}

\newcommand{\arr}[1]{\raisebox{1.5ex}[0cm][0cm]{#1}}
\newcommand{\diagscale}{0.3\textwidth}%
\newcommand{\subseteqab}{\stackrel{ab}{\subseteq}}
\newcommand{\subseteqcer}{\stackrel{cer}{\subseteq}}
\newcommand{\Cpp}{C\nolinebreak\hspace{-.05em}\raisebox{.4ex}{\tiny\textbf{+}}\nolinebreak\hspace{-.10em}\raisebox{.4ex}{\tiny\textbf{+}}}

% ==============================================================================
% HIPERVÍNCULOS
% ==============================================================================
\usepackage{xcolor}
\usepackage{listings}

\usepackage{hyperref}
\hypersetup{
  colorlinks=true,
  linkcolor=blue,
  citecolor=blue,
  urlcolor=blue,
  bookmarks=true,
  bookmarksopen=true,
  bookmarksnumbered=true,
  pdfpagemode=UseOutlines,
}
\let\oldtableofcontents\tableofcontents
\renewcommand{\tableofcontents}{%
  \pdfbookmark[0]{\contentsname}{toc}%
  \oldtableofcontents%
}

% ==============================================================================
% CLEVEREF Y CONFIGURACIÓN (SIEMPRE AL FINAL)
% ==============================================================================
\usepackage{etoolbox}
\usepackage[spanish,noabbrev]{cleveref}

% =========================================================
% CÓDIGO FUENTE (listings con paleta bosque)
% =========================================================

% Paleta bosque
\definecolor{bosqueOscuro}{HTML}{0B3D2E}
\definecolor{bosqueMusgo}{HTML}{6A994E}
\definecolor{bosqueTierra}{HTML}{BC6C25}
\definecolor{bosqueCrema}{HTML}{FEFAE0}
\definecolor{bosqueTexto}{HTML}{283618}
% Versiones oscurecidas para texto (contraste >= 4.5:1 sobre el fondo crema)
\colorlet{bosqueMusgoTexto}{bosqueMusgo!80!black}   % #557A3E
\colorlet{bosqueHoja}{bosqueMusgo!70!black}         % #4A6B37
\colorlet{bosqueTierraTexto}{bosqueTierra!85!black} % #A05C1F

\lstdefinestyle{cpp}{
  language=C++,
  basicstyle=\ttfamily\small\color{bosqueTexto},
  keywordstyle=\color{bosqueTierraTexto}\bfseries,
  commentstyle=\color{bosqueMusgoTexto},
  stringstyle=\color{bosqueHoja},
  directivestyle=\color{bosqueOscuro}\bfseries,
  morekeywords={vector,string,ifstream,ofstream},
  backgroundcolor=\color{bosqueCrema},
  frame=single, rulecolor=\color{bosqueMusgo}, framesep=4pt,
  numbers=left, numberstyle=\tiny\color{bosqueMusgoTexto}, numbersep=8pt,
  xleftmargin=2.2em, xrightmargin=4.4pt, framexleftmargin=1.6em,
  breaklines=true, breakatwhitespace=true,
  postbreak=\mbox{\textcolor{bosqueMusgo}{$\hookrightarrow$}\space},
  tabsize=3, showstringspaces=false, columns=fullflexible, keepspaces=true,
  captionpos=b,
  escapeinside={(*@}{@*)},
  literate={á}{{\'a}}1 {é}{{\'e}}1 {í}{{\'i}}1 {ó}{{\'o}}1 {ú}{{\'u}}1
           {Á}{{\'A}}1 {É}{{\'E}}1 {Í}{{\'I}}1 {Ó}{{\'O}}1 {Ú}{{\'U}}1
           {ñ}{{\~n}}1 {Ñ}{{\~N}}1 {ü}{{\"u}}1 {¿}{{?`}}1 {¡}{{!`}}1
}
\lstset{style=cpp}

% Nombres en español
\renewcommand{\lstlistingname}{Código}
\renewcommand{\lstlistlistingname}{Lista de códigos}
\AtBeginDocument{%
  \crefname{listing}{código}{códigos}%
  \Crefname{listing}{Código}{Códigos}%
}

% Nombres de referencias
\crefname{Ejm}{ejemplo}{ejemplos}
\Crefname{Ejm}{Ejemplo}{Ejemplos}

\crefname{Prop}{proposición}{proposiciones}
\Crefname{Prop}{Proposición}{Proposiciones}

\crefname{Def}{definición}{definiciones}
\Crefname{Def}{Definición}{Definiciones}

\crefname{Tma}{teorema}{teoremas}
\Crefname{Tma}{Teorema}{Teoremas}

\crefname{Lema}{lema}{lemas}
\Crefname{Lema}{Lema}{Lemas}

\crefname{Cor}{corolario}{corolarios}
\Crefname{Cor}{Corolario}{Corolarios}

\crefname{Ejc}{ejercicio}{ejercicios}
\Crefname{Ejc}{Ejercicio}{Ejercicios}

\crefname{Not}{notación}{notaciones}
\Crefname{Not}{Notación}{Notaciones}

\crefname{Obs}{observación}{observaciones}
\Crefname{Obs}{Observación}{Observaciones}

\crefname{table}{tabla}{tablas}
\Crefname{table}{Tabla}{Tablas}

\crefname{section}{sección}{secciones}
\Crefname{section}{Sección}{Secciones}

\crefname{line}{línea}{líneas}%
\Crefname{line}{Línea}{Líneas}%

% Vinculación de los entornos compartidos con el contador 'equation'
\AtBeginEnvironment{Ejm}{\crefalias{equation}{Ejm}}
\AtBeginEnvironment{Prop}{\crefalias{equation}{Prop}}
\AtBeginEnvironment{Def}{\crefalias{equation}{Def}}
\AtBeginEnvironment{Tma}{\crefalias{equation}{Tma}}
\AtBeginEnvironment{Lema}{\crefalias{equation}{Lema}}
\AtBeginEnvironment{Cor}{\crefalias{equation}{Cor}}
\AtBeginEnvironment{Ejc}{\crefalias{equation}{Ejc}}
\AtBeginEnvironment{Not}{\crefalias{equation}{Not}}
\AtBeginEnvironment{Obs}{\crefalias{equation}{Obs}}

\usepackage{nccmath}
\allowdisplaybreaks %Para romper entornos align entre cambios de página
% ==============================================================================
% INICIO DEL DOCUMENTO
% ==============================================================================
\begin{document}

\pagenumbering{gobble}
\begin{center}
\thispagestyle{empty}
\includegraphics[height=4cm]{HojaTitulo/EscudoUN.pdf}\\[2.6cm]
\textbf{\huge
Topologías de Alexandroff sobre grafos simples}\\[4.2cm]
\Large\textbf{Juan Camilo Lozano Suárez}\\[4.2cm]
\small Universidad Nacional de Colombia\\
Facultad de Ciencias\\
Departamento de Matemáticas\\
Bogotá, Colombia\\
2026\\
\end{center}

\newpage{\pagestyle{empty}\cleardoublepage}

\newpage
\begin{center}
\thispagestyle{empty} \vspace*{0cm} \textbf{\huge
Topologías de Alexandroff sobre grafos simples}\\[2.5cm]
\Large\textbf{Juan Camilo Lozano Suárez}\\[2.5cm]
\small Trabajo de grado presentado como requisito parcial para optar al título de:\\
\textbf{Matemático}\\[2cm]
Directora:\\
\textbf{Ibeth Marcela Rubio Perilla}\\[1.8cm]
Líneas de Investigación:\\
Topología\\
Teoría de grafos\\[2cm]
Universidad Nacional de Colombia\\
Facultad de Ciencias\\
Departamento de Matemáticas\\
Bogotá, Colombia\\
2026\\
\end{center}
% ---------------- Página en blanco ----------------
\clearpage
\thispagestyle{empty}
\null
\clearpage

% ---------------- Dedicatoria ----------------
\thispagestyle{empty}
\begin{flushright}
   \begin{minipage}{0.6\textwidth}
      \raggedleft
      \textit{A mis padres, por su apoyo, paciencia y comprensión, que no tienen cardinal.}
   \end{minipage}
\end{flushright}
\clearpage

% ---------------- Página en blanco ----------------
\thispagestyle{empty}
\null
\clearpage% ---------------- Epígrafes ----------------
\clearpage
\thispagestyle{empty}
\begin{flushright}
   \begin{minipage}{0.6\textwidth}
      \raggedleft
      \textit{Para sobrevivir al horror hay que inventar algún tipo de belleza.}\par
      \vspace{1ex}
      ---Fernando Zalamea\footnotemark
   \end{minipage}
\end{flushright}
\footnotetext{Refiriéndose a la invención del concepto de \textit{haz} por parte de Jean Leray, en 1942, cuando este se encontraba recluido en un campo de concentración \cite{zalamea2018}.}

\vspace{3cm}

\begin{flushright}
   \begin{minipage}{0.6\textwidth}
      \raggedleft
      \textit{Nadie puede leer dos mil libros. En los cuatro siglos que vivo no habré pasado de media docena.}\par
      \vspace{1ex}
      ---Jorge Luis Borges, \textit{Utopía de un hombre que está cansado}
   \end{minipage}
\end{flushright}
\vspace*{\fill}
\clearpage
% ---------------- Página en blanco ----------------
\clearpage
\thispagestyle{empty}
\null
\clearpage

\newpage
\textbf{\LARGE Resumen}\\\\
El resumen es una presentaci\'{o}n abreviada y precisa (la NTC 1486 de 2008 recomienda revisar la norma ISO 214 de 1976). Se debe usar una extensi\'{o}n m\'{a}xima de 12 renglones. Se recomienda que este resumen sea anal\'{\i}tico, es decir, que sea completo, con informaci\'{o}n cuantitativa y cualitativa, generalmente incluyendo los siguientes aspectos: objetivos, dise\~{n}o, lugar y circunstancias, pacientes (u objetivo del estudio), intervenci\'{o}n, mediciones y principales resultados, y conclusiones. Al final del resumen se deben usar palabras claves tomadas del texto (m\'{\i}nimo 3 y m\'{a}ximo 7 palabras), las cuales permiten la recuperaci\'{o}n de la informaci\'{o}n.\\

\textbf{\small Palabras clave: (m\'{a}ximo 10 palabras, preferiblemente seleccionadas de las listas internacionales que permitan el indizado cruzado)}.\\

\textbf{\LARGE Abstract}\\\\
Es el mismo resumen pero traducido al ingl\'{e}s. Se debe usar una extensi\'{o}n m\'{a}xima de 12 renglones. Al final del Abstract se deben traducir las anteriores palabras claves tomadas del texto (m\'{\i}nimo 3 y m\'{a}ximo 7 palabras), llamadas keywords. Es posible incluir el resumen en otro idioma diferente al espa\~{n}ol o al ingl\'{e}s, si se considera como importante dentro del tema tratado en la investigaci\'{o}n, por ejemplo: un trabajo dedicado a problemas ling\"{u}\'{\i}sticos del mandar\'{\i}n seguramente estar\'{\i}a mejor con un resumen en mandar\'{\i}n.\\

\textbf{\small Keywords: palabras clave en ingl\'{e}s(m\'{a}ximo 10 palabras, preferiblemente seleccionadas de las listas internacionales que permitan el indizado cruzado)}\\


\tableofcontents

\newpage
\null
\thispagestyle{empty}
\newpage
\pagenumbering{arabic}

\addchap{Agradecimientos}
Esta secci\'{o}n es opcional, en ella el autor agradece a las personas o instituciones que colaboraron en la realizaci\'{o}n de la tesis  o trabajo de investigaci\'{o}n. Si se incluye esta secci\'{o}n, deben aparecer los nombres completos, los cargos y su aporte al documento.\\

\chapter{Preliminares}\label{chap:Preliminares}
%En este capítulo se presentan los conceptos básicos que serán usados a lo largo del trabajo organizados en dos secciones; una referente a teoría de grafos y la otra referente a topología.
\section{Teoría de Grafos}
\subsection{Introducción}
Intuitivamente, podemos pensar en un \textit{grafo simple} como una colección de círculos o puntos llamados \textit{vértices}, conectados con segmentos que llamamos \textit{aristas}.
\begin{figure}[H]
\centering
\resizebox{0.3\textwidth}{!}{
\begin{tikzpicture}[
  scale=1.0,
  vnode/.style={circle, draw=black, fill=blue!20, thick, inner sep=0pt, minimum size=6.5mm},
  edge/.style={thick, draw=black}
]
  % Vértices exteriores (cuadrado C4)
  \node[vnode] (v1) at (-2, 0) {$v_1$};
  \node[vnode] (v2) at (0, 2) {$v_2$};
  \node[vnode] (v3) at (2, 0) {$v_3$};
  \node[vnode] (v4) at (0, -2) {$v_4$};

  % Vértices polo interior/exterior (conectados a todo el C4)
  \node[vnode] (v5) at (-0.7, 0.5) {$v_5$};
  \node[vnode] (v6) at (0.7, -0.5) {$v_6$};

  % Aristas del cuadrado exterior
  \draw[edge] (v1) -- (v2) -- (v3) -- (v4) -- (v1);

  % Conexiones del vértice v5 (polo superior/interior)
  \draw[edge] (v5) -- (v1);
  \draw[edge] (v5) -- (v2);
  \draw[edge] (v5) -- (v3);
  \draw[edge] (v5) -- (v4);

  % Conexiones del vértice v6 (polo inferior/interior)
  \draw[edge] (v6) -- (v1);
  \draw[edge] (v6) -- (v2);
  \draw[edge] (v6) -- (v3);
  \draw[edge] (v6) -- (v4);

\end{tikzpicture}}%
\caption{Un grafo simple}
\end{figure}

En los grafos simples cada arista conecta exactamente dos vértices, de modo que no puede haber aristas que conecten un vértice consigo mismo, ni puede haber más de una arista entre dos vértices. Siguiendo a \cite{Foulds1991}, \cite{Gross2006} y \cite{Pirzada2007}, vemos que, a nivel práctico, los grafos simples son la herramienta estándar para modelar ciertas redes de contactos o relaciones sociales entre personas, organizar la asignación óptima de personal a diferentes tareas y planificar la infraestructura de redes físicas de computadores o servidores. También se utilizan para resolver problemas de asignación de frecuencias en antenas de telefonía móvil evitando interferencias electrónicas, e incluso en biología computacional y ciberseguridad para detectar conflictos genéticos o aislar virus informáticos en servidores clave mediante el análisis de coberturas mínimas. Según las fuentes, la gran ventaja de los grafos simples radica en la limpieza y eficiencia de su estructura. Al no admitir conexiones repetidas ni bucles sobre un mismo punto, pueden almacenarse en computadoras mediante matrices o listas de conexiones muy compactas, lo que mejora significativamente el uso de memoria y a su vez facilita el diseño de algoritmos de búsqueda y optimización extremadamente rápidos.\\\\
Una ligera variante en la noción de grafo simple la presentan los llamados \textit{grafos orientados}. Podemos pensar en un grafo orientado como en un grafo simple en el cual a cada arista se le asigna una ``dirección'', de modo que ahora tenemos registro del vértice del cual ``sale'' una arista y al cual ``llega''. Las aristas son ahora representadas por flechas apuntando desde el vértice de salida al vértice de llegada. Si el grafo orientado $D$ resulta de asignar dirección a las aristas del grafo simple $G$, decimos que $D$ es una \textit{orientación} de G.
\begin{figure}[H]
   \centering
   \resizebox{0.4\textwidth}{!}{\input{Preliminares/Grafos/Diagramas/Diag38.tex}} \hspace{2em}
   \resizebox{0.4\textwidth}{!}{\input{Preliminares/Grafos/Diagramas/Diag39.tex}} 
   \caption{$D$ es una orientación de $G$}
\end{figure}

En \cite{Foulds1991} y \cite{Roberts1978} se nos muestra que esta noción de dirección única permite organizar de forma eficiente el flujo vehicular en ciudades mediante calles de un solo sentido sin generar zonas atrapadas o inaccesibles, o representar las jerarquías de victoria y derrota en competencias de tipo ``todos contra todos''. También se nos afirma que el beneficio de un grafo dirigido frente al grafo simple es que añade una jerarquía a las conexiones; sin embargo, al prohibir arcos de ida y vuelta simultáneos entre los mismos dos puntos, conserva una estructura ligera que aún simplifica los cálculos.\\\\
Permitiendo más de una arista entre dos vértices en un grafo orientado, así como aristas de un vértice en sí mismo, obtenemos la noción de \textit{grafo dirigido} o \textit{digrafo}. Las aristas que conectan un mismo par de vértices (sin importar en qué dirección lo hagan) son llamadas \textit{aristas múltiples}. Si un conjunto de aristas múltiples tiene la misma dirección, estas son llamadas \textit{aristas paralelas}. Las aristas que unen un vétice consigo mismo son llamados \textit{bucles}. En la representación de un grafo dirigido solemos curvar ligeramente las aristas múltiples para evitar que estas se superpongan. Ahora, también etiquetamos las aristas del grafo dirigido para poder hacer distinción entre aristas paralelas.
\begin{figure}[H]
  \centering
    \resizebox{0.5\textwidth}{!}{% Recuerda incluir en el preámbulo:
% \usetikzlibrary{arrows.meta, babel}

\begin{tikzpicture}[
  scale=0.85, 
  auto=left,
  >={Stealth[length=2.5mm]},
  every node/.style={circle, fill=blue!20, draw=black, line width=1pt, inner sep=2pt, minimum size=6mm}
]

  % Posicionamiento de vértices
  \node (v1) at (3,6) {$v_1$};
  \node (v2) at (0,4) {$v_2$};
  \node (v3) at (6,4) {$v_3$};
  \node (v4) at (3,3) {$v_4$};
  \node (v5) at (0,1) {$v_5$};
  \node (v6) at (6,1) {$v_6$};
  \node (v7) at (9,1.5) {$v_7$};

  % Bucles en v1 (e1 y e2)
  % e1 bajado un poco más (yshift=-5pt)
  \path[->, line width=1pt] (v1) edge[loop above, out=120, in=60, looseness=8] node[draw=none, fill=none, above] {$e_1$} (v1);
  \path[->, line width=1pt] (v1) edge[loop above, out=150, in=30, looseness=18] node[draw=none, fill=none, above] {$e_2$} (v1);

  % Conexiones v1 <-> v2
  \draw[->, line width=1pt] (v2) to[bend left=15] node[draw=none, fill=none, above] {$e_3$} (v1);
  \draw[->, line width=1pt] (v1) to[bend left=15] node[draw=none, fill=none, below] {$e_4$} (v2);

  % Arista v1 -> v4
  \draw[->, line width=1pt] (v1) -- node[draw=none, fill=none, right] {$e_5$} (v4);

  % Arista v4 -> v5
  \draw[->, line width=1pt] (v4) -- node[draw=none, fill=none, above] {$e_6$} (v5);

  % Bucle en v5 (e7) - Ajustado para alejarlo menos (yshift/xshift reducido a -2pt)
  \path[->, line width=1pt] (v5) edge[loop left, out=210, in=150, looseness=8] node[draw=none, fill=none, xshift=-1pt, yshift=-1pt] {$e_7$} (v5);

  % Múltiples aristas v4 -> v6 (e8, e9, e10)
  \draw[->, line width=1pt] (v4) to[bend left=10] node[draw=none, fill=none, above] {$e_8$} (v6);
  \draw[->, line width=1pt] (v4) to[bend right=25] node[draw=none, fill=none, below] {$e_9$} (v6);
  \draw[->, line width=1pt] (v4) to[bend right=60] node[draw=none, fill=none, below] {$e_{10}$} (v6);

  % Arista v6 -> v4 (e11)
  \draw[->, line width=1pt] (v6) to[bend right=45] node[draw=none, fill=none, above] {$e_{11}$} (v4);

  % Arista v7 -> v3 (e12) - Mantiene orientación horizontal/vertical (sin sloped) y centrada con pos=0.5
  \draw[->, line width=1pt] (v7) -- node[draw=none, fill=none, pos=0.35, above] {$e_{12}$} (v3);

\end{tikzpicture}
} 
  \caption{Un grafo dirigido}
  \label{fig:digrafo}
\end{figure}
En el grafo dirigido de la \Cref{fig:digrafo} podemos señalar, por ejemplo, que $e_1, e_2$ y $e_7$ son bucles, $e_3$ y $e_4$ son aristas múltiples entre $v_1$ y $v_2$, y $e_8,e_9$ y $e_{10}$ son aristas paralelas. Tal como se señala en \cite{Foulds1991}, \cite{Gross2006} y \cite{Roberts1978}, esta flexibilidad hace que los digrafos sean indispensables para coordinar las etapas de proyectos complejos, determinando qué actividades deben terminarse antes de iniciar otras, como reconstruir la secuencia original de fragmentos de ADN o ARN en biología y modelar sistemas con probabilidades de transición entre estados en economía y sociología, por ejemplo. Las fuentes nos muestran que la ventaja de los grafos dirigidos sobre los grafos orientados es que permiten flujos bidireccionales y recorridos que vuelven al mismo punto, algo esencial para representar procesos dinámicos con retroalimentación o ciclos de repetición.\\\\
Si en un grafo dirigido ignoramos la dirección de sus aristas, obtenemos un \textit{multigrafo}.
\begin{figure}[H]
  \centering
  \resizebox{0.45\textwidth}{!}{% Recuerda incluir en el preámbulo:
% \usetikzlibrary{arrows.meta, babel}

\begin{tikzpicture}[
  scale=0.85, 
  auto=left,
  every node/.style={circle, fill=blue!20, draw=black, line width=1pt, inner sep=2pt, minimum size=6mm}
]

  % Posicionamiento horizontal de vértices
  \node (v1) at (0,3) {$v_1$};
  \node (v3) at (0,-3) {$v_3$};
  \node (v2) at (0,0) {$v_2$};
  \node (v4) at (2.5,0) {$v_4$};
  \node (v5) at (5,0) {$v_5$};

  % Bucles en v1, v2 y v3
  \path[line width=1pt] (v1) edge[loop left, out=135, in=225, looseness=7] node[draw=none, fill=none, left, inner sep=2pt] {$e_{12}$} (v1);
  \path[line width=1pt] (v2) edge[loop left, out=135, in=225, looseness=7] node[draw=none, fill=none, left, inner sep=2pt] {$e_{13}$} (v2);
  \path[line width=1pt] (v3) edge[loop left, out=135, in=225, looseness=7] node[draw=none, fill=none, left, inner sep=2pt] {$e_{14}$} (v3);

  % Aristas verticales iniciales
  \draw[line width=1pt] (v1) -- node[draw=none, fill=none, right, inner sep=1pt] {$e_1$} (v2);
  \draw[line width=1pt] (v2) -- node[draw=none, fill=none, right, inner sep=1pt] {$e_2$} (v3);

  % Conexiones curvas v1 -> v5 y v3 -> v5
  \draw[line width=1pt] (v1) to[bend left=20] node[draw=none, fill=none, above, inner sep=1pt, pos=0.4] {$e_3$} (v5);
  \draw[line width=1pt] (v3) to[bend right=20] node[draw=none, fill=none, below, inner sep=1pt, pos=0.4] {$e_4$} (v5);

  % Múltiples aristas entre v2 y v4
  \draw[line width=1pt] (v2) to[bend left=40] node[draw=none, fill=none, above, yshift=-2pt] {$e_5$} (v4);
  \draw[line width=1pt] (v2) -- node[draw=none, fill=none, above, yshift=-3pt] {$e_6$} (v4);
  \draw[line width=1pt] (v2) to[bend right=40] node[draw=none, fill=none, below, yshift=2pt] {$e_7$} (v4);

  % Múltiples aristas entre v4 y v5
  \draw[line width=1pt] (v4) to[bend left=40] node[draw=none, fill=none, above, yshift=-2pt] {$e_8$} (v5);
  \draw[line width=1pt] (v4) -- node[draw=none, fill=none, above, yshift=-3pt] {$e_9$} (v5);
  \draw[line width=1pt] (v4) to[bend right=40] node[draw=none, fill=none, below, yshift=5.5pt] {$e_{10}$} (v5);

  % Bucle en v5
  \path[line width=1pt] (v5) edge[loop right, out=45, in=-45, looseness=7] node[draw=none, fill=none, right, pos=0.5, inner sep=2pt] {$e_{11}$} (v5);

\end{tikzpicture}
} 
  \caption{Un multigrafo}
  \label{fig:multigrafo}
\end{figure}
En un multigrafo aún se permiten aristas múltiples entre vértices, así como bucles. En el multigrafo de la \Cref{fig:multigrafo} podemos ver, por ejemplo, que $e_8,e_9$ y $e_{10}$ son aristas múltiples entre $v_4$ y $v_5$, que $e_{12}$ es un bucle sobre $v_1$ y que $e_{14}$ es un bucle sobre $v_3$. De acuerdo con \cite{Andrei2008}, \cite{Butler2023}, \cite{Gross2006} y \cite{Shafie2012}, los multigrafos se aplican con éxito en el análisis de redes sociales para estudiar simultáneamente distintos tipos de lazos entre las mismas personas (como amistad, trabajo o préstamos de dinero), en la representación de redes de transporte con varias alternativas paralelas (como tren y autobús entre dos estaciones) y en química, para describir enlaces moleculares múltiples, entre otros. Se menciona que la principal ventaja de los multigrafos reside en que pueden registrar varios canales de interacción al mismo tiempo, sin necesidad de simplificarlos o fusionarlos en un solo enlace, preservando la riqueza del fenómeno real.\\\\
Si ahora permitimos que las aristas sean subconjuntos de vértices, no vacíos, y de cualquier cardinalidad, obtenemos la noción de \textit{hipergrafo}. Las aristas de un hipergrafo se representan como diagramas de Venn, encerrando únicamente los vértices que contiene.
\begin{figure}[H]
  \centering
    \resizebox{0.6\textwidth}{!}{% Recuerda incluir en el preámbulo:
% \usetikzlibrary{babel}

\begin{tikzpicture}[
  scale=0.9,
  every node/.style={circle, fill=blue!20, draw=black, line width=1pt, inner sep=2pt, minimum size=7mm}
]

  % ==========================================
  % 1. HIPERARISTAS (dibujadas al fondo)
  % ==========================================
  
  % e1 = {v1, v3, v4, v5} (Rojo oscuro - trazado orgánico y suave)
  \draw[red!40!black, line width=1.5pt, fill=red!10, fill opacity=0.3]
    plot [smooth cycle, tension=0.65] coordinates {
      (-1.3, 3.0)
      (-0.2, 4.4)
      (3.0, 4.6)
      (6.5, 4.3)
      (9.3, 5.8)
      (10.3, 3.0)
      (10.0, 0.7)
      (8.8, -0.1)
      (5.8, 1.8)
      (4.2, 3.7)
      (1.5, 3.7)
      (-0.4, 1.7)
    };
  \node[draw=none, fill=none, red!40!black, font=\large] at (1.5, 5.1) {$e_1$};

  % e2 = {v1, v6, v7} (Azul)
  \draw[blue!80!black, line width=1.5pt, fill=blue!10, fill opacity=0.3, rounded corners=20pt]
    (-1.2, 3) -- (-0.5, 4.4) -- (1.2, 4) -- (3.5, 1.2) -- (5.5, 0.2) -- 
    (5.8, -2.2) -- (0.8, -2.2) -- (-0.8, 0.5) -- cycle;
  \node[draw=none, fill=none, blue!80!black, font=\large] at (3, -2.6) {$e_2$};

  % e3 = {v2, v3, v4, v6} (Verde)
  \draw[green!60!black, line width=1.5pt, fill=green!10, fill opacity=0.3, rounded corners=20pt]
    (0.2, -1.2) -- (0.5, 0.5) -- (2.2, 4.4) -- (6, 4.4) -- (9.8, 6.4) -- 
    (10.6, 4.2) -- (6.5, 1.5) -- (3.5, 1.5) -- (2.5, -2) -- (0.8, -2) -- cycle;
  \node[draw=none, fill=none, green!60!black, font=\large] at (7.5, 2.8) {$e_3$};

  % e5 = {v2, v7} (Amarillo oscuro/Dorado)
  \draw[yellow!40!black, line width=1.5pt, fill=yellow!30, fill opacity=0.4, rounded corners=18pt]
    (1.8, 3.6) -- (4, 4.4) -- (6, -1) -- (3.8, -2.2) -- cycle;
  \node[draw=none, fill=none, yellow!40!black, font=\large] at (5.8, 0.8) {$e_5$};

  % e4 = {v5} (Naranja)
  \draw[orange!70!black, line width=1.5pt, fill=orange!10, fill opacity=0.3]
    (9,1) circle (8mm);
  \node[draw=none, fill=none, orange!70!black, font=\large] at (9, 2.2) {$e_4$};


  % ==========================================
  % 2. VÉRTICES (dibujados sobre las áreas)
  % ==========================================
  \node (v1) at (0,3) {$v_1$};
  \node (v2) at (3,3) {$v_2$};
  \node (v3) at (6,3) {$v_3$};
  \node (v4) at (9,5) {$v_4$};
  \node (v5) at (9,1) {$v_5$};
  \node (v6) at (1.5,-1) {$v_6$};
  \node (v7) at (4.5,-1) {$v_7$};
  \node (v8) at (12,3) {$v_8$};

\end{tikzpicture}
} 
  \caption{Un hipergrafo}
\end{figure}
Según lo señalado en \cite{Bretto2013} y \cite{Molnar2014}, los hipergrafos son muy útiles en informática para diseñar y organizar grandes bases de datos relacionales, coordinar horarios concurrentes sin solapamientos y analizar colaboraciones de grupo en redes de personas (como comités o artículos científicos escritos por múltiples autores). Nos muestran que el máximo beneficio de este tipo de grafos radica en romper la barrera del enlace por parejas, agrupando a un número arbitrario de participantes en una misma relación e independizando al modelo de simplificaciones bidireccionales o binarias artificiales.

\subsection{Conceptos básicos}
En el presente estudio trabajaremos únicamente con grafos simples, así que de ahora en adelante se entenderá \textit{grafo} y \textit{grafo simple} de manera indistinta. Iniciamos formalizando la noción de grafo que hemos introducido. Las definiciones presentadas en lo que sigue de esta sección son estándar en la teoría; las tomamos siguiendo a \cite{diestel2017}, \cite{gomez2025} y \cite{jafarian2013}.
\begin{Not}
Dado un conjunto $S$ y un número natural $k$, denotamos por $[S]^k$ a la colección de todos los subconjuntos de $S$ con exactamente $k$ elementos.
\end{Not}
Así por ejemplo, tomando $S=\left\lbrace 1,2,3,4,5 \right\rbrace$, tenemos que 
$$
\begin{aligned}
   [S]^3 &= \left\lbrace \left\lbrace 1,2,3\right\rbrace, \left\lbrace 1,2,4\right\rbrace, \left\lbrace 1,2,5\right\rbrace, \left\lbrace 1,3,4\right\rbrace, \left\lbrace 1,3,5\right\rbrace, \left\lbrace 1,4,5\right\rbrace, \left\lbrace 2,3,4\right\rbrace, \left\lbrace 2,3,5\right\rbrace \right.,\\
         &\phantom{ooo} \left. \left\lbrace 2,4,5\right\rbrace, \left\lbrace 3,4,5\right\rbrace\right\rbrace. 
\end{aligned}
$$
Los elementos de $[S]^k$ son llamados $k-subconjuntos$ de $S$. Con esto, damos la definición de grafo.
\begin{Def}
   Un \textbf{grafo} (simple) es una pareja $G=\left( V,E\right)$ de conjuntos tales que $E\subseteq [V]^2$. Los elementos de $V$ son llamados \textbf{vértices} y los elementos de $E$ son llamados \textbf{aristas}.
\end{Def}
De este modo, las aristas de un grafo son $2-subconjuntos$ de vértices. Si es necesario, dado un grafo $G$, también se denota por $V(G)$ a su conjunto de vértices, y por $E(G)$ a su conjunto de aristas. Un grafo con conjunto de vértices $V$ también es llamado un \textit{grafo sobre V}. Para evitar ambiguedades de notación, asumimos también que $V(G)\cap E(G)=\emptyset$ para todo grafo $G$. Una arista $\left\lbrace u,v\right\rbrace$ de un grafo también puede ser representada mediante $uv$. Una manera estándar y conveniente de representar visualmente un grafo es dibujar un círculo o punto por cada vértice, y unir dos de estos círculos por una línea si, y solo si, los correspondientes vértices forman una arista del grafo. La distribución particular de los círculos en el espacio, así como la forma en que las líneas son dibujadas es irrelevante; la única información que importa es qué par de vértices están conectados mediante una arista y cuáles no. Solemos referirnos por \textit{grafo} también a la respectiva representación visual.
\begin{Ejm} \label{Ejm:cuatro_grafos}
   Los siguientes son ejemplos de grafos sobre $V=\left\lbrace v_1,v_2,v_3,v_4,v_5, v_6 \right\rbrace$:
  \phantom{o}\\
{\centering
  % ================= FILA 1 =================
  \makebox[0.48\textwidth][c]{%
    \input{Preliminares/Grafos/Diagramas/Diag1.tex}%
  }%
  \hfill
  \makebox[0.48\textwidth][c]{%
    \input{Preliminares/Grafos/Diagramas/Diag2.tex}%
  }\par

  \vspace{1.5em}% Separación vertical entre filas de diagramas

  % ================= FILA 2 =================
  \makebox[0.48\textwidth][c]{%
    \input{Preliminares/Grafos/Diagramas/Diag3.tex}%
  }%
  \hfill
  \makebox[0.48\textwidth][c]{%
    \input{Preliminares/Grafos/Diagramas/Diag4.tex}%
  }\par

  \vspace{1.2em}% Espacio entre los diagramas y el título/caption
  \captionof{figure}{Ejemplos de grafos sobre $V$}
  \label{fig:ejemplos_grafos_conjuntos}
\par}
\phantom{o}\\

   Un ejemplo de grafo sobre $V'=\left\lbrace v_1, v_2, v_3, \dots\right\rbrace$ es $G_5=(V',\left\lbrace v_1v_2, v_2v_3, v_3v_4, v_4v_5, \dots\right\rbrace)$.
   \begin{figure}[H]
     \centering
       \resizebox{0.45\textwidth}{!}{\begin{tikzpicture}
  [scale=.8,auto=left,every node/.style={circle,fill=blue!20,draw=black,thick}]
  \node (n1) at (0,0) {$v_1$};
  \node (n2) at (2,0)  {$v_2$};
  \node (n3) at (4,0)  {$v_3$};
  \node (n4) at (6,0) {$v_4$};
  \node (n5) at (8,0)  {$v_5$};

  \foreach \from/\to in {n1/n2,n2/n3, n3/n4, n4/n5}
    \draw (\from) -- (\to);

   %Flecha de (n5) a un nodo invisible en (10,0)   
   \draw (n5) -- (10,0);
   
   %Etiqueta de puntos suspensivos
   \node[draw=none, fill=none, scale=1.5] at (11,0) {$\cdots$};
   % Etiqueta G_1 debajo del grafo
%   \node[draw=none, fill=none, scale=1.5] at (6,-1.5) {$G_5$};
\end{tikzpicture}
} 
     \caption{El grafo $G_5$ sobre $V'$}
   \end{figure}
\end{Ejm}
\begin{Not}
   El grafo $G=(\emptyset,\emptyset)$ en el cual tanto el conjunto de vértices como el de aristas son vacíos, lo llamamos \textbf{grafo vacío} y lo representamos simplemente por $\emptyset$.
\end{Not}
Notemos que en un grafo $G$, dada $uv\in E(G)$, se tiene que $uv$ y $vu$ hacen referencia a la misma arista, pues $uv=\left\lbrace u,v \right\rbrace=\left\lbrace v, u\right\rbrace=vu$. De esta forma, en un grafo, el orden de los vértices en la representación de una arista no es tenido en cuenta, lo cual está alineado con nuestra idea de no considerar orientación en ella. Además, la exigencia de que las aristas de un grafo sean $2-subconjuntos$ de vértices implica que si $uv\in E(G)$ entonces $u\neq v$, es decir, no hay aristas que unan un vértice con sí mismo, de modo que los grafos que consideramos no tienen bucles. Así mismo, entre dos vértices distintos $u,v\in V(G)$ podrá haber a lo más una arista, dependiendo de si $uv$ pertenece a $E(G)$ o no, es decir, se excluye la posibilidad de que hayan aristas múltiples entre un par de vértices.\\\\
Frecuentemente en matemáticas, uno de los primeros pasos al analizar una estructura formal es estudiar sus subestructuras, es decir, objetos del mismo tipo contenidos en ella. Para esto introducimos la noción de \textit{subgrafo}.
\begin{Def}
   Dados dos grafos $G=(V,E)$ y $G'=(V',E')$, decimos que $G'$ es un \textbf{subgrafo} de $G$ si $V'\subseteq V$ y $E'\subseteq E$. En este caso denotamos $G'\subseteq G$.
\end{Def}
Para cualquier grafo $G$ se tiene $\emptyset\subseteq G$. Notemos además que si $G_1$, $G_2$ y $G_3$ son grafos, con $G_1\subseteq G_2$ y $G_2\subseteq G_3$ entonces $G_1\subseteq G_3$, pues $V(G_1)\subseteq V(G_2)$ y $V(G_2)\subseteq V(G_3)$, implicando $V(G_1)\subseteq V(G_3)$, y también $E(G_1)\subseteq E(G_2)$ y $E(G_2)\subseteq E(G_3)$, implicando $E(G_1)\subseteq E(G_3)$.  Damos a continuación más ejemplos de contenencia de grafos.
\begin{Ejm}\label{Ejm:Contenencia}
   En cada fila de la \Cref{tab:grafos} se tiene $G'\subseteq G$.
   \par\vspace{1em}
\begin{longtable}{|c|c|c|}
  \hline
  \rule[-0.5em]{0pt}{1.8em} & \rule[-0.5em]{0pt}{1.8em} $G$ & \rule[-0.5em]{0pt}{1.8em} $G'$ \\
  \hline
  \endhead
  \hline
  \endfoot
  \caption{Ejemplos de contenencia de grafos} \addtocounter{table}{-1}\refstepcounter{table}\label{tab:grafos} \\
  \endlastfoot

  $G'_1 \subseteq G_1$ &
  \begin{minipage}{0.3\textwidth} \centering \vspace{8pt} \resizebox{0.85\linewidth}{!}{\input{Preliminares/Grafos/Diagramas/Diag9.tex}} \vspace{6pt} \end{minipage} &
  \begin{minipage}{0.3\textwidth} \centering \vspace{8pt} \resizebox{0.85\linewidth}{!}{\input{Preliminares/Grafos/Diagramas/Diag10.tex}} \vspace{6pt} \end{minipage} \\
  \hline

  $G'_2 \subseteq G_2$ &
  \begin{minipage}{0.3\textwidth} \centering \vspace{8pt} \resizebox{0.8\linewidth}{!}{\input{Preliminares/Grafos/Diagramas/Diag11.tex}} \vspace{6pt} \end{minipage} &
  \begin{minipage}{0.3\textwidth} \centering \vspace{8pt} \resizebox{0.8\linewidth}{!}{\input{Preliminares/Grafos/Diagramas/Diag12.tex}} \vspace{6pt} \end{minipage} \\
  \hline

  $G'_3 \subseteq G_3$ &
  \begin{minipage}{0.3\textwidth} \centering \vspace{8pt} \resizebox{0.8\linewidth}{!}{\input{Preliminares/Grafos/Diagramas/Diag13.tex}} \vspace{6pt} \end{minipage} &
  \begin{minipage}{0.3\textwidth} \centering \vspace{8pt} \resizebox{0.8\linewidth}{!}{\input{Preliminares/Grafos/Diagramas/Diag14.tex}} \vspace{6pt} \end{minipage} \\
  \hline
\end{longtable}
\par\vspace{-0.3em}

\end{Ejm}
En el \Cref{Ejm:Contenencia} se cumple que $G'_3$ conserva todas las aristas entre vértices de $V(G'_3)$ que aparecían originalmente en $G_3$, es decir, para cualesquiera $u,v\in V(G'_3)$ se tiene que $uv\in E(G_3)$ implica $uv\in E(G'_3)$. Notemos que en los ejemplos $G_1$ y $G_2$ esto no se cumple. Por ejemplo, para $G_1$ se tiene $v_3,v_5\in V(G'_1)$ y $v_3v_5\in E(G_1)$, pero $v_3v_5 \notin E(G'_1)$; para $G_2$ se tiene $v_1,v_3\in V(G'_2)$ y $v_1v_3\in E(G_2)$, pero $v_1v_3\notin E(G'_2)$. El comportamiento que se da con $G_3$ y $G'_3$ lo resumimos diciendo que $G'_3$ es el \textit{subgrafo de $G_3$ inducido} por $V(G'_3)$, y denotamos $G'_3=G[V(G'_3)]$. Más precisamente:
\begin{Def}
   Sean $G$ un grafo y $W\subseteq V(G)$. El \textbf{subgrafo de $G$ inducido} por $W$, denotado por $G[W]$, se define como el subgrafo de $G$ tal que $V(G[W])=W$ y $E(G[W])=\left\lbrace uv\in E(G)\mid u,v\in W\right\rbrace$.
\end{Def}
\begin{Ejm}
   \phantom{o}
   \begin{itemize}
      \item $G_2$ es el subgrafo de $G_1$ inducido por $W=\left\lbrace v_2,v_3,v_4,v_5,v_6\right\rbrace$:
         \begin{figure}[H]
  \centering
  \resizebox{0.2\textwidth}{!}{\input{Preliminares/Grafos/Diagramas/Diag15.tex}} \hspace{2cm}
    \resizebox{0.2\textwidth}{!}{\input{Preliminares/Grafos/Diagramas/Diag16.tex}} 
    \caption{$G_2=G_1[W]$}
\end{figure}
  
      \item $G_4$ es el subgrafo de $G_3$ inducido por $W'=\left\lbrace v_1,v_2,v_3,v_5\right\rbrace$:
         \vspace{-4em}
\begin{figure}[H]
  \hspace*{3.9cm}%
  \resizebox{0.23\textwidth}{!}{\input{Preliminares/Grafos/Diagramas/Diag17.tex}}
  \hspace*{1.15cm}
    \resizebox{0.23\textwidth}{!}{\input{Preliminares/Grafos/Diagramas/Diag18.tex}} 
    \caption{$G_4=G_3[W']$}
\end{figure}
  
   \end{itemize}
\end{Ejm}
\begin{Def}
   El \textbf{orden} de un grafo $G$, representado mediante $|G|$, es el cardinal del conjunto $V(G)$.
\end{Def}
Así, en el \Cref{Ejm:cuatro_grafos}, el grado de los cuatro primeros grafos es $6$, mientras que $|G_5|=\aleph _0$. Los grafos pueden ser \textit{finitos, infinitos, contables}, etc., dependiendo del respectivo tamaño del conjunto $G$.
\begin{Def}
   Dado un grafo $G$, decimos que dos vértices $u,v\in V(G)$ son \textbf{adyacentes} si $uv\in E(G)$. Igualmente, decimos que el vértice $v\in V(G)$ es \textbf{incidente} con la arista $e\in E(G)$ si $v\in e$. 
\end{Def}
En el grafo $G_1$ del \Cref{Ejm:cuatro_grafos}, los vértices $v_3$ y $v_4$ son adyacentes, los vértices $v_2$ y $v_5$ no son adyacentes, el vértice $v_6$ es incidente con la arista $v_1v_6$ y el vértice $v_4$ no es incidente con la arista $v_2v_3$. Los grafos finitos permiten una representación mediante matrices basada en las relaciones de adyacencia de sus vértices.
\begin{Def}
   Sea $G$ un grafo con $V(G)=\left\lbrace v_1,\dots ,v_n\right\rbrace$ para algún $n\in\mathbb{Z}^+$. La \textbf{matriz de adyacencia} $A=(a_{ij})_{n\times n}$ de $G$ se define mediante
   $$
   a_{ij} := \begin{cases} 
     1 & \text{si } v_iv_j \in E(G) \\ 
     0 & \text{en otro caso.} 
   \end{cases}
   $$
   Si $G=\emptyset$, definimos la matriz de adyacencia de $G$ como la matriz vacía de tamaño $0\times 0$.
\end{Def}
El uso de matrices de adyacencia para la representación de grafos finitos presenta utilidades para su manejo mediante computadoras. Como afirma \cite{marimont1960}, aunque los grafos representados visualmente son ``intuitivamente sugestivos y fáciles de comprender'', su manejo directo mediante computadoras resulta inviable, no se presta a manipulaciones cuantitativas inmediatas, y a medida que crecen en tamaño y complejidad, se va perdiendo la intuición que en principio proporcionan. Las matrices en cambio, pueden ser manipuladas fácilmente por computadoras, las cuales sacan vasto provecho del procesamiento cuantitativo que ofrecen, a pesar del poco atractivo intuitivo que pueda presentar al ojo humano frente a la representación visual del grafo. En las \Cref{sec:computacional_adyacencia,sec:computacional_incidencia} hacemos uso de matrices de adyacencia en algunas implementaciones computacionales.
\begin{Ejm}
   En cada fila de la \Cref{tab:grafos_matrices} se muestra un grafo finito $G$ y su respectiva matriz de adyacencia $M$.
   \par\vspace{1em}
\begin{longtable}{|c|c|c|}
  \hline
  \rule[-0.5em]{0pt}{1.8em} & \rule[-0.5em]{0pt}{1.8em} $G$ & \rule[-0.5em]{0pt}{1.8em} $M$ \\
  \hline
  \endhead
  \hline
  \endfoot
  \caption{Grafos y sus matrices de adyacencia} \addtocounter{table}{-1}\refstepcounter{table}\label{tab:grafos_matrices} \\
  \endlastfoot

  $G_1$ y $M_1$ &
  \begin{minipage}{0.3\textwidth} \centering \vspace{8pt} \resizebox{0.85\linewidth}{!}{\input{Preliminares/Grafos/Diagramas/Diag70.tex}} \vspace{6pt} \end{minipage} &
  \begin{minipage}{0.3\textwidth} \centering \vspace{8pt}
    $\begin{pmatrix}
      0 & 1 & 1 & 1 & 0 \\
      1 & 0 & 1 & 1 & 0 \\
      1 & 1 & 0 & 1 & 1 \\
      1 & 1 & 1 & 0 & 0 \\
      0 & 0 & 1 & 0 & 0
    \end{pmatrix}$
  \vspace{6pt} \end{minipage} \\
  \hline

  $G_2$ y $M_2$ &
  \begin{minipage}{0.3\textwidth} \centering \vspace{8pt} \resizebox{0.85\linewidth}{!}{\input{Preliminares/Grafos/Diagramas/Diag71.tex}} \vspace{6pt} \end{minipage} &
  \begin{minipage}{0.3\textwidth} \centering \vspace{8pt}
    $\begin{pmatrix}
      0 & 1 & 1 & 0 & 0 & 0 \\
      1 & 0 & 1 & 0 & 0 & 0 \\
      1 & 1 & 0 & 0 & 0 & 0 \\
      0 & 0 & 0 & 0 & 1 & 0 \\
      0 & 0 & 0 & 1 & 0 & 1 \\
      0 & 0 & 0 & 0 & 1 & 0
    \end{pmatrix}$
  \vspace{6pt} \end{minipage} \\
  \hline

  $G_3$ y $M_3$ &
  \begin{minipage}{0.3\textwidth} \centering \vspace{8pt} \resizebox{0.85\linewidth}{!}{\input{Preliminares/Grafos/Diagramas/Diag72.tex}} \vspace{6pt} \end{minipage} &
  \begin{minipage}{0.3\textwidth} \centering \vspace{8pt}
    $\begin{pmatrix}
      0 & 1 & 1 & 1 & 1 \\
      1 & 0 & 1 & 1 & 1 \\
      1 & 1 & 0 & 1 & 1 \\
      1 & 1 & 1 & 0 & 1 \\
      1 & 1 & 1 & 1 & 0
    \end{pmatrix}$
  \vspace{6pt} \end{minipage} \\
  \hline
\end{longtable}
\par\vspace{-0.3em}

\end{Ejm}
La ausencia de bucles en los grafos hace que todas las entradas en la diagonal principal de una matriz de adyacencia sean siempre $0$. Igualmente, el hecho de que en un grafo $G$ se tenga $v_iv_j\in E(G)$ si, y solo si, $v_jv_i\in E(G)$, hace que cualquier matriz de adyacencia sea simétrica.
\begin{Def}
   Sea $G$ un grafo. Dado un vértice $v\in V(G)$, denotamos por $\bm{A_v}$ al conjunto de todos los vértices adyacentes a $v$, es decir $A_v=\left\lbrace u\in V(G) \mid uv\in E(G)\right\rbrace$, y por $\bm{E_v}$ al conjunto de todas las aristas con las cuales es incidente $v$, es decir $E_v=\left\lbrace e\in E(G)\mid v\in e\right\rbrace$.
\end{Def}
De nuestra definicón de grafo se sigue inmediatamente que para cualquier par de vértices $u,v$ de un grafo, $u\in A_v$ si, y solo si, $v\in A_u$. En el grafo $G_2$ del \Cref{Ejm:cuatro_grafos}, tenemos que $A_{v_1}=\lbrace v_2,v_3,v_4\rbrace$, $A_{v_6}=\lbrace v_5\rbrace$, $E_{v_1}=\lbrace v_1v_2,v_1v_3,v_1v_4 \rbrace$ y $E_{v_6}=\lbrace v_5v_6\rbrace$. En el grafo $G_5$ del mismo ejemplo se tiene, para todo $i\in \left\lbrace 2,3,4,\dots\right\rbrace$, $A_{v_i}=\left\lbrace v_{i-1}, v_{i+1}\right\rbrace$ y $E_{v_i}=\left\lbrace v_{i-1}v_i, v_iv_{i+1}\right\rbrace$.\\
La siguiente proposición nos muestra el hecho natural de que la cantidad de vértices adyacentes a $v$ y la cantidad de aristas con las cuales es incidente $v$ es siempre la misma, para cualquier vértice $v$ en un grafo arbitrario.
\begin{Prop} \label{Prop:AvEv}
   En cualquier grafo $G$, dado $v\in V(G)$, se tiene $|A_v|=|E_v|$.
\end{Prop}
\begin{proof}
Definamos $f:A_v\to E_v: u\mapsto uv$. Nótese que si $u\in A_v$, entonces $uv\in E(G)$, y $v\in uv$, luego $uv\in E_v$, de modo que $f$ está bien definida. Se tiene que $f$ es inyectiva, pues si $u_1,u_2 \in A_v$ son tales que $f(u_1)=f(u_2)$ entonces $u_1v=u_2v$, es decir $\left\lbrace u_1,v\right\rbrace=\left\lbrace u_2,v\right\rbrace$, y como $u_1\neq v$ y $u_2\neq v$, entonces $u_1=u_2$. Además $f$ es sobreyectiva, pues dada $e\in E_v$, tenemos $v\in e$ y por tanto $e=uv$ para algún $u\in V(G)$; en particular $uv\in E(G)$, luego $u\in A_v$, y $f(u)=uv=e$. Por tanto $f$ es una biyección, con lo cual $|A_v|=|E_v|$.
\end{proof}
\begin{Def}
   Sea $G$ un grafo. El \textbf{grado} de un vértice $v\in V(G)$ se define como $\operatorname{deg}(v):=|A_v|$. Un vértice de grado $0$ se llama \textbf{aislado}, y uno de grado $1$ se llama una \textbf{hoja}.
\end{Def}
De este modo, el grado de un vértice se define como la cantidad de vértices adyacentes con él, que por la \Cref{Prop:AvEv}, es igual a la cantidad de aristas con las cuales es incidente. En el grafo $G_4$ del \Cref{Ejm:cuatro_grafos} se tiene $\operatorname{deg}(v_2)=3$ y $\operatorname{deg}(v_3)=4$, mientras que $v_1$ es un vértice aislado y $v_4$ es una hoja. En el grafo $G_5$ del mismo ejemplo se tiene $\operatorname{deg}(v_1)=1$ y $\operatorname{deg}(v_i)=2$ para todo $i\in \left\lbrace 2,3,4,\dots \right\rbrace$.
\begin{Def}
   Sea $G$ un grafo. Definimos el \textbf{grado mínimo} de $G$ como $\delta (G):= \min \left\lbrace \operatorname{deg}(v)\mid v\in V\right\rbrace$ y el \textbf{grado máximo} de $G$ como $\Delta (G):= \max \left\lbrace \operatorname{deg}(v)\mid v\in V\right\rbrace$, siempre que estos valores estén bien definidos.
\end{Def}
En el \Cref{Ejm:cuatro_grafos}, $\delta (G_1)=2=\Delta (G_1)$, $\delta (G_2)=1$, $\Delta (G_2)=3$, $\delta (G_3)=1$, $\Delta (G_3)=5$, $\delta (G_4)=0$, $\Delta (G_4)=4$, $\delta (G_5)=1$, $\Delta (G_5)=2$.
\begin{Def}
   Un grafo en el cual todos los vértices tienen el mismo grado $k$ se dice $\boldsymbol{k}$\textbf{-regular}, o simplemente \textbf{regular}.
\end{Def}
El grafo $G_1$ del \Cref{Ejm:cuatro_grafos} es $2$-regular. Naturalmente, en cualquier grafo $k$-regular $G$ se tiene $\delta(G)=\Delta(G)=k$.
\begin{Def}
   Un grafo se dice \textbf{localmente finito} si todos sus vértices tienen grado finito. 
\end{Def}
Todos los grafos finitos son localmente finitos. El grafo $G_5$ del \Cref{Ejm:cuatro_grafos} es un ejemplo de grafo infinito localmente finito. Se tienen otras situaciones, como se muestra a continuación.
\begin{Ejm}
   \phantom{o}
   \begin{figure}[H]
     \centering
       \resizebox{0.4\textwidth}{!}{\begin{tikzpicture}
  [scale=.8,auto=left,every node/.style={circle,fill=blue!20,draw=black,thick}]
  \node (n1) at (0,0) {$v_{1,1}$};
  \node (n2) at (2,0)  {$v_{2,1}$};
  \node (n3) at (2,2)  {$v_{2,2}$};
  \node (n4) at (4,0) {$v_{3,1}$};
  \node (n5) at (4,2)  {$v_{3,2}$};
  \node (n6) at (4,4)  {$v_{3,3}$};
  \node (n7) at (6,0)  {$v_{4,1}$};
  \node (n8) at (6,2)  {$v_{4,2}$};
  \node (n9) at (6,4)  {$v_{4,3}$};
  \node (n10) at (6,6)  {$v_{4,4}$};

  \foreach \from/\to in {n1/n2,n1/n3, n3/n4, n3/n5, n3/n6, n6/n7, n6/n8, n6/n9, n6/n10}
    \draw (\from) -- (\to);

   \draw (n10) -- (9,0);
   \draw (n10) -- (9,2);
   \draw (n10) -- (9,4);
   \draw (n10) -- (9,6);
   \draw (n10) -- (9,9);
   
   %Etiqueta de puntos suspensivos
   \node[draw=none, fill=none, scale=1.5] at (10,4) {$\cdots$};
\end{tikzpicture}
} 
     \caption{Un grafo infinito localmente finito}
   \end{figure}
\end{Ejm}
\begin{Ejm}
   El siguiente grafo infinito no es localmente finito, pues $\operatorname{deg}(v_1)$ no es finito.
   \begin{figure}[H]
     \centering
       \resizebox{0.45\textwidth}{!}{\begin{tikzpicture}
  [scale=.8,auto=left,every node/.style={circle,fill=blue!20,draw=black,thick}]
  \node (n1) at (0,2)  {$v_1$};
  \node (n2) at (0,0)  {$v_2$};
  \node (n3) at (2,0)  {$v_3$};
  \node (n4) at (4,0)  {$v_4$};
  \node (n5) at (6,0)  {$v_5$};


   \foreach \from/\to in {n1/n2,n1/n3, n1/n4}
      \draw (\from) -- (\to);
   
   \draw (n1) -- (9,0);

   \node[draw=none, fill=none, scale=1.5] at (10,0) {$\cdots$};
   
\end{tikzpicture}

} 
     \caption{Un grafo infinito que no es localmente finito}
   \end{figure}
\end{Ejm}
Un concepto importante en el presente trabajo es el de \textit{grafo conexo}. Para ello introducimos primero la definición de \textit{camino}.
\begin{Def}
   Un \textbf{camino} es un grafo no vacío $P=(V,E)$ de la forma
   $$
   V=\left\lbrace x_0, x_1, \dots, x_k\right\rbrace, \quad E=\left\lbrace x_0x_1, x_1x_2, \dots , x_{k-1}x_k\right\rbrace,
   $$
   para $k\in \mathbb{N}$, y donde todos los $x_i$, con $i\in \left\lbrace 0,1,\dots,k\right\rbrace$, son distintos. En este caso denotamos a $P$ mediante $x_0x_1\dots x_k$ o $x_kx_{k-1}\dots x_0$. Decimos que los vértices $x_0$ y $x_k$ están \textbf{conectados} por el camino $P$, y que $x_0$ y $x_k$ son los \textbf{extremos} de $P$. El número de aristas de un camino es su \textbf{longitud}. Denotamos por $\bm{P^k}$ a cualquier camino con $k$ vértices.
\end{Def}
\begin{Ejm}
   \phantom{o}
   \begin{itemize}
      \item $P^1$ es el grafo conformado por un solo vértice; es un camino de longitud $0$.
      \item Un camino de longitud $6$ es el siguiente:
      \begin{figure}[H]
        \centering
          \resizebox{0.3\textwidth}{!}{\begin{tikzpicture}
  [scale=.8,auto=left,every node/.style={circle,fill=blue!20,draw=black,thick}]
  \node (n1) at (2,4) {$v_1$};
  \node (n2) at (0,2)  {$v_2$};
  \node (n3) at (2,0)  {$v_3$};
  \node (n4) at (4,0) {$v_4$};
  \node (n5) at (6,0)  {$v_5$};
  \node (n6) at (5,2)  {$v_6$};
  \node (n7) at (3,2)  {$v_7$};

  \foreach \from/\to in {n1/n2,n2/n3, n3/n4, n4/n5, n5/n6, n6/n7}
    \draw (\from) -- (\to);
\end{tikzpicture}
}
        \caption{$P^7=v_1v_2v_3v_4v_5v_6v_7$}
      \end{figure}
   \end{itemize}
\end{Ejm}
\begin{Def}
   Sea $G$ un grafo. Decimos que dos vértices $u,v\in V(G)$ están \textbf{conectados por el camino $P$} en $G$ si $P\subseteq G$ y $P$ tiene a $u$ y $v$ como extremos.
\end{Def}
\begin{Ejm}
   \phantom{o}
   \begin{itemize}
      \item En el grafo $G_4$ del \Cref{Ejm:cuatro_grafos}, los vértices $v_6$ y $v_4$ están conectados, por ejemplo, por el camino $v_6v_2v_3v_4$, mientras que $v_1$ y $v_6$ no están conectados por un camino.
      \begin{figure}[H]
        \centering
          \resizebox{0.2\textwidth}{!}{\begin{tikzpicture}
  [scale=.8,auto=left,every node/.style={circle,fill=blue!20,draw=black,thick}]
  \node (n1) at (2,4) {$v_1$};
  \node (n2) at (4,3)  {$v_2$};
  \node (n3) at (4,1)  {$v_3$};
  \node (n4) at (2,0) {$v_4$};
  \node (n5) at (0,1)  {$v_5$};
  \node (n6) at (0,3)  {$v_6$};

  \foreach \from/\to in {n3/n4,n2/n3, n2/n5, n2/n6, n3/n5, n3/n6, n5/n6}
  \draw[opacity=0.3] (\from) -- (\to);

  \foreach \from/\to in {n6/n2, n2/n3, n3/n4}
  \draw[ultra thick] (\from) -- (\to);

%\draw[ultra thick, red] (n2) -- (n6);
\end{tikzpicture}
}
        \caption{Un camino entre $v_6$ y $v_4$}
      \end{figure}
      \item En el grafo $G_1$ del \Cref{Ejm:Contenencia}, los vértices $v_4$ y $v_7$ están conectados por el camino $v_4v_5v_7$, pero $v_2$ y $v_6$ no están conectados por un camino.
      \begin{figure}[H]
        \centering
          \resizebox{0.3\textwidth}{!}{\begin{tikzpicture}
  [scale=.8,auto=left,every node/.style={circle,fill=blue!20,draw=black,thick}]
  \node (n1) at (0,5) {$v_1$};
  \node (n2) at (3,6) {$v_2$};
  \node (n3) at (2,3) {$v_3$};
  \node (n4) at (2,0) {$v_4$};
  \node (n5) at (5,3) {$v_5$};
  \node (n6) at (5,0) {$v_6$};
  \node (n7) at (7,2) {$v_7$};

  \foreach \from/\to in {n1/n2, n3/n5, n5/n6, n6/n4, n4/n3, n3/n6, n4/n5, n5/n7}
  \draw[opacity=0.3] (\from) -- (\to);


  \foreach \from/\to in {n4/n5, n5/n7}
  \draw[ultra thick] (\from) -- (\to);
\end{tikzpicture}
}
        \caption{Un camino entre $v_4$ y $v_7$}
      \end{figure}
   \end{itemize}
\end{Ejm}
\begin{Not}
   Siempre que en un camino aparezca una expresión de la forma $p_i\dots p_j$ con $i<j$, entendemos que entre $p_i$ y $p_j$ están, aumentando su subíndice de uno en uno, todos los $p_k$ tales que $i<k<j$. Esta notación nos será útil para trabajar con caminos satisfaciendo la exigencia, dada en la definición, de que todos los $p_i$ en un camino sean distintos.
\end{Not}
\begin{Prop}\label{Prop:Conexion_caminos}
    Sean $G$ un grafo y $x,y,z\in V(G)$ cualesquiera, tales que, tanto $x$ y $y$, como $x$ y $z$, están conectados por un camino en $G$. Entonces $y$ y $z$ están conectados por un camino en $G$.
\end{Prop}
\begin{proof}
   Si $x,y,z$ no son distintos dos a dos es inmediato que $y$ y $z$ están conectados por un camino en $G$. Supongamos $x\ne y$, $x\ne z$ y $y\ne z$. Entonces existen $p_0\dots p_n$ y $p_n\dots p_m$ caminos en $G$, con $n\ge 1$, $m\ge n+1$, $p_0=y$, $p_n=x$ y $p_m=z$. Si $p_i\ne p_j$ para cualesquiera $i\in \{0,\dots,n-1\},\, j\in\{n+1,\dots,m\}$, entonces $p_0\dots p_n\dots p_m$ es un camino en $G$ que conecta $y$ y $z$. Si $p_i=p_j$ para algunos $i\in\{0,\dots ,n-1\}, \, j\in\{n+1,\dots ,m\}$ entonces tomamos $i^*$ como el menor $i\in \{0,\dots ,n-1\}$ tal que existe $j\in \{n+1,\dots ,m\}$ con $p_i=p_j$, y tomamos $j^*$ como el único $j\in \{n+1,\dots,m\}$ tal que $p_{i^*}=p_j$. Si $i^*=0$ tomamos $P=p_{j^*}\dots p_m$. Supongamos $i^*>0$. Si $j^*=m$  tomamos $P=p_0\dots p_{i^*}$. Si $j^*=m-1$, tomamos $P=p_0\dots p_{i^*}p_m$. Si $j^*<m-1$ tomamos $P=p_0\dots p_{i^*}p_{j^*+1}\dots p_m$. En cualquier caso $P$ es un camino en $G$ que conecta $y$ y $z$.
\end{proof}
\begin{Def}
   Sean $G$ un grafo y $x\in V(G)$ cualquiera. Para cualquier $y\in V(G)$ decimos que $y$ es \textbf{alcanzable} desde $x$ en $G$ si existe un camino $P\subseteq G$ que conecta $x$ y $y$.  
\end{Def}
De esta manera, en el grafo $G_4$ del \Cref{Ejm:cuatro_grafos}, se tiene que $v_4$ es alcanzable desde $v_6$ pero $v_1$ no es alcanzable desde $v_6$.
\begin{Def}
   Un grafo $G$ es llamado \textbf{conexo} si $G\neq \emptyset$ y cualesquiera par de vértices  de $G$ están conectados por un camino en $G$. Si $G$ no es conexo, decimos que es \textbf{disconexo}. 
\end{Def}
Así, en el \Cref{Ejm:cuatro_grafos}, los grafos $G_1$, $G_3$ y $G_5$ son conexos, mientras que $G_2$ y $G_4$ son disconexos. Los grafos con un solo vértice son conexos. Dado un grafo conexo $G$, podemos embeber propiamente a $G$ en un grafo conexo de la siguiente manera.
\begin{Prop}\label{Prop:Conexo+1}
   Sean $G$ un grafo conexo, $x\in V(G)$ y $z\notin V(G)$. Entonces $G':=(V(G)\cup \{z\}, E(G)\cup \{xz\})$ es un grafo conexo y $G\subsetneq G'$. 
\end{Prop}
\begin{proof}
   Es claro que $G'$ es un grafo y $G\subsetneq G'$. Veamos que $G'$ es conexo. Como $G$ es conexo, basta ver que $z$ es alcanzable en $G'$ desde cualquier vértice de $G$. Sea $y\in V(G)$ cualquiera. Como $x,y\in V(G)$, existe un camino en $G$ (y por tanto en $G'$) que conecta $x$ y $y$. Además, $xz$ es un camino en $G'$ que conecta $x$ y $z$. Por la \Cref{Prop:Conexion_caminos}, existe un camino en $G'$ que conecta $y$ y $z$. Esto completa la prueba.
\end{proof}
\begin{Def}
   Sea $G$ un grafo. Un subgrafo conexo-maximal de $G$ es llamado una \textbf{componente} de $G$. 
\end{Def}
Lo anterior nos quiere decir que un subgrafo $H\subseteq G$ es una componente de $G$ si, y solo si, $H$ es conexo y además $H$ no está propiamente contenido en otro subgrafo conexo de $G$. Naturalmente, un grafo es conexo si, y solo si, tiene una única componente. Además, como los grafos conexos son no vacíos, se sigue que el grafo vacío no tiene componentes.
\begin{Ejm}
   En los siguientes grafos se encierra con una línea punteada cada una de sus componentes.
   \begin{figure}[H]
     \centering
       \resizebox{0.45\textwidth}{!}{\begin{tikzpicture}[scale=.6,auto=left,every node/.style={circle,fill=blue!20,draw=black,thick}]
  \node (n1) at (0,0) {};
  \node (n2) at (2,3) {};
  \node (n3) at (2,1) {};
  \node (n4) at (4,3) {};
  \node (n5) at (4,0) {};
  \node (n6) at (6,0) {};
  \node (n7) at (8,0) {};
  \node (n8) at (9,2) {};
  \node (n9) at (10,0) {};
  
  \foreach \from/\to in {n1/n2,n1/n3,n1/n5,n2/n3,n2/n5,n3/n4,n3/n5,n4/n5,n5/n6,n6/n7,n7/n8,n7/n9,n8/n9}
    \draw (\from) -- (\to);

  % --- Envolvente estilo mano alzada para la componente conexa ---
  \begin{scope}[
    dashed,
    black!85,
    thick,
    decorate,
    decoration={random steps, segment length=5mm, amplitude=1.2mm}
  ]
    \draw[rounded corners=8pt] 
      (-0.8,3.8) -- (10.8,3.8) -- (10.8,-0.8) -- (-0.8,-0.8) -- cycle;
  \end{scope}

\end{tikzpicture}
}
     \caption{Un grafo con una componente}
   \end{figure}
   \begin{figure}[H]
     \centering
       \resizebox{0.35\textwidth}{!}{\begin{tikzpicture}
[scale=.6,auto=left,every node/.style={circle,fill=blue!20,draw=black,thick}]

  % Componente izquierda (Grafo completo K_4)
  \node (n1) at (0,3) {};
  \node (n2) at (3,3) {};
  \node (n3) at (0,0) {};
  \node (n4) at (3,0) {};

  \foreach \from/\to in {n1/n2,n1/n3,n1/n4,n2/n3,n2/n4,n3/n4}
    \draw (\from) -- (\to);

  % Componente derecha (Camino P_3)
  \node (n5) at (5,0) {};
  \node (n6) at (7,0) {};
  \node (n7) at (7,3) {};

  \foreach \from/\to in {n5/n6,n6/n7}
    \draw (\from) -- (\to);

  % --- Envolventes estilo mano alzada rectangulares ---
  \begin{scope}[
    dashed,
    black!85,
    thick,
    decorate,
    decoration={random steps, segment length=5mm, amplitude=1.2mm}
  ]
    % Envolvente izquierda (K_4)
    \draw[rounded corners=8pt] 
      (-0.8,3.8) -- (3.8,3.8) -- (3.8,-0.8) -- (-0.8,-0.8) -- cycle;

    % Envolvente derecha (P_3)
    \draw[rounded corners=8pt] 
      (4.2,3.8) -- (7.8,3.8) -- (7.8,-0.8) -- (4.2,-0.8) -- cycle;
  \end{scope}

\end{tikzpicture}
}
     \caption{Un grafo con dos componentes}
   \end{figure}
   \begin{figure}[H]
     \centering
       \resizebox{0.55\textwidth}{!}{\begin{tikzpicture}
[scale=.6,auto=left,every node/.style={circle,fill=blue!20,draw=black,thick}]

  % Componente izquierda
  \node (n1) at (0,2) {};
  \node (n2) at (2,2) {};
  \node (n3) at (2,0) {};
  \node (n4) at (4,2) {};

  \foreach \from/\to in {n1/n2,n1/n3,n2/n3,n3/n4}
    \draw (\from) -- (\to);

  % Vértice aislado central
  \node (n5) at (6,2) {};

  % Componente derecha
  \node (n6) at (8,2) {};
  \node (n7) at (10,2) {};
  \node (n8) at (10,0) {};
  \node (n9) at (12,2) {};

  \foreach \from/\to in {n6/n7,n6/n8,n7/n8,n7/n9,n8/n9}
    \draw (\from) -- (\to);

  % --- Envolventes estilo mano alzada rectangulares ---
  \begin{scope}[
    dashed,
    black!85, % Color más oscuro
    thick,
    decorate,
    decoration={random steps, segment length=5mm, amplitude=1.2mm}
  ]
    % Envolvente izquierda (cuadrada/rectangular)
    \draw[rounded corners=8pt] 
      (-0.8,2.8) -- (4.8,2.8) -- (4.8,-0.8) -- (-0.8,-0.8) -- cycle;

    % Envolvente centro (alrededor de n5)
    \draw[rounded corners=8pt] 
      (5.2,2.8) -- (6.8,2.8) -- (6.8,1.2) -- (5.2,1.2) -- cycle;

    % Envolvente derecha (cuadrada/rectangular)
    \draw[rounded corners=8pt] 
      (7.2,2.8) -- (12.8,2.8) -- (12.8,-0.8) -- (7.2,-0.8) -- cycle;
  \end{scope}

\end{tikzpicture}
}
     \caption{Un grafo con tres componentes}
   \end{figure}
\end{Ejm}
La colección de componentes de un grafo induce una partición del conjunto de vértices del grafo en el siguiente sentido.
\begin{Prop}\label{Prop:Componentes_particion}
   Sea $G$ un grafo. Denotamos por $\mathcal{C}$ a la colección de todas las componentes de $G$. Entonces $\{V(C)\mid C\in\mathcal{C}\}$ es una partición de $V(G)$. Más aún, para cualesquiera $C_1,C_2\in\mathcal{C}$, si $V(C_1)\cap V(C_2)\ne \emptyset$ entonces $C_1=C_2$.
\end{Prop}
\begin{proof}
   Para todo $C\in\mathcal{C}$ se tiene, por definición de componente de un grafo, que $C\ne \emptyset$, y por tanto $V(C)\ne\emptyset$. Además $C\subseteq G$ y con esto $V(C)\subseteq V(G)$. Así, $\{V(C)\mid C\in\mathcal{C}\}$ es una colección de subconjuntos no vacíos de $V(G)$.\\\\
   Veamos que $\bigcup_{C\in\mathcal{C}}V(C)=V(G)$. La contenencia $\bigcup_{C\in\mathcal{C}}V(C)\subseteq V(G)$ se sigue de que $V(C)\subseteq V(G)$ para todo $C\in\mathcal{C}$. Ahora, sea $x\in V(G)$ cualquiera. Definimos $V_x=\{y\in V(G)\mid y\text{ es alcanzable desde }x\text{ en }G\}$ y $C_x=G[V_x]$. Tenemos que $C_x\subseteq G$ y $C_x\ne \emptyset$ pues, como $x$ es alcanzable desde $x$ en $G$ por el camino $P^1$, entonces $x\in V_x=V(C_x)$. Dados $y,z\in V(C_x)=V_x$ cualesquiera, existen caminos $P,Q\subseteq G$ tales que $P$ conecta $x$ y $y$, y $Q$ conecta $x$ y $z$. Por la \Cref{Prop:Conexion_caminos} se tiene que existe un camino en $G$ que conecta $y$ y $z$. Por lo tanto $C_x$ es conexo. Sea $H$ un subgrafo conexo de $G$ cualquiera tal que $C_x\subseteq H$. Sea $w\in V(H)$ cualquiera. Como $x\in V(C_x)\subseteq V(H)$ entonces $x\in V(H)$. Como $H$ es conexo, existe un camino en $G$ que conecta $x$ y $w$, de modo que $w$ es alcanzable desde $x$ en $G$, es decir, $w\in V_x=V(C_x)$. Por tanto $V(H)\subseteq V(C_x)$. Sea $e\in E(H)$ cualquiera. Existen $w,t\in V(H)$ tales que $e=wt$. En particular $w,t\in V(C_x)=V_x$  y $wt\in E(G)$. Por tanto $e\in\{uv\in E(G)\mid u,v \in V_x\}=E(C_x)$. Así, $E(H)\subseteq E(C_x)$, con lo cual $H\subseteq C_x$ y $H=C_x$. Con lo anterior, $C_x$ es conexo-maximal en $G$ y por tanto $C_x$ es una componente de $G$. Así, $x\in V(C_x)$ y $C_x\in \mathcal{C}$, de donde $x\in \bigcup_{C\in\mathcal{C}}V(C)$ y $V(G)\subseteq \bigcup_{c\in\mathcal{C}}V(C)$. Obtenemos $\bigcup_{C\in\mathcal{C}}V(C)=V(G)$.\\\\
   Veamos que los elementos de $\{V(C)\mid C\in\mathcal{C}\}$ son disyuntos dos a dos. Tomemos $V(C_1),V(C_2)\in \{V(C)\mid C\in\mathcal{C}\}$ cualesquiera, con $C_1,C_2\in\mathcal{C}$, y supongamos $V(C_1)\ne V(C_2)$. Razonando por contradicción, supongamos que existe $x\in V(C_1)\cap V(C_2)$. Se tiene $V(C_1)\nsubseteq V(C_2)$ o $V(C_2)\nsubseteq V(C_1)$. Sin pérdidad de generalidad supongamos $V(C_2)\nsubseteq V(C_1)$. Sea $y\in V(C_2)\setminus V(C_1)$. Como $x,y\in V(C_2)$ y $C_2$ es conexo, entonces existe un camino $p_0\dots p_n\subseteq C_2$ con $n\ge 1$, $p_0=x$ y $p_n=y$. Sea $j$ el menor índice en $\{1,\dots, n\}$ tal que $p_j \in V(C_2)\setminus V(C_1)$. Notemos que $p_{j-1}\in V(C_1)$. Como además $C_1$ es conexo, entonces, por la \Cref{Prop:Conexo+1}, $(V(C_1)\cup \{p_j\},E(C_1)\cup \{p_{j-1}p_j\})$ es un subgrafo conexo de $G$ que contiene propiamente a $C_1$, contradiciendo que $C_1$ es un subgrafo conexo-maximal de $G$. Por tanto $V(C_1)\cap V(C_2)=\emptyset$, y obtenemos que $\{V(C)\mid C\in\mathcal{C}\}$ es una partición de $V(G)$.\\\\
   Ahora, sean $C_1,C_2\in\mathcal{C}$ cualesquiera, y supongamos $V(C_1)\cap V(C_2)\ne \emptyset$. De lo desarrollado anteriormente se tiene que $V(C_1)=V(C_2)$. Sea $e\in E(C_1)$ cualquiera. Entonces $e=uv$ para algunos $u,v\in V(C_1)$, luego $u,v\in V(C_2)$ y $C_3:=\left( V(C_2),E(C_2)\cup \{e\}\right)$ es un grafo, con $C_2\subseteq C_3\subseteq G$. Como $C_2$ es conexo, se sigue que $C_3$ es conexo. Pero $C_2$ es un subconjunto conexo-maximal de $G$, así que $C_2=C_3$ y por tanto $e\in E(C_2)$, con lo cual $E(C_1)\subseteq E(C_2)$. De forma similar se prueba que $E(C_2)\subseteq E(C_1)$ y por tanto $E(C_1)=E(C_2)$, con lo que concluimos que $C_1=C_2$. 
\end{proof}
Como consecuencia de la \Cref{Prop:Componentes_particion} se obtiene lo siguiente.
\begin{Prop}\label{Prop:Componentes_no_conectan}
   Sean $G$ un grafo y $C_1, C_2$ componentes distintas de $G$. Para cualesquiera $x_1\in V(C_1), x_2\in V(C_2)$ se tiene que $x_1$ y $x_2$ no están conectados por un camino en $G$.
\end{Prop}
\begin{proof}
   Por la \Cref{Prop:Componentes_particion} se tiene $x_2\notin V(C_1)$. Razonando por contradicción, supongamos que existe un camino $p_0\dots p_n$ en $G$, con $n\ge 1$, $p_0=x_1$, y $p_n=x_2$. Sea $j$ el menor índice en $\{1,\dots ,n\}$ tal que $p_j\notin V(C_1)$. Notemos que $p_{j-1}\in V(C_1)$. Por la \Cref{Prop:Conexo+1}, $\left( V(C_1)\cup \{p_{j}\},E(C_1)\cup\{p_{j-1}p_j\}\right)$ es un subgrafo conexo de $G$ que contiene propiamente a $C_1$, contradiciendo que $C_1$ es un subgrafo conexo-maximal de $G$. Por tanto $x_1$ y $x_2$ no están conectados por un camino en $G$.
\end{proof}
\begin{Obs}\label{Obs:Componentes_no_adyacentes}
   En particular, vértices de distintas componentes de un grafo son no adyacentes. 
\end{Obs}
Dado un grafo, se puede inducir un nuevo grafo al eliminar adecuadamente un subconjunto de vértices.
\begin{Def}
   Sean $G$ un grafo y $S\subseteq V(G)$. Denotamos por $\bm{G-S}$ al subgrafo de $G$ inducido por $V(G)\setminus S$, es decir, $G-S=G[V(G)\setminus S]$. 
\end{Def}
\begin{Ejm}\label{Ejm:Supresion}
   \phantom{0}
   \begin{itemize}
      \item Si $S=\left\lbrace v_3,v_4\right\rbrace$ entonces:  
         \begin{figure}[H]
  \centering
  \resizebox{0.2\textwidth}{!}{\input{Preliminares/Grafos/Diagramas/Diag15.tex}} \hspace{2cm}
  \resizebox{0.2\textwidth}{!}{\input{Preliminares/Grafos/Diagramas/Diag24.tex}} 
    \caption{$G_1$ y $G_1-S$}
\end{figure}
  
      \item Si $S'=\left\lbrace v_2,v_3,v_5\right\rbrace$ entonces:
         \vspace{1em}
         \par\vspace*{-5em}
{\centering
  \resizebox{0.23\textwidth}{!}{\input{Preliminares/Grafos/Diagramas/Diag17.tex}}%
  \hspace{2cm}%
  \resizebox{0.23\textwidth}{!}{\input{Preliminares/Grafos/Diagramas/Diag25.tex}}\par
  \vspace{0.3em}%
  \captionof{figure}{$G_3$ y $G_3-S'$}
\par}
\vspace{1em}
  
   \end{itemize}
\end{Ejm}
Como se ve con el grafo $G_3$ del \Cref{Ejm:Supresion}, la supresión de algunos vértices en un grafo puede aumentar su número de componentes. La siguiente definición trata de casos especiales donde se da esta situación.
\begin{Def}\label{Def:cortes}
   Sea $G$ un grafo. Decimos que $v\in V(G)$ es un \textbf{vértice de corte} de $G$ si $G-\left\lbrace v\right\rbrace$ tiene más componentes que $G$. Si $G$ es conexo, decimos que un subconjunto de vértices $C\subseteq V(G)$ es un \textbf{conjunto de corte} de $G$ si $G-C$ tiene más de una componente. Si además, ningún subconjunto propio de $C$ es un conjunto de corte de $G$, decimos que $C$ es un \textbf{conjunto de corte minimal} de $G$.  
\end{Def}
\begin{Ejm}
   \phantom{o}
   \begin{itemize}
      \item En el grafo $G_1$, $v_1$ es un vértice de corte, pues $G_1-\left\lbrace v_1\right\rbrace$ tiene cuatro componentes, mientras que $G_1$ tiene dos componentes.
         \begin{figure}[H]
  \centering
  \resizebox{0.35\textwidth}{!}{\input{Preliminares/Grafos/Diagramas/Diag26.tex}} \hspace{2cm}
  \resizebox{0.35\textwidth}{!}{\input{Preliminares/Grafos/Diagramas/Diag27.tex}} 
  \caption{Un vértice de corte}
\end{figure}
\vspace{-0.5cm}
      \item En el grafo $G_2$ que se muestra en la siguiente figura, tomemos $S=\{u,v,w\}$ y $S'=\{u,v\}$. Se tiene que $S$ es un conjunto de corte, pues $G_2-S$ tiene dos componentes.\vspace{0.1cm}
      \begin{figure}[H]
        \centering
        \resizebox{0.3\textwidth}{!}{\definecolor{customred}{HTML}{E7180B}

\begin{tikzpicture}[scale=0.85, auto=left]
  \tikzset{
    blue_node/.style={circle, fill=blue!20, draw, line width=1pt, inner sep=2pt, minimum size=6mm},
    red_node/.style={circle, fill=customred!20, draw=customred, line width=1pt, inner sep=2pt, minimum size=6mm}
  }

  % Vértices según la cuadrícula
  % Vértices con etiqueta
  \node[blue_node] (v) at (3,2) {$v$};
  \node[blue_node] (u) at (2,1) {$u$};
  \node[blue_node] (w) at (4,1) {$w$};

  % Vértices sin etiqueta (círculos vacíos)
  \node[blue_node] (top) at (3,4) {};
  \node[blue_node] (bottom) at (3,0) {};
  \node[blue_node] (left_bottom) at (0,0) {};
  \node[blue_node] (right_bottom) at (6,0) {};

  % Aristas del diamante interno (u, v, w, bottom)
  \draw (v) -- (u);
  \draw (v) -- (w);
  \draw (u) -- (w);
  \draw (u) -- (bottom);
  \draw (w) -- (bottom);

  % Conexiones del vértice superior (top)
  \draw (top) -- (v);
  \draw (top) -- (left_bottom);
  \draw (top) -- (right_bottom);

  % Conexión del vértice inferior izquierdo
  \draw (left_bottom) -- (u);

  % Bounding box para asegurar la alineación de la figura
  %\useasboundingbox (-0.5,-1.6) rectangle (6.5,6.5);

  % Etiqueta del grafo
  \node[draw=none, fill=none, scale=1.5] at (3,-1.5) {$G_2$};
\end{tikzpicture}
}\hspace{2cm}
        \resizebox{0.3\textwidth}{!}{\definecolor{customred}{HTML}{E7180B}

\begin{tikzpicture}[scale=0.85, auto=left]
  \tikzset{
    blue_node/.style={circle, fill=blue!20, draw, line width=1pt, inner sep=2pt, minimum size=6mm},
    red_node/.style={circle, fill=customred!20, draw=customred, line width=1pt, inner sep=2pt, minimum size=6mm}
  }

  % Vértices según la cuadrícula
  % Vértices con etiqueta
  \node[blue_node] (v) at (3,2) [opacity=0]{$v$};
  \node[blue_node] (u) at (2,1) [opacity=0]{$u$};
  \node[blue_node] (w) at (4,1) [opacity=0]{$w$};

  % Vértices sin etiqueta (círculos vacíos)
  \node[blue_node] (top) at (3,4) {};
  \node[blue_node] (bottom) at (3,0) {};
  \node[blue_node] (left_bottom) at (0,0) {};
  \node[blue_node] (right_bottom) at (6,0) {};

  % Aristas del diamante interno (u, v, w, bottom)

  % Conexiones del vértice superior (top)
  \draw (top) -- (left_bottom);
  \draw (top) -- (right_bottom);

  % Conexión del vértice inferior izquierdo

  % Bounding box para asegurar la alineación de la figura
  %\useasboundingbox (-0.5,-1.6) rectangle (6.5,6.5);

  % Etiqueta del grafo
  \node[draw=none, fill=none, scale=1.5] at (3,-1.5) {$G_2-S$};
\end{tikzpicture}
}
        \caption{$S$ es un conjunto de corte de $G_2$}
     \end{figure}\vspace{-0.4cm}
      Sin embargo, $S$ no es un conjunto de corte minimal de $G_2$, ya que $S'\subsetneq S$ y $S'$ es un conjunto de corte de $G_2$, pues $G_2-S'$ tiene dos componentes.
      \begin{figure}[H]
        \centering
          \resizebox{0.3\textwidth}{!}{\definecolor{customred}{HTML}{E7180B}

\begin{tikzpicture}[scale=0.85, auto=left]
  \tikzset{
    blue_node/.style={circle, fill=blue!20, draw, line width=1pt, inner sep=2pt, minimum size=6mm},
    red_node/.style={circle, fill=customred!20, draw=customred, line width=1pt, inner sep=2pt, minimum size=6mm}
  }

  % Vértices según la cuadrícula
  % Vértices con etiqueta
  \node[blue_node] (v) at (3,2) [opacity=0]{$v$};
  \node[blue_node] (u) at (2,1) [opacity=0]{$u$};
  \node[blue_node] (w) at (4,1) {$w$};

  % Vértices sin etiqueta (círculos vacíos)
  \node[blue_node] (top) at (3,4) {};
  \node[blue_node] (bottom) at (3,0) {};
  \node[blue_node] (left_bottom) at (0,0) {};
  \node[blue_node] (right_bottom) at (6,0) {};

  % Aristas del diamante interno (u, v, w, bottom)

  % Conexiones del vértice superior (top)
  \draw (top) -- (left_bottom);
  \draw (top) -- (right_bottom);
  \draw (w) -- (bottom);
  % Conexión del vértice inferior izquierdo

  % Bounding box para asegurar la alineación de la figura
  %\useasboundingbox (-0.5,-1.6) rectangle (6.5,6.5);

  % Etiqueta del grafo
  %\node[draw=none, fill=none, scale=1.5] at (3,-1.5) {$G_2-S'$};
\end{tikzpicture}
}
        \caption{El grafo $G_2-S'$}
      \end{figure}
      Además, $S'$ sí es un conjunto de corte minimal de $G_2$, pues ni $\{u\}$ ni $\{v\}$ son conjuntos de corte de $G_2$, ya que $G_2-\{u\}$ y $G_2-\{v\}$ son grafos conexos.\vspace{-1.2cm}
      \begin{figure}[H]
        \centering
           \resizebox{0.3\textwidth}{!}{\definecolor{customred}{HTML}{E7180B}

\begin{tikzpicture}[scale=0.85, auto=left]
  \tikzset{
    blue_node/.style={circle, fill=blue!20, draw, line width=1pt, inner sep=2pt, minimum size=6mm},
    red_node/.style={circle, fill=customred!20, draw=customred, line width=1pt, inner sep=2pt, minimum size=6mm}
  }

  % Vértices según la cuadrícula
  % Vértices con etiqueta
  \node[blue_node] (v) at (3,2) {$v$};
  \node[blue_node] (u) at (2,1) [opacity=0]{$u$};
  \node[blue_node] (w) at (4,1) {$w$};

  % Vértices sin etiqueta (círculos vacíos)
  \node[blue_node] (top) at (3,4) {};
  \node[blue_node] (bottom) at (3,0) {};
  \node[blue_node] (left_bottom) at (0,0) {};
  \node[blue_node] (right_bottom) at (6,0) {};

  % Aristas del diamante interno (u, v, w, bottom)
  \draw (v) -- (w);
  \draw (w) -- (bottom);

  % Conexiones del vértice superior (top)
  \draw (top) -- (v);
  \draw (top) -- (left_bottom);
  \draw (top) -- (right_bottom);

  % Conexión del vértice inferior izquierdo

  % Bounding box para asegurar la alineación de la figura
  \useasboundingbox (-0.5,-1.6) rectangle (6.5,6.5);

  % Etiqueta del grafo
  \node[draw=none, fill=none, scale=1.5] at (3,-1.5) {$G_2-\left\lbrace u\right\rbrace$};
\end{tikzpicture}
}\hspace{2cm}
          \resizebox{0.3\textwidth}{!}{\definecolor{customred}{HTML}{E7180B}

\begin{tikzpicture}[scale=0.85, auto=left]
  \tikzset{
    blue_node/.style={circle, fill=blue!20, draw, line width=1pt, inner sep=2pt, minimum size=6mm},
    red_node/.style={circle, fill=customred!20, draw=customred, line width=1pt, inner sep=2pt, minimum size=6mm}
  }

  % Vértices según la cuadrícula
  % Vértices con etiqueta
  \node[blue_node] (v) at (3,2) [opacity=0]{$v$};
  \node[blue_node] (u) at (2,1) {$u$};
  \node[blue_node] (w) at (4,1) {$w$};

  % Vértices sin etiqueta (círculos vacíos)
  \node[blue_node] (top) at (3,4) {};
  \node[blue_node] (bottom) at (3,0) {};
  \node[blue_node] (left_bottom) at (0,0) {};
  \node[blue_node] (right_bottom) at (6,0) {};

  % Aristas del diamante interno (u, v, w, bottom)
  \draw (u) -- (w);
  \draw (u) -- (bottom);
  \draw (w) -- (bottom);

  % Conexiones del vértice superior (top)
  \draw (top) -- (left_bottom);
  \draw (top) -- (right_bottom);

  % Conexión del vértice inferior izquierdo
  \draw (left_bottom) -- (u);

  % Bounding box para asegurar la alineación de la figura
  \useasboundingbox (-0.5,-1.6) rectangle (6.5,6.5);

  % Etiqueta del grafo
  \node[draw=none, fill=none, scale=1.5] at (3,-1.5) {$G_2-\left\lbrace v\right\rbrace$};
\end{tikzpicture}
}
          \vspace{0.5cm}
        \caption{Los grafos $G_2-\{u\}$ y $G_2-\{v\}$ son conexos}
      \end{figure}
   \end{itemize}
\end{Ejm}
\begin{Def}
   Sean $G$ un grafo y $S\subseteq V(G)$. Decimos que $S$ es \textbf{independiente} si para cualesquiera $u,v\in S$ se tiene $uv\notin E(G)$. En cambio, decimos que $S$ es un \textbf{clique} si para cualesquiera $u,v\in S$ se tiene $uv\in E(G)$.
\end{Def}
De esta forma, un subconjunto de vértices de un grafo es independiente si ningún par de estos vértices son adyacentes, y es un clique si cualquier par de estos vértices son adyacentes. En el \Cref{Ejm:cuatro_grafos}, $\left\lbrace v_2,v_4,v_6\right\rbrace$ es un conjunto independiente de $G_1$, $\left\lbrace v_1,v_2,v_3,v_4\right\rbrace$ es un clique de $G_2$, $\left\lbrace v_2,v_3,v_4,v_5,v_6\right\rbrace$ es un conjunto independiente de $G_3$ y $\left\lbrace v_2,v_3,v_5,v_6\right\rbrace$ es un clique de $G_4$.
\begin{Def}
   Dado un grafo $G=(V,E)$, su \textbf{complemento}, denotado $\overline{G}$, es el grafo con conjunto de vértices $V$ y conjunto de aristas $[V]^2\setminus E$. 
\end{Def}
Así, $\overline{G}$ mantiene los vértices de $G$, pero dos vértice son adyacentes en $\overline{G}$ si, y solo si, no son adyacentes en $G$.
\begin{Ejm}
   La siguiente figura muestra ejemplos de grafos y sus complementos.
   \par\vspace{1.5em}
{\centering
 % Fila superior (Diag33 y Diag34)
 \makebox[0.3\textwidth]{\resizebox{0.2\textwidth}{!}{\input{Preliminares/Grafos/Diagramas/Diag33.tex}}}%
 \hspace{2cm}%
 \makebox[0.3\textwidth]{\resizebox{0.2\textwidth}{!}{\input{Preliminares/Grafos/Diagramas/Diag34.tex}}}\par
 \vspace{0.5em}%
 % Fila inferior (Diag35 y Diag36)
 \makebox[0.3\textwidth]{\resizebox{0.25\textwidth}{!}{\input{Preliminares/Grafos/Diagramas/Diag35.tex}}}%
 \hspace{2cm}%
 \makebox[0.3\textwidth]{\resizebox{0.25\textwidth}{!}{\input{Preliminares/Grafos/Diagramas/Diag36.tex}}}\par
 \captionof{figure}{Grafos y sus complementos}
\par}
\vspace{1em}

\end{Ejm}
Ahora procedemos a enunciar algunas clases especiales de grafos.
\begin{Def}
   Sea $G$ un grafo. Decimos que $G$ es \textbf{completo} si $V(G)$ es un clique. Un grafo completo con $n$ vértices se denota como $\bm{K^n}$.
\end{Def}
Por tanto, un grafo es completo si todos sus vértices son adyacentes dos a dos.
\begin{Ejm}
   El grafo vacío y el grafo con un solo vértice son grafos completos, que se denotan por $K^0$ y $K^1$, respectivamente. A continuación se muestran grafos completos con $3$, $5$ y $7$ vértices.
   \begin{figure}[H]
  \renewcommand{\diagscale}{0.2\textwidth}%
   \centering
   \resizebox{\diagscale}{!}{\input{Preliminares/Grafos/Diagramas/Diag42.tex}}\hspace{4em}
   \resizebox{\diagscale}{!}{\input{Preliminares/Grafos/Diagramas/Diag43.tex}}\hspace{4em}
   \resizebox{\diagscale}{!}{\input{Preliminares/Grafos/Diagramas/Diag44.tex}}\vspace{0.5em}
 \caption{Ejemplos de grafos completos}
\end{figure}

\end{Ejm}
Para introducir los grafos conocidos como \textit{ciclos}, primero damos la siguiente definición.
\begin{Def}
   Sean $G$ un grafo y $F\subseteq [V(G)]^2$. Definimos el nuevo grafo $\bm{G+F}:=(V,E\cup F)$.
\end{Def}
De esta manera, lo que hacemos es adicionar a $G$, como nuevas aristas, las parejas de vértices que contenga $F$.
\begin{Ejm}
   \phantom{0}
   \begin{itemize}
      \item Si $F=\left\lbrace \left\lbrace v_2,v_6\right\rbrace, \left\lbrace v_2,v_5\right\rbrace, \left\lbrace v_3,v_6\right\rbrace, \left\lbrace v_3,v_5\right\rbrace\right\rbrace$ entonces:  
         \begin{figure}[H]
  \centering
  \resizebox{0.2\textwidth}{!}{\input{Preliminares/Grafos/Diagramas/Diag2.1.tex}} \hspace{2cm}
  \resizebox{0.2\textwidth}{!}{\input{Preliminares/Grafos/Diagramas/Diag45.tex}} 
    \caption{$G_2$ y $G_2+F$}
\end{figure}
  
      \item Si $F'=\left\lbrace \left\lbrace v_2,v_3\right\rbrace, \left\lbrace v_3,v_4\right\rbrace, \left\lbrace v_4,v_5\right\rbrace, \left\lbrace v_5,v_6 \right\rbrace\right\rbrace$ entonces:
         \vspace{1em}
         \begin{figure}[H]
  \centering
  \resizebox{0.2\textwidth}{!}{\input{Preliminares/Grafos/Diagramas/Diag3.1.tex}} \hspace{2cm}
  \resizebox{0.2\textwidth}{!}{\input{Preliminares/Grafos/Diagramas/Diag46.tex}} 
    \caption{$G_3$ y $G_3+F'$}
\end{figure}
  
   \end{itemize}
\end{Ejm}
\begin{Def}
   Sea $P=x_0\dots x_{k-1}$ un camino, con $k\geq 3$. Decimos que el grafo $C:=P+\left\lbrace \left\lbrace x_{k-1}, x_0\right\rbrace\right\rbrace$ es llamado un \textbf{ciclo}. En este caso denotamos a $C$ mediante $x_0x_1\dots x_{k-1}x_0$. La \textbf{longitud} de un ciclo es su número de aristas (que es igual a su número de vértices). Un ciclo de longitud $k$ es llamado un $\bm{k-ciclo}$ y es denotado por $\bm{C^k}$. 
\end{Def}
\begin{Ejm}
   Se tiene que $C^3=K^3$; éste es el ciclo de menor longitud posible. A continuación mostramos los ciclos de longitud $4$, $6$ y $8$.
   \begin{figure}[H]
  \renewcommand{\diagscale}{0.2\textwidth}%
   \centering
   \resizebox{\diagscale}{!}{\input{Preliminares/Grafos/Diagramas/Diag47.tex}}\hspace{4em}
   \resizebox{\diagscale}{!}{\input{Preliminares/Grafos/Diagramas/Diag48.tex}}\hspace{4em}
   \resizebox{\diagscale}{!}{\input{Preliminares/Grafos/Diagramas/Diag49.tex}}\vspace{1.3em}
 \caption{Ejemplos de ciclos}
\end{figure}

\end{Ejm}
Los ciclos nos permiten definir fácilmente lo que es un \textit{árbol}.
\begin{Def}
   Un grafo conexo que no contiene ningún ciclo es llamado un \textbf{árbol}. 
\end{Def}
\begin{Ejm}
   Todos los caminos son ejemplos de árboles. Los únicos grafos completos que son árboles son $K^1$ y $K^2$; el grafo $K^0=\emptyset$ no es un árbol porque no es conexo, y cualquier grafo completo con $3$ o más vértices contiene a $C^3$. En la \Cref{fig:arboles_ejemplos} los grafos $T_1$, $T_2$ y $T_3$ son ejemplos de árboles.
   \begin{figure}[H]
  \centering
  \input{Preliminares/Grafos/Diagramas/Diag50.tex}\hfill
  \input{Preliminares/Grafos/Diagramas/Diag51.tex}\hfill
  \input{Preliminares/Grafos/Diagramas/Diag52.tex}
  \caption{Ejemplos de árboles}
  \label{fig:arboles_ejemplos}
\end{figure}

   En la \Cref{fig:no_arboles}, el grafo $S_1$ no es un árbol porque no es conexo; el grafo $S_2$ no es un árbol porque contiene el ciclo $C^5$.
   \begin{figure}[H]
  \centering
  \input{Preliminares/Grafos/Diagramas/Diag53.tex}\hspace{6em}
  \input{Preliminares/Grafos/Diagramas/Diag54.tex}
  \caption{Ejemplos de grafos que no son árboles}
  \label{fig:no_arboles}
\end{figure}

\end{Ejm}
\begin{Def}
   Sea $G$ un grafo de orden mayor o igual que $2$. Decimos que $G$ es \textbf{bipartito} si $V(G)$ admite una partición en dos clases $V_1$ y $V_2$ tales que $V_1$ y $V_2$ son conjuntos independientes. Si además cualesquiera dos vértices de distintas clases de la partición son adyacentes, decimos que $G$ es \textbf{bipartito completo}, y denotamos a $G$ por $K_{n_1,n_2}$, donde $n_1=|V_1|$ y $n_2=|V_2|$. Grafos de la forma $K_{1,n}$ con $n\in \mathbb{Z}^+$ son llamados \textbf{estrellas}.
\end{Def}
De esta forma, en un grafo bipartito con partición de vértices $\left\lbrace V_1,V_2\right\rbrace$, cada vértice de $V_1$ puede ser adyacente únicamente a vértices de $V_2$, y cada vértice de $V_2$ puede ser adyacente únicamente a vértices de $V_1$. Si el grafo es además bipartito completo, cada vértice de $V_1$ es adyacente a todos los vértices de $V_2$ (y por tanto cada vértice de $V_2$ es adyacente a todos los vértices de $V_1$).\\\\
Además, una estrella es un grafo con dos o más vértices en que un único vértice es adyacente a todos los demás, y las únicas aristas son las que permiten estas adyacencias.
\begin{Obs}
    Por convención, consideramos que el grafo vacío y el grafo $K^1$ son bipartitos. Esta precisión técnica es sumamente útil, pues evita tener que arrastrar condiciones de no trivialidad en los enunciados y demostraciones de teoremas en la teoría (por ejemplo, al caracterizar los grafos bipartitos como aquellos que no contienen ciclos de longitud impar, según la Proposición~1.6.1 de \cite{diestel2017}).
\end{Obs}
\begin{Ejm}
   La \Cref{fig:bipartitos_ejemplos} muestra el grafo bipartito $G_1$, que no es bipartito completo; el grafo bipartito completo $K_{3,5}$ y la estrella $K_{1,7}$. En cada grafo las respectivas particiones del conjunto de vértices son encerradas como en un diagrama de Venn.
   \begin{figure}[H]
  \centering
  \input{Preliminares/Grafos/Diagramas/Diag56.tex}\hfill
  \input{Preliminares/Grafos/Diagramas/Diag57.tex}\hfill
  \input{Preliminares/Grafos/Diagramas/Diag58.tex}
  \caption{Ejemplos de grafos bipartitos}
  \label{fig:bipartitos_ejemplos}
\end{figure}

\end{Ejm}
La siguiente definición aborda el problema de la coloración de grafos.
\begin{Def}
   Sea $G$ un grafo. Una \textbf{coloración} de $G$ es una función $c: V(G) \to S$, donde $S$ es un conjunto cuyos elementos son llamados \textbf{colores}, tal que $c(u) \neq c(v)$ para cualesquiera dos vértices adyacentes $u, v \in V(G)$. Si el conjunto de colores $S$ tiene cardinalidad $k$, decimos que $c$ es una \textbf{$\bm{k}$-coloración}. Si $G$ tiene una \textbf{$\bm{k}$-coloración} decimos que $G$ es \textbf{$\bm{k}$-coloreable}. El menor cardinal $k$ para el cual existe una $k$-coloración de $G$ es el \textbf{número cromático} de $G$, el cual se denota por $\chi(G)$.
\end{Def}
Notemos que, para cualquier grafo $G$, la función identidad en $V(G)$ es una coloración de $G$; por tanto, todo grafo admite al menos una coloración y el número cromático $\chi(G)$ está bien definido. En particular, la función vacía de $\emptyset$ en $\emptyset$ es una coloración del grafo vacío, por lo que $\chi(\emptyset) = 0$. Para representar visualmente la coloración $c$ de un grafo, colocaremos la etiqueta correspondiente a $c(v)$ junto al círculo que representa al vértice $v$ (y no en su interior, para no confundirla con el vértice en sí). A continuación, mostramos algunos ejemplos de coloración de grafos.
\begin{Ejm}
   En la \Cref{tab:coloraciones_grafos}, para cada grafo $G_i$, con $i=1,2,3$, se tiene que $c_i$ es una coloración del grafo.
   \begin{longtable}{|S{c}|S{c}|S{c}|}
  % 1. DEFINIMOS EL PIE DE LA TABLA (AQUÍ VA EL CAPTION PARA QUE SALGA ABAJO)
  \hline
   \caption{Ejemplos de grafos con coloraciones}
  \addtocounter{table}{-1}\refstepcounter{table}% <-- Notifica a cleveref que es una tabla
  \label{tab:coloraciones_grafos} \\
  \endlastfoot
  
  % =========================================================
  % 2. A PARTIR DE AQUÍ VA EL CUERPO DE LA TABLA
  % =========================================================
  \hline
  % FILA 1: Nombres de los grafos
  \textbf{$G_1$} & \textbf{$G_2$} & \textbf{$G_3$} \\
  \hline
  
  % FILA 2: Grafos centrados en sus casillas
  \input{Preliminares/Grafos/Diagramas/Diag59.tex} &
  \input{Preliminares/Grafos/Diagramas/Diag60.tex} &
  \input{Preliminares/Grafos/Diagramas/Diag61.tex} \\
  \hline
  
  % FILA 3: Funciones (diagramas de flechas para c_1 y c_2)
  \begin{tikzpicture}[
    baseline=(t.base),
    conj/.style={draw=gray!70, thick},
    vert/.style={inner sep=1pt},
    num/.style={inner sep=1pt},
    flecha/.style={-{Stealth[length=1.8mm]}, thick},
    tramo/.style={thick},
    % Estilo de las flechas según su imagen (color y trazo distintos para impresión en blanco y negro)
    img1/.style={draw=blue!70!black, fill=blue!70!black, solid},
    img2/.style={draw=red!75!black, fill=red!75!black, dashed},
    img3/.style={draw=green!50!black, fill=green!50!black, dash pattern=on 6pt off 2pt},
    img4/.style={draw=orange!90!black, fill=orange!90!black, dash dot}
  ]
    % Encabezado: V(G_1) sobre los vértices y el codominio sobre los números
    \node (t) at (0,0) {$V(G_1)$};
    \node (cod) at (2.16,0) {$\{1, 2, 3, 4\}$};
    \node[anchor=base east, inner sep=0pt] at (t.base west) {$c_1 \!:\phantom{o}$};
    \draw[-{To[length=1.5mm]}] (t.east) -- (cod.west);
    \draw[conj] (0,-3.1) ellipse (0.6 and 2.5);
    \draw[conj] (2.16,-3.1) ellipse (0.45 and 1.95);
    \foreach \i in {1,...,7} {
      \node[vert] (v\i) at (0,{-0.7-0.6*\i}) {$v_{\i}$};
    }
    \foreach \k in {1,...,4} {
      \node[num] (n\k) at (2.16,{-3.1+(2.5-\k)*0.9}) {$\k$};
    }
    % Flechas: las que llegan al mismo número tienen el mismo color y trazo.
    % Primero se trazan los segmentos y luego, encima, las puntas en la mitad.
    \foreach \i/\k in {1/1, 2/2, 3/3, 4/4, 5/1, 6/3, 7/1} {
      \draw[tramo, img\k, shorten <=1pt, shorten >=1pt] (v\i.east) -- (n\k.west);
    }
    \foreach \i/\k in {1/1, 2/2, 3/3, 4/4, 5/1, 6/3, 7/1} {
      \path (v\i.east) -- (n\k.west) coordinate[pos=0.5] (m);
      \draw[flecha, img\k, shorten <=1pt] (v\i.east) -- (m);
    }
  \end{tikzpicture}
  &
  \begin{tikzpicture}[
    baseline=(t.base),
    conj/.style={draw=gray!70, thick},
    vert/.style={inner sep=1pt},
    num/.style={inner sep=1pt},
    flecha/.style={-{Stealth[length=1.8mm]}, thick},
    tramo/.style={thick},
    % Estilo de las flechas según su imagen (color y trazo distintos para impresión en blanco y negro)
    img1/.style={draw=blue!70!black, fill=blue!70!black, solid},
    img2/.style={draw=red!75!black, fill=red!75!black, dashed},
    img3/.style={draw=green!50!black, fill=green!50!black, dash pattern=on 6pt off 2pt},
    img4/.style={draw=orange!90!black, fill=orange!90!black, dash dot}
  ]
    % Encabezado: V(G_2) sobre los vértices y el codominio sobre los números
    \node (t) at (0,0) {$V(G_2)$};
    \node (cod) at (2.16,0) {$\{1, 2, 3\}$};
    \node[anchor=base east, inner sep=0pt] at (t.base west) {$c_2 \!:\phantom{o}$};
    \draw[-{To[length=1.5mm]}] (t.east) -- (cod.west);
    \draw[conj] (0,-4) ellipse (0.6 and 3.45);
    \draw[conj] (2.16,-4) ellipse (0.45 and 1.6);
    \foreach \i in {1,...,10} {
      \node[vert] (v\i) at (0,{-0.7-0.6*\i}) {$v_{\i}$};
    }
    \foreach \k in {1,2,3} {
      \node[num] (n\k) at (2.16,{-4+(2-\k)*1.0}) {$\k$};
    }
    % Flechas: las que llegan al mismo número tienen el mismo color y trazo.
    % Primero se trazan los segmentos y luego, encima, las puntas en la mitad.
    \foreach \i/\k in {1/1, 2/2, 3/3, 4/1, 5/2, 6/3, 7/1, 8/2, 9/2, 10/2} {
      \draw[tramo, img\k, shorten <=1pt, shorten >=1pt] (v\i.east) -- (n\k.west);
    }
    \foreach \i/\k in {1/1, 2/2, 3/3, 4/1, 5/2, 6/3, 7/1, 8/2, 9/2, 10/2} {
      \path (v\i.east) -- (n\k.west) coordinate[pos=0.45] (m);
      \draw[flecha, img\k, shorten <=1pt] (v\i.east) -- (m);
    }
  \end{tikzpicture}
  &
  $\begin{aligned}[t]
     c_3 \!:\phantom{o} V(G_3) &\to \{1, 2\} \\
     v_i &\mapsto 
     \begin{cases} 
       1 & \begin{array}[t]{@{}l@{}} \text{si } i \text{ es} \\[-1ex] \text{impar} \end{array} \\[0.6ex]
       2 & \begin{array}[t]{@{}l@{}} \text{si } i \text{ es} \\[-1ex] \text{par} \end{array} 
     \end{cases}
  \end{aligned}$ \\
  \hline
\end{longtable}

\end{Ejm}

\subsection{Homomorfismos}
\begin{comment}
En lo que resta de esta sección nos ocuparemos de relacionar parejas de grafos mediante funciones.
\begin{Def}\label{Def:homomorfismo}
   Sean $G=(V,E)$ y $G'=(V',E')$ dos grafos. Decimos que $\varphi:V\to V'$ es un homomorfismo de $G$ en $G'$ si $\left\lbrace u,v\right\rbrace\in E$ implica $\left\lbrace \varphi(u),\varphi(v) \right\rbrace\in E'$ para cualesquiera vértices $u,v\in V$.
\end{Def}
El hecho de que $\left\lbrace u,v \right\rbrace\in E$ implique $\left\lbrace \varphi(u),\varphi(v)\right\rbrace\in E'$ lo resumimos diciendo que $\varphi$ preserva la adyacencia entre vértices, puesto que $\varphi$ envía vértices adyacentes de $G$ a vértices adyacentes de $G'$.
\begin{Ejm}\label{Ejm:homomorfismos}
   La función
      \begin{center}
   \begin{tikzpicture}[
     conj/.style={draw=gray!70, thick},
     vert/.style={inner sep=1pt},
     num/.style={inner sep=1pt},
     flecha/.style={-{Stealth[length=1.8mm]}, thick},
     tramo/.style={thick},
     % Estilo de las flechas según su imagen (color y trazo distintos para impresión en blanco y negro)
     img1/.style={draw=blue!70!black, fill=blue!70!black, solid},
     img2/.style={draw=red!75!black, fill=red!75!black, dash dot},
     img3/.style={draw=green!50!black, fill=green!50!black, dash pattern=on 6pt off 2pt}
   ]
     % Encabezado: V(G) sobre los vértices de G y V(G') sobre los de G'
     \node (t) at (0,0) {$V(G)$};
     \node (cod) at (2.6,0) {$V(G')$};
     \node[anchor=base east, inner sep=0pt] at (t.base west) {$\varphi_1 \!:\phantom{o}$};
     \draw[-{To[length=1.5mm]}] (t.east) -- (cod.west);
     \draw[conj] (0,-2.5) ellipse (0.6 and 1.9);
     \draw[conj] (2.6,-2.5) ellipse (0.6 and 1.7);
     \foreach \i in {1,...,5} {
       \node[vert] (v\i) at (0,{-0.7-0.6*\i}) {$v_{\i}$};
     }
     \foreach \k in {1,...,4} {
       \node[num] (n\k) at (2.6,{-2.5+(2.5-\k)*0.75}) {$w_{\k}$};
     }
     % Flechas: las que llegan al mismo vértice tienen el mismo color y trazo.
     % Primero se trazan los segmentos y luego, encima, las puntas.
     \foreach \i/\k in {1/1, 2/2, 3/3, 4/2, 5/1} {
       \draw[tramo, img\k, shorten <=1pt, shorten >=1pt] (v\i.east) -- (n\k.west);
     }
     \foreach \i/\k in {1/1, 2/2, 3/3, 4/2, 5/1} {
       \path (v\i.east) -- (n\k.west) coordinate[pos=0.62] (m);
       \draw[flecha, img\k, shorten <=1pt] (v\i.east) -- (m);
     }
   \end{tikzpicture}
   \end{center}
   es un homomorfismo de $G$ en $G'$, donde $G$ y $G'$ son los grafos que se muestran a continuación.
   \begin{figure}[H]
  \centering
  % Grafos G y G' junto con las flechas de la función \varphi_1
  \scalebox{0.7}{% Ajusta este valor (ej. 0.6, 0.7, 0.8) para controlar el tamaño
    \input{Preliminares/Grafos/Diagramas/Diag73.tex}
  }
  \caption{$\varphi_1$ es un homomorfismo de $G$ en $G'$}
  \label{fig:homomorfismo_grafos}
\end{figure}

   Para verificar que $\varphi_1$ es un homomorfismo notemos que $\varphi_1$ preserva la adyacencia entre vértices de $G$, revisando cada arista de $G$ y su correspondiente imagen. Por ejemplo:
   \begin{center}
   $\hfill
   \begin{array}{rcl}
      \{v_1, v_2\} \in E(G) &\text{ y }& \{\varphi_1(v_1), \varphi_1(v_2)\} = \{w_1, w_2\} \in E(G'), \\
      \{v_1, v_3\} \in E(G) &\text{ y }& \{\varphi_1(v_1), \varphi_1(v_3)\} = \{w_1, w_3\} \in E(G').
   \end{array}
   \hfill$
   \end{center}
   En cambio la función
   \begin{center}
   $\hfill
   \begin{array}{rcl}
      \varphi_2 \colon V(G) &\to& V(G') \\
      v_1 &\mapsto& w_4 \\
      v_2 &\mapsto& w_2 \\
      v_3 &\mapsto& w_4 \\
      v_4 &\mapsto& w_1 \\
      v_5 &\mapsto& w_3
   \end{array}
   \hfill$
   \end{center}
   no es un homomorfismo de $G$ en $G'$, pues $\left\lbrace v_1,v_3\right\rbrace\in E(G)$ pero $\left\lbrace \varphi(v_1),\varphi(v_3)\right\rbrace=\left\lbrace w_4\right\rbrace\notin E(G')$.\\
   Podemos ver también que la función
   \begin{center}
   $\hfill
   \begin{array}{rcl}
      \psi \colon V(H) &\to& V(H') \\
      x_1 &\mapsto& y_1 \\
      x_2 &\mapsto& y_2 \\
      x_3 &\mapsto& y_3 \\
      x_4 &\mapsto& y_4 
   \end{array}
   \hfill$
   \end{center}
   es un homomorfismo entre los grafos $H$ y $H'$, que se presentan en la \Cref{fig:homomorfismo_grafos_2}.
   \begin{figure}[H]
  \centering
  % Primer grafo (G_1)
  \begin{minipage}[t]{0.49\textwidth}
    \centering
    \scalebox{0.7}{% Ajusta este valor (ej. 0.6, 0.7, 0.8) para controlar el tamaño
      \input{Preliminares/Grafos/Diagramas/Diag64.tex}
    }
  \end{minipage}%
  \hspace{-3.5em}%
  % Segundo grafo (G'_1)
  \begin{minipage}[t]{0.49\textwidth}
    \centering
    \scalebox{0.7}{% Usa el mismo factor de escala para que guarden proporción
      \input{Preliminares/Grafos/Diagramas/Diag65.tex}
    }
  \end{minipage}
  
  \caption{$H$ y $H'$}
  \label{fig:homomorfismo_grafos_2}
\end{figure}

\end{Ejm}
Estamos interesados en encontrar funciones entre vértices que permitan determinar si los grafos en cuestión son ``indistinguibles'' para la teoría construida. Un homomorfismo no es suficiente para esto: notemos que en la \Cref{Def:homomorfismo}, un homomorfismo puede enviar vértices no adyacentes en el grafo de partida a vértices adyacentes en el grafo de llegada. Como muestra de esto, en el \Cref{Ejm:homomorfismos}, se tiene que $\varphi_1$ es un homomorfismo de $G$ en $G'$, y a pesar de que $v_3$ y $v_5$ no son adyacentes en $G$, vemos que $\varphi_1(v_3)=w_3$ y $\varphi_1(v_5)=w_1$ sí son adyacentes en $G'$. Un primer paso sería exigir que la función, además de preservar la adyacencia de los vértices del grafo de partida, no envíe vértices no adyacentes en vértices adyacentes. Sin embargo, esta propiedad la cumple la función $\psi$ del \Cref{Ejm:homomorfismos}, conectando al grafo $H$ de $4$ vértices y $4$ aristas con el grafo $H'$ de $8$ vértices y $12$ aristas. Incluir la exigencia de que la función sea biyectiva solventa el problema y se obtiene la definición de \textit{isomorfismo} entre grafos.
\begin{Def}
   Sean $G=(V,E)$ y $G'=(V',E')$ dos grafos. Decimos que una función $\varphi:V\to V'$ es un \textbf{isomorfismo} de $G$ en $G'$ si $\varphi$ es biyectiva y $\left\lbrace u,v\right\rbrace\in E$, si y sólo si, $\left\lbrace \varphi(u), \varphi(v)\right\rbrace\in E'$, para cualesquiera vértices $u,v\in V$. Si existe un isomorfismo de $G$ en $G'$ decimos que $G$ y $G'$ son \textbf{isomorfos} y escribimos $G\simeq G'$. Un isomorfismo de $G$ en sí mismo es llamado un \textbf{automorfismo} de $G$.
\end{Def}
Usualmente no hacemos distinción alguna entre grafos isomorfos. Esto es lo que nos permite hablar de ``\textit{el} camino de $6$ vértices'', ``\textit{el} grafo completo de $4$ vértices'', etc. Nótese que en cualquier grafo $G$ la función identidad de $V(G)$ es un automorfismo de $G$.
\begin{Ejm}
   En la \Cref{fig:primer_isomorfismo} los grafos $G$ y $G'$ son isomorfos. 
   \begin{figure}[H]
  \centering
  % Primer grafo (G_1)
  \begin{minipage}[t]{0.49\textwidth}
    \centering
    \scalebox{0.7}{% Ajusta este valor (ej. 0.6, 0.7, 0.8) para controlar el tamaño
      \input{Preliminares/Grafos/Diagramas/Diag66.tex}
    }
  \end{minipage}%
  \hspace{-3em}%
  % Segundo grafo (G'_1)
  \begin{minipage}[t]{0.49\textwidth}
    \centering
    \scalebox{0.7}{% Usa el mismo factor de escala para que guarden proporción
      \input{Preliminares/Grafos/Diagramas/Diag67.tex}
    }
  \end{minipage}
  
  \caption{$G$ y $G'$ son isomorfos}
  \label{fig:primer_isomorfismo}
\end{figure}

   Consideremos la función $\varphi:V(G)\to V(G'):v_i\mapsto w_i$ para $i=1,2,\dots ,7$. Se tiene que $\varphi$ es biyectiva. Además
   \begin{center}
   $\hfill
   \begin{array}{rcl}
      \{v_1, v_2\} \in E(G) &\text{ y }& \{\varphi(v_1), \varphi(v_2)\} = \{w_1, w_2\} \in E(G'), \\
      \{v_1, v_3\} \in E(G) &\text{ y }& \{\varphi(v_1), \varphi(v_3)\} = \{w_1, w_3\} \in E(G'), \\
      \{v_2, v_3\} \in E(G) &\text{ y }& \{\varphi(v_2), \varphi(v_3)\} = \{w_2, w_3\} \in E(G'), \\
      \{v_3, v_7\} \in E(G) &\text{ y }& \{\varphi(v_3), \varphi(v_7)\} = \{w_3, w_7\} \in E(G'), \\
      \{v_3, v_6\} \in E(G) &\text{ y }& \{\varphi(v_3), \varphi(v_6)\} = \{w_3, w_6\} \in E(G'), \\
      \{v_6, v_7\} \in E(G) &\text{ y }& \{\varphi(v_6), \varphi(v_7)\} = \{w_6, w_7\} \in E(G'), \\
      \{v_4, v_7\} \in E(G) &\text{ y }& \{\varphi(v_4), \varphi(v_7)\} = \{w_4, w_7\} \in E(G'), \\
      \{v_4, v_5\} \in E(G) &\text{ y }& \{\varphi(v_4), \varphi(v_5)\} = \{w_4, w_5\} \in E(G'), \\
      \{v_5, v_6\} \in E(G) &\text{ y }& \{\varphi(v_5), \varphi(v_6)\} = \{w_5, w_6\} \in E(G').
   \end{array}
   \hfill$
   \end{center}
   Notemos que en la anterior lista se verifica tanto que $\left\lbrace u,v\right\rbrace\in E(G)$ implica $\left\lbrace \varphi(u), \varphi(v) \right\rbrace \allowbreak \in E(G')$, como que $\left\lbrace \varphi(u), \varphi(v)\right\rbrace\in E(G')$ implica $\left\lbrace u,v\right\rbrace\in E(G)$, para cualesquiera $u,v\in V(G)$. Por tanto $\varphi$ es un isomorfismo de $G$ en $G'$.
\end{Ejm}
El siguiente ejemplo muestra que no es suficiente que una función sea un homomorfismo biyectivo para que sea un isomorfismo.
\begin{Ejm}
   La función $\varphi:V(G)\to V(G'):v_i\mapsto w_i$ para $i=1,2$, es un homomorfismo biyectivo entre los grafos $G$ y $G'$ de la \Cref{fig:homomorfo_no_isomorfo}. Vacíamente se tiene que $\varphi$ preserva la adyacencia entre vértices.
   \begin{figure}[H]
\centering
  \resizebox{0.4\textwidth}{!}{\input{Preliminares/Grafos/Diagramas/Diag68.tex}}%
  \caption{$G$ y $G'$ no son isomorfos}
  \label{fig:homomorfo_no_isomorfo}
\end{figure}

   Sin embargo $\varphi$ no es un isomorfismo, pues $\left\lbrace \varphi(v_1),\varphi(v_2)\right\rbrace=\left\lbrace w_1,w_2\right\rbrace\in E(G')$ pero $\left\lbrace v_1,v_2\right\rbrace\notin E(G)$.
\end{Ejm}
Tenemos en cambio que si una función $\varphi$ es un homomorfismo biyectivo y su inversa $\varphi^{-1}$ es también un homomorfismo, entonces $\varphi$ sí es un isomorfismo. De hecho, esto caracteriza a los isomorfismos entre grafos, como muestra la siguiente proposición.
\begin{Prop}\label{Prop:Caracterizacion_isomorfismo}
   Sean $G$ y $G'$ dos grafos y $\varphi:V(G)\to V(G')$ una función. Las siguientes afirmaciones son equivalentes:
   \begin{itemize}
      \item[1.] $\varphi$ es un homomorfismo biyectivo de $G$ en $G'$ y $\varphi^{-1}$ es un homomorfismo de $G'$ en $G$.
      \item[2.] $\varphi$ es un isomorfismo de $G$ en $G'$.
   \end{itemize}
\end{Prop}
\begin{proof}
   $1\Rightarrow 2$ : Supongamos que $\varphi$ es un homomorfismo biyectivo de $G$ en $G'$ y $\varphi^{-1}$ es un homomorfismo de $G'$ en $G$. Sean $u,v \in V(G)$. Si $\left\lbrace u,v\right\rbrace\in E(G)$, como $\varphi$ es un homomorfismo, $\left\lbrace \varphi(u),\varphi(v)\right\rbrace\in E(G')$. Si $\left\lbrace \varphi(u),\varphi(v)\right\rbrace\in E(G')$, como $\varphi^{-1}$ es un homomorfismo y $\varphi$ es biyectivo, entonces $\left\lbrace u,v\right\rbrace=\left\lbrace \varphi^{-1}(\varphi(u)),\varphi^{-1}(\varphi(v))\right\rbrace\in E(G)$. Por tanto $\varphi$ es un isomorfismo de $G$ en $G'$.\par
   $2\Rightarrow 1$ : Supongamos que $\varphi$ es un isomorfismo de $G$ en $G'$. Por definición de isomorfismo, $\varphi$ es biyectiva. Si $u,v\in V(G)$ son tales que $\left\lbrace u,v\right\rbrace\in E(G)$ entonces $\left\lbrace \varphi(u),\varphi(v)\right\rbrace\in E(G')$, y por tanto $\varphi$ es un homomorfismo de $G$ en $G'$. Supongamos que $x,y\in V(G')$ son tales que $\left\lbrace x,y\right\rbrace\in E(G')$. Como $\varphi$ es biyectiva, existen $w_x,w_y\in V(G)$ tales que $\varphi(w_x)=x$ y $\varphi(w_y)=y$, luego $\left\lbrace \varphi(w_x),\varphi(w_y)\right\rbrace\in E(G')$. Como $\varphi$ es un isomorfismo, se tiene $\left\lbrace \varphi^{-1}(x),\varphi^{-1}(y)\right\rbrace=\left\lbrace w_x,w_y\right\rbrace\in E(G)$. Por tanto $\varphi^{-1}$ es un homomorfismo de $G'$ en $G$.
\end{proof}
Como consecuencia inmediata de la \Cref{Prop:Caracterizacion_isomorfismo} se tiene lo siguiente.
\begin{Obs}\label{Cor:Inversa_isomorfismo}
   Si $\varphi$ es un isomorfismo de $G$ en $G'$ entonces $\varphi^{-1}$ es un isomorfismo de $G'$ en $G$.
\end{Obs}
La siguiente proposición nos muestra que un isomorfismo preserva el grado de los vértices.
\begin{Prop}\label{Prop:Isomorfismo_grado}
   Sean $G$ y $G'$ dos grafos y $\varphi$ un isomorfismo de $G$ en $G'$. Para todo $v\in V(G)$ se tiene $\operatorname{deg}(v)=\operatorname{deg}(\varphi(v))$.    
\end{Prop}
\begin{proof}
   Sea $v\in V(G)$ cualquiera. Basta probar que $|A_v|=|A_{\varphi(v)}|$. Veamos que $\left.\varphi\right|_{A_v}$ es una biyección de $A_v$ en $A_{\varphi(v)}$. La inyectividad de $\left.\varphi\right|_{A_v}$ se hereda de la inyectividad de $\varphi$. Sea $w\in A_{\varphi(v)}$ cualquiera. Como $\varphi$ es sobreyectiva, existe $u\in V(G)$ tal que $w=\varphi(u)$, de donde $\varphi(u)\in A_{\varphi(v)}$, es decir $\varphi(u)\varphi(v)\in E(G')$. Como $\varphi$ es un isomorfismo, se tiene $uv\in E(G)$, es decir, $u\in A_v$. Con esto, $\left.\varphi\right|_{A_v}(u)=\varphi(u)=w$ y por tanto $\left.\varphi\right|_{A_v}$ es una función sobreyectiva de $A_v$ en $A_{\varphi(v)}$. Así, $\left.\varphi\right|_{A_v}$ es una biyección de $A_v$ en $A_{\varphi(v)}$, y por tanto $|A_v|=|A_{\varphi(v)}|$.
\end{proof}
Un corolario de la \Cref{Prop:Isomorfismo_grado} es el siguiente.
\begin{Cor}\label{Cor:Isomorfismo_grados_iguales}
   Sean $G$ y $G'$ dos grafos isomorfos. Entonces $G$ tiene un vértice de grado $n$ si, y solo si, $G'$ tiene un vértice de grado $n$. 
\end{Cor}
\begin{proof}
   Existe un isomorfismo $\varphi$ de $G$ en $G'$.\\
   ($\Rightarrow$) Supongamos que $G$ tiene un vértice $v$ de grado $n$. Entonces $\operatorname{deg} (\varphi(v)) = \operatorname{deg} (v) = n$, con $\varphi(v)\in V(G')$.\\
   ($\Leftarrow$) Supongamos que $G'$ tiene un vértice $w$ de grado $n$. Como $\varphi ^{-1}$ es un isomorfismo de $G'$ en $G$ entonces $\operatorname{deg}\left( \varphi^{-1}(w)\right)=\operatorname{deg} (w)=n$, con $\varphi ^{-1}(w)\in V(G)$.
\end{proof}
\begin{Def}
   Decimos que un grafo $G$ es \textbf{vértice transitivo} si para cualquier par de vértices $u,v\in V(G)$ existe un automorfismo de $G$ que envía $u$ a $v$.    
\end{Def}
\begin{Ejm}
   Para todo $n\in\mathbb{N}$ con $n\ge 3$ se tiene que $C^n$ es un grafo vértice transitivo. Podemos suponer que $V(C^n)$ es de la forma $\left\lbrace v_0,v_1, \dots, v_{n-1}\right\rbrace$ con $E(C^n)=\left\lbrace v_0v_1, v_1v_2, \dots ,v_{n-2}v_{n-1}, v_{n-1}v_0\right\rbrace$. Sean $v_i,v_j\in V(C^n)$ cualesquiera. La función $\varphi :V(C^n)\to V(C^n):v_k\mapsto v_{(k+j-i)\bmod n}$ para cada $k\in\left\lbrace 0,1,\dots ,n-1\right\rbrace$, es un automorfismo que envía $v_i$ a $v_j$. Si representamos a $C^n$ como en la \Cref{fig:ciclo_n}, podemos pensar en $\varphi$ como una función que ``rota'' al grafo en sentido horario hasta llevar al vértice $v_i$ a la posición del vértice $v_j$.
   \begin{figure}[H]
\centering
\resizebox{0.35\textwidth}{!}{
   \begin{tikzpicture}[scale=0.9, auto=left, every node/.style={circle, fill=blue!20, draw=black, thick, inner sep=0pt, minimum size=0.9cm}]

     % Vértices visibles de v_0 a v_8 en sentido horario
     \node (v0) at (90:3) {$v_0$};
     \node (v1) at (54:3) {$v_1$};
     \node (v2) at (18:3) {$v_2$};
     \node (v3) at (-18:3) {$v_3$};
     \node (v4) at (-54:3) {$v_4$};
     \node (v5) at (-90:3) {$v_5$};
     \node (v6) at (-126:3) {$v_6$};
     \node (v7) at (-162:3) {$v_7$};
     \node (v8) at (198:3) {$v_8$};

     % Último vértice v_{n-1}
     \node (vn1) at (126:3) {$v_{n-1}$};

     % Vértices invisibles y puntos para la discontinuidad
     \node[draw=none, fill=none, minimum size=0pt] (v9_inv) at (180:3) {};
     \node[draw=none, fill=none, minimum size=0pt] (vn2_inv) at (144:3) {};
     \node[draw=none, fill=none, minimum size=0pt] (dots) at (162:3) {$\dots$};

     % Aristas continuas principales
     \draw (v0) -- (v1);
     \draw (v1) -- (v2);
     \draw (v2) -- (v3);
     \draw (v3) -- (v4);
     \draw (v4) -- (v5);
     \draw (v5) -- (v6);
     \draw (v6) -- (v7);
     \draw (v7) -- (v8);

     % Arista normal entre v_{n-1} y v_0
     \draw (vn1) -- (v0);

     % Aristas hacia/desde la zona discontinua
     \draw (v8) -- (v9_inv);
     \draw (vn2_inv) -- (vn1);

   \end{tikzpicture}}%
   \caption{El grafo $C^n$}
   \label{fig:ciclo_n}
\end{figure}

\end{Ejm}

\section{Topología}
En esta sección se presentan los principales conceptos y resultados referentes a topología que serán usados en los capítulos posteriores. Para los conceptos básicos de topología general se considera principalmente \cite{munkres2000}. Para la teoría referente a \textit{espacios de Alexandroff}, que no se suele incluir en cursos estándar de topología, se sigue a \cite{alexandroff1937}, \cite{arenas1999}, \cite{mesa2019} y \cite{speer2007}.
\subsection{Conceptos básicos}\label{subsec:topologia_general}
\begin{Def}
   Sea $X$ un conjunto. Una \textbf{topología} sobre $X$ es una colección $\tau$ de subconjuntos de $X$ que satisface las siguientes propiedades:
   \begin{itemize}
      \item[1.] $\emptyset$ y $X$ pertenecen a $\tau$.
      \item[2.] La unión arbitraria de elementos de $\tau$ pertenece a $\tau$.
      \item[3.] La intersección finita de elementos de $\tau$ pertenece a $\tau$.
   \end{itemize}
   Los elementos de $\tau$ son llamados \textbf{conjuntos abiertos}. Si $\tau$ es una topología sobre $X$, decimos que $(X,\tau)$ es un \textbf{espacio topológico}.
\end{Def}
\begin{Obs}
   Decimos que el espacio topológico $(X,\tau)$ es \textit{finito} o \textit{infinito} dependiendo de si $X$ es finito o infinito, respectivamente. Si es clara cuál es la topología sobre $X$ con la que se está trabajando, o si no hace falta hacer referencia explícita a ella, nos referiremos por \textit{espacio topológico} simplemente a $X$. También es común referirnos a los elementos de la topología meramente como \textit{``abiertos''}. De esta forma, podemos decir que una topología sobre $X$ es una colección de abiertos tal que $\emptyset$ y $X$ son abiertos, la unión arbitraria de abiertos es abierto y la intersección finita de abiertos es abierto.
\end{Obs}
\begin{Ejm}\label{Ejm:Topologias}
   \phantom{o}\par
   \begin{itemize}
      \item Para cualquier conjunto $X$, la colección $\{\emptyset, X\}$ es una topología sobre $X$ llamada la \textbf{topología trivial}. Si $X$ está dotado de esta topología, decimos que $X$ es un \textbf{espacio trivial}.
      \item Para cualquier conjunto $X$, la colección de todos los subconjuntos de $X$, $\wp(X)$, es una topología sobre $X$, llamada \textbf{topología discreta}. Si $X$ está dotado de esta topología, decimos que $X$ es un \textbf{espacio discreto}.
      \item Tomando $X=\left\lbrace a,b,c\right\rbrace$ se tiene que $\tau_1=\left\lbrace \emptyset, \left\lbrace a\right\rbrace, \left\lbrace a,b\right\rbrace, X\right\rbrace$, $\tau_2=\{\emptyset,\{b\},\{a,b\},\allowbreak\{b,c\},X\}$ y $\tau_3=\{\emptyset,\{b\},X\}$ son topologías sobre $X$.
      \item Para cualquier conjunto $X$, la colección $\tau_{\text{cof}}=\{U\subseteq X\mid X\setminus U \text{ es finito}\}$ es una topología sobre $X$, llamada la \textbf{topología cofinita}.
      \item La colección $\tau_4=\{(-n,n)\mid n\in \mathbb{Z}^+\}\cup\{\emptyset, \mathbb{R}\}$ es una topología sobre $\mathbb{R}$. Notemos que $\bigcap_{n\in \mathbb{Z}^+}(-n,n)=\{0\}\notin \tau_4$, lo cual nos muestra que en una topología la intersección arbitraria de abiertos no siempre es un abierto.
   \end{itemize}
\end{Ejm}
\begin{Not}
    Sea $X$ un espacio topológico. Denotamos $U\subseteqab X$ para expresar que $U$ es un subconjunto abierto de $X$.
\end{Not}
\begin{Def}
   Sea $X$ un espacio topológico. Si $x\in X$ y $U\subseteqab X$ con $x\in U$, entonces decimos que $U$ es una \textbf{vecindad} de $x$. 
\end{Def}
La siguiente es una caracterización útil de los abiertos de un espacio topológico.
\begin{Prop}\label{Prop:Caracterizacion_abiertos}
   Sean $X$ un espacio topológico y $U\subseteq X$. Entonces $U\subseteqab X$ si, y solo si, para todo $x\in U$ existe una vecindad $V$ de $x$ con $V\subseteq U$.
\end{Prop}
\begin{Def}
   Sean $X$ un espacio topológico y $B\subseteq X$. Decimos que $B$ es un subconjunto \textbf{cerrado} de $X$ si $X\setminus B \subseteqab X$. Si $C\subseteq X$ es simultáneamente abierto y cerrado, decimos que $C$ es un subconjunto \textbf{clopen} de $X$. 
\end{Def}
\begin{Not}
   Sea $X$ un espacio topológico. Denotamos $B\subseteqcer X$ para expresar que $B$ es un subconjunto cerrado de $X$.
\end{Not}
En el \Cref{Ejm:Topologias}, todos los conjuntos de la forma $(-\infty,-n]\cup [n,\infty)$ con $n\in \mathbb{Z}^+$, son subconjuntos cerrados de $(\mathbb{R},\tau_4)$. En cualquier espacio topológico $X$ se tiene que $X$ y $\emptyset$ son subconjuntos clopen de $X$. En cualquier espacio discreto $Y$, todo subconjunto de $Y$ es clopen.\\\\
Como consecuencia de las leyes de De Morgan se tiene la siguiente proposición.
\begin{Prop}
   En cualquier espacio topológico la intersección arbitraria de subconjuntos cerrados es un subconjunto cerrado, y la unión finita de subconjuntos cerrados es un subconjunto cerrado.
\end{Prop}
Las nociones de \textit{bases} y \textit{subbases} nos permiten tanto construir topologías para un conjunto a partir de ciertas familias de subconjuntos, como capturar la información necesaria para reconstruir la topología de un espacio a partir de un subconjunto de abiertos con algunas propiedades especiales. Esto se formaliza a continuación. 
\begin{Def}
   Sean $X$ un conjunto y $\mathcal{B}$ una colección de subconjuntos de $X$. Decimos que $\mathcal{B}$ es \textbf{base para una topología} sobre $X$ si $\tau_\mathcal{B}:=\left\lbrace \bigcup A \mid A\subseteq \mathcal{B}\right\rbrace$ (la colección de todas las uniones de elementos de $\mathcal{B}$) es una topología sobre $X$. En este caso llamamos \textbf{abiertos básicos} de $X$ a los elementos de $\mathcal{B}$, y decimos que $\tau_\mathcal{B}$ es la \textbf{topología generada} por $\mathcal{B}$. Si $B\in\mathcal{B}$ y $x\in B$, decimos que $B$ es una \textbf{vecindad básica} de $x$.
\end{Def}
La siguiente proposición establece varias condiciones equivalentes para que una familia de subconjuntos sea base para una topología.
\begin{Prop}\label{Prop:Caracterizacion_base}
   Sean $X$ un conjunto y $\mathcal{B}$ una colección de subconjuntos de $X$. Las siguientes afirmaciones son equivalentes:
   \begin{itemize}
      \item[1.] $\mathcal{B}$ es base para una topología sobre $X$.
      \item[2.] $\mathcal{B}$ satisface las siguientes dos propiedades:
         \begin{itemize}
            \item Para todo $x\in X$ existe $B\in \mathcal{B}$ tal que $x\in B$.
            \item Dados $x\in X$ y $B_1,B_2\in \mathcal{B}$ cualesquiera, si $x\in B_1\cap B_2$, entonces existe $B_3\in \mathcal{B}$ tal que $x\in B_3$ y $B_3\subseteq B_1\cap B_2$.
         \end{itemize}
      \item[3.] $\mathcal{B}$ satisface las siguientes dos propiedades:
         \begin{itemize}
            \item $X = \bigcup \mathcal{B}$.
            \item La intersección de cualesquiera dos elementos de $\mathcal{B}$ es una unión (posiblemente vacía) de elementos de $\mathcal{B}$.
         \end{itemize}
   \end{itemize}
\end{Prop}
\begin{Obs}
   Si $(X,\tau)$ es un espacio topológico y $\mathcal{B}$ es base para una topología sobre $X$ tal que $\tau = \tau_\mathcal{B}$, se dice indistintamente que $\mathcal{B}$ es \textbf{base para} $\tau$ o que $\mathcal{B}$ es \textbf{base para} $X$, si entendemos a $X$ como espacio topológico dotado de la topología $\tau$.
\end{Obs}
Con lo anterior se tiene que $\mathcal{B}$ es base para $\tau$ si, y solo si, cualquier elemento de $\tau$ se puede escribir como unión de elementos de $\mathcal{B}$, es decir, cualquier abierto del espacio es unión de abiertos básicos. La siguiente proposición ofrece otra caracterización para que $\mathcal{B}$ sea una base para $\tau$.
\begin{Prop}\label{Prop:base_para_tau}
   Sean $(X,\tau)$ un espacio topológico y $\mathcal{B}$ una colección de subconjuntos de $X$. Se tiene que $\mathcal{B}$ es base para $\tau$ si, y solo si, $\mathcal{B}\subseteq\tau$ y se cumple la siguiente propiedad:
   \begin{itemize}
      \item Para cualquier $U\in \tau$, si $x\in U$, entonces existe $B\in \mathcal{B}$ tal que $x\in B$ y $B\subseteq U$.
   \end{itemize}
\end{Prop}
\begin{Ejm}\label{Ejm:Bases}
   \phantom{o}\par
   \begin{itemize}
      \item La colección de intervalos abiertos $\{(a,b)\mid a,b\in \mathbb{R} \text{ y } a<b\}$ es base para una topología sobre $\mathbb{R}$, llamada la \textbf{topología usual} de $\mathbb{R}$. De esta forma, en $\mathbb{R}$ con la topología usual, un subconjunto es abierto si, y solo si, es unión (posiblemente vacía) de intervalos abiertos.
      \item Para cualquier conjunto $X$, la colección $\{\{ x\}\mid x\in X\}$ es base para la topología discreta sobre $X$.
   \end{itemize}
\end{Ejm}
\begin{Def}
   Sean $X$ un conjunto y $\mathcal{S}$ una colección de subconjuntos de $X$. Decimos que $\mathcal{S}$ es \textbf{subbase para una topología} sobre $X$ si $X=\bigcup \mathcal{S}$. 
\end{Def}
\begin{Prop}
   Sean $X$ un conjunto y $\mathcal{S}$ una subbase para una topología sobre $X$. Entonces $\mathcal{B}_\mathcal{S}$, la colección de todas las intersecciones finitas no vacías de elementos de $\mathcal{S}$, es base para una topología sobre $X$.
\end{Prop}
\begin{Def}
   Sean $X$ un conjunto y $\mathcal{S}$ una subbase para una topología sobre $X$. La topología generada por $\mathcal{B}_\mathcal{S}$, denotada por $\tau_\mathcal{S}$, es llamada la \textbf{topología generada por la subbase} $\mathcal{S}$.
\end{Def}
De esta forma, una manera de construir una topología sobre $X$ a partir de una subbase $\mathcal{S}$ es tomar intersecciones finitas, y no vacías, de elementos de $\mathcal{S}$ para formar $\mathcal{B}_\mathcal{S}$, y luego, tomar uniones arbitrarias de elementos de $\mathcal{B}_\mathcal{S}$, lo cual produce la topología $\tau_\mathcal{S}$. Con lo anterior, se tiene en particular que los elementos de $\mathcal{S}$ pertenecen a $\tau_\mathcal{S}$; a los abiertos de este tipo los llamaremos \textit{abiertos subbásicos}. El procedimiento para construir $\tau_\mathcal{S}$ se ejemplifica a continuación.
\begin{Ejm}
   Sea $X=\{a,b,c,d\}$. La colección de subconjuntos de $X$ 
   $$
   \mathcal{S} = \{\{a, b\}, \{b, c\}, \{c, d\}\}
   $$
   es una subbase para una topología sobre $X$, pues
   $$
   \bigcup \mathcal{S} = \{a, b\} \cup \{b, c\} \cup \{c, d\} = \{a, b, c, d\} = X.
   $$
   Tomando las intersecciones de elementos de $\mathcal{S}$ obtenemos
   $$
   \mathcal{B}_\mathcal{S} = \{\emptyset, \{b\}, \{c\}, \{a, b\}, \{b, c\}, \{c, d\}\}.
   $$
   Tomando las uniones de elementos de $\mathcal{B}_\mathcal{S}$ cubrimos todas las uniones arbitrarias de $\mathcal{B}_\mathcal{S}$, para obtener
   $$
   \tau_\mathcal{S} = \{\emptyset, \{b\}, \{c\}, \{b, c\}, \{a, b\}, \{c, d\}, \{a, b, c\}, \{b, c, d\}, X\}.
   $$
\end{Ejm}
\begin{Obs}
   Si $(X,\tau)$ es un espacio topológico y $\mathcal{S}$ es subbase para una topología sobre $X$ tal que $\tau=\tau_\mathcal{S}$, se dice indistintamente que $\mathcal{S}$ es \textbf{subbase para} $\tau$ o que $\mathcal{S}$ es \textbf{subbase para} $X$, si entendemos a $X$ como espacio topológico dotado de la topología $\tau$.
\end{Obs}
Dos nociones que son de uso común en topología son el \textit{interior} y la \textit{clausura} de un subconjunto.
\begin{Def}\label{Def:Int_Cl}
   Sean $X$ un espacio topológico y $A\subseteq X$. El \textbf{interior} de $A$, denotado $\operatorname{Int} A$, es el conjunto $\bigcup \left\lbrace U\mid U\subseteqab X \text{ y } U\subseteq A \right\rbrace$. Los elementos de $\operatorname{Int}A$ se llaman \textbf{puntos interiores} de $A$. La \textbf{clausura} de $A$, denotada $\overline{A}$ es el conjunto \\$\bigcap \left\lbrace B \mid B\subseteqcer X \text{ y } A\subseteq B\right\rbrace$. Los elementos de $\overline{A}$ se llaman \textbf{puntos de adherencia} de $A$.
\end{Def}
De este modo, el \textit{interior} de $A$ es la unión de todos los abiertos contenidos en $A$, así, es el mayor abierto contenido en $A$, y la \textit{clausura} de $A$ es la intersección de todos los cerrados que contienen a $A$, de modo que es el menor cerrado que contiene a $A$. Además, se cumple $\operatorname{Int} A\subseteq A\subseteq \overline{A}$. 
\begin{Ejm}
   Consideremos $\mathbb{R}$ dotado de la topología usual. Sea $A=[-1,0)\cup(0,1]$. Entonces $\operatorname{Int}A=(-1,0)\cup (0,1)$ y $\overline{A}=[-1,1]$. 
\end{Ejm}
La siguiente proposición nos muestra que la clausura \textit{se comporta bien} respecto a la contenencia y la unión finita de conjuntos.
\begin{Prop}\label{Prop:Contenencia_Clausura}
   Sean $X$ un espacio topológico y $A,B\subseteq X$. Entonces se cumple lo siguiente.
   \begin{itemize}
      \item[1.] Si $A\subseteq B$ entonces $\overline{A}\subseteq\overline{B}$.
      \item[2.] $\overline{A\cup B}=\overline{A}\cup\overline{B}$.
   \end{itemize}
\end{Prop}
Usando inducción sobre el ítem 2, tenemos que para cualquier colección finita de subconjuntos de $X$ $\left\lbrace A_1,\dots,A_n\right\rbrace$ con $n\in\mathbb{Z}^+$, se cumple que $\overline{\bigcup_{i=1}^{n}A_i}=\bigcup_{i=1}^{n}\overline{A_i}$. Sin embargo, la igualdad no siempre se tiene para uniones infinitas, como lo muestra el siguiente ejemplo.
\begin{Ejm}
   En $\mathbb{R}$ con la topología usual se tiene que
   $$
   \begin{aligned}
      \overline{\bigcup_{i=1}^{\infty}\left(-1+\frac{1}{i},1-\frac{1}{i}\right)} &= \overline{(-1,1)}\\
         &=[-1,1]
   \end{aligned}
   $$
   mientras que
   $$
   \begin{aligned}
      \bigcup_{i=1}^{\infty}\overline{\left( -1+\frac{1}{i},1-\frac{1}{i}\right)} &= \bigcup_{i=1}^{\infty}\left[-1+\frac{1}{i},1-\frac{1}{i}\right]\\
          &= (-1,1).
   \end{aligned}
   $$
\end{Ejm}
En general, solo se tiene un contenencia:
\begin{Prop}\label{Prop:Clausura_uniones_finitas}
   Sea $X$ un espacio topológico. Para cualquier colección $\left( U_i\right)_{i\in I}$ de subconjuntos abiertos de $X$ se tiene $\bigcup_{i\in I}\overline{U_i}\subseteq \overline{\bigcup_{i\in I}U_i}$.
\end{Prop}
La \Cref{Def:Int_Cl} junto con las siguiente proposición nos permite caracterizar los subconjuntos abiertos y cerrados de un espacio topológico.
\begin{Prop}\label{Carac_abierto_cerrado}
   Sean $X$ un espacio topológico y $A\subseteq X$. Se tiene que $A\subseteqab X$ si, y solo si, $A=\operatorname{Int}A$, y $A\subseteqcer X$ si, y solo si, $A=\overline{A}$. 
\end{Prop}
La siguiente proposición nos da algunas caracterizaciones para que un elemento de un espacio topológico sea un punto de adherencia de un subconjunto.
\begin{Prop}\label{Prop:Caracterizacion_punto_adherencia}
   Sean $X$ un espacio topológico, $A\subseteq X$ y $x\in X$.
   \begin{itemize}
      \item[1.] Se tiene que $x\in \overline{A}$ si, y solo si, cualquier vecindad de $x$ interseca a $A$.
      \item[2.] Si la topología de $X$ está dada por una base, entonces $x\in\overline{A}$ si, y solo si, cualquier vecindad básica de $x$ interseca a $A$. 
   \end{itemize}
\end{Prop}
Cuando decimos que un conjunto $B$ interseca a $A$ nos referimos a que $B\cap A\neq \emptyset$.
%EJEMPLOS
\begin{Ejm}\label{Ejm:ejemplo_denso}
   Consideremos $\mathbb{R}$ con la topología usual. Sabemos que la colección de todos los intervalos abiertos es una base para este espacio topológico (\Cref{Ejm:Bases}). Para cualquier $x\in\mathbb{R}$ y cualquier intervalo $(a,b)$ con $a,b\in\mathbb{R}$ y $a<b$, tal que $x\in(a,b)$, como consecuencia de la propiedad arquimediana de los números reales se tiene que existe $q\in\mathbb{Q}$ tal que $q\in(a,b)$ (puede verse \cite{apostol1974} para una prueba de este hecho), es decir, $(a,b)$ interseca a $\mathbb{Q}$. Por la \Cref{Prop:Caracterizacion_punto_adherencia} se sigue que $x\in\overline{\mathbb{Q}}$ y por tanto $\mathbb{R}\subseteq\overline{\mathbb{Q}}$. Como además, por la definición de clausura se tiene $\overline{\mathbb{Q}}\subseteq\mathbb{R}$, entonces $\overline{\mathbb{Q}}=\mathbb{R}$.
\end{Ejm}
Los subconjuntos que, como en el \Cref{Ejm:ejemplo_denso}, al tomarles su clausura son iguales a todo el espacio que los contiene, reciben un nombre especial.
\begin{Def}
   Sean $X$ un espacio topológico y $A\subseteq X$. Decimos que $A$ es un subconjunto \textbf{denso} de $X$ si $\overline{A}=X$. 
\end{Def}
La siguiente proposición da algunas caracterizaciones para que un subconjunto de un espacio sea denso.
\begin{Prop}\label{Prop:Caracterizacion_denso}
   Sean $X$ un espacio topológico, $A\subseteq X$ y $\mathcal{B}$ una base para $X$. Las siguientes afirmaciones son equivalentes:
   \begin{itemize}
      \item[1.] $A$ es un subconjunto denso de $X$.
      \item[2.] $\operatorname{Int}(X\setminus A)=\emptyset$. 
      \item[3.] Para todo $U\subseteqab X$, con $U\ne\emptyset$, se tiene $A\cap U\ne\emptyset$.
      \item[4.] Para todo $B\in \mathcal{B}$, con $B\ne \emptyset$, se tiene $A\cap B\ne \emptyset$.
   \end{itemize}
\end{Prop}
\begin{Def}
   Sean $X$ un espacio topológico y $A$ un subconjunto denso de $X$. Decimos que $A$ es un subconjunto \textbf{denso minimal} de $X$, si para cualquier subconjunto denso $B$ de $X$, la inclusión $B\subseteq A$ implica $B=A$.   
\end{Def}
Equivalentemente, $A$ es un subconjunto denso minimal de $X$ si, y solo si, ningún subconjunto denso de $X$ está propiamente contenido en $A$.
\begin{Ejm}
   Consideremos a $\mathbb{R}$ dotado de la topología $\tau=\{(-n,n)\mid n\in \mathbb{Z}^+\}\cup\{\emptyset, \mathbb{R}\}$. Cualquier abierto no vacío de $(\mathbb{R},\tau)$ contiene a $0$, de modo que, por el ítem 3. de la \Cref{Prop:Caracterizacion_denso}, $\{0\}$ es un subconjunto denso de $(\mathbb{R},\tau)$. Si $B$ es un subconjunto denso de $X$ entonces, por ejemplo, $(-1,1)\cap B\ne \emptyset$, así que $B\ne \emptyset$, y por tanto $B$ no está propiamente contenido en $\{0\}$. Se sigue así que $\{0\}$ es un subconjunto denso minimal de $(\mathbb{R},\tau)$. 
\end{Ejm}
Dos propiedades de gran interés al trabajar con espacios topológicos son la \textit{conexidad} y la \textit{compacidad}.
\begin{Def}
   Sea $X$ un espacio topológico. Una \textbf{separación} de $X$ es un par $\{U,V\}$ de subconjuntos abiertos de $X$, no vacíos y disyuntos, tales que $U\cup V=X$. El espacio $X$ se dice \textbf{conexo} si no existe ninguna separación de $X$. En caso contrario decimos que $X$ es \textbf{disconexo}.
\end{Def}
Una forma equivalente de \textit{conexidad} se presenta en la siguiente caracterización.
\begin{Prop}
   Un espacio topológico $X$ es conexo si, y solo si, los únicos subconjuntos clopen de $X$ son $\emptyset$ y $X$. 
\end{Prop}
\begin{Ejm}
   \phantom{o}
   \begin{itemize}
      \item Cualquier espacio trivial es conexo.
      \item Cualquier espacio discreto $X$ con dos o más elementos es disconexo, pues si $a\in X$ entonces $\{a\}$ es un subconjunto clopen de $X$ distinto de $\emptyset$ y de $X$.
   \end{itemize}
\end{Ejm}
Para definir un espacio topológico \textit{compacto}, primero presentamos la definición de \textit{cubrimiento abierto} y \textit{subcubrimiento finito}.
\begin{Def}
   Sea $X$ un espacio topológico. Un \textbf{cubrimiento abierto} de $X$ es una familia $\left\lbrace U_i\right\rbrace_{i\in I}$ con $U_i\subseteqab X$ para todo $i\in I$, tal que $\bigcup_{i\in I}U_i=X$. Si $\left\lbrace U_i\right\rbrace_{i\in I}$ es un cubrimiento abierto de $X$, un \textbf{subcubrimiento finito} de $\left\lbrace U_i\right\rbrace_{i\in I}$ es una familia $\left\lbrace U_{i_j}\right\rbrace_{j=1}^{n}$ con $n\in \mathbb{Z}^+$ e $i_1,\dots,i_n\in I$, tal que $\bigcup_{j=1}^{n}U_{i_j}=X$.
\end{Def}
Es de precisar que todo espacio topológico $X$ admite al menos un cubrimiento abierto: la familia cuyo único elemento es $X$ (o toda la topología).
\begin{Def}
   Sea $X$ un espacio topológico. Decimos que $X$ es \textbf{compacto} si todo cubrimiento abierto de $X$ tiene un subcubrimiento finito.
\end{Def}
\begin{Ejm}
   \phantom{o}
   \begin{itemize}
      \item Todo espacio topológico finito $X$ es compacto, pues su número de subconjuntos (y por tanto, también su número de abiertos) es finito, con lo cual todo cubrimiento abierto de $X$ es en sí mismo un subcubrimiento finito.
      \item Todo espacio discreto infinito $X$ no es compacto, pues $\left\lbrace \{x\}\mid x\in X\right\rbrace$ es un cubrimiento abierto de $X$ que no tiene un subcubrimiento finito.
      \item El espacio $\mathbb{R}$ con la topología usual no es compacto, pues $\{(-n,n)\mid n\in \mathbb{Z}^+\}$ es un cubrimiento abierto de $\mathbb{R}$ que no tiene un subcubrimiento finito.
   \end{itemize}
\end{Ejm}
Ahora enunciamos algunas propiedades conocidas como \textit{propiedades de separación}.
\begin{Def}
   Sea $X$ un espacio topológico.
   \begin{itemize}
      \item Decimos que $X$ es $\bm{T_0}$ si para cualesquiera $x,y\in X$, con $x\ne y$, existe una vecindad de $x$ que no contiene a $y$, o una vecindad de $y$ que no contiene a $x$. 
      \item Decimos que $X$ es $\bm{T_1}$ si para cualesquiera $x,y\in X$, con $x\ne y$, existen una vecindad de $x$ que no contiene a $y$, y una vecidad de $y$ que no contiene a $x$.
      \item Decimos que $X$ es \textbf{Hausdorff} (o $\bm{T_2}$) si para cualesquiera $x,y\in X$, con $x\ne y$, existen una vecindad $U_x$ de $x$ y una vecindad $U_y$ de $y$ tales que $U_x\cap U_y=\emptyset$.
   \end{itemize}
\end{Def}
Es claro que todo espacio Hausdorff es $T_1$, y a su vez, todo espacio $T_1$ es $T_0$.
\begin{Ejm}
   \phantom{o}
   \begin{itemize}
      \item Cualquier espacio trivial con dos o más elementos no es $T_0$.
      \item El conjunto $X=\{a,b\}$ dotado de la topología $\tau=\{\emptyset,\{a\},X\}$ es $T_0$ pero no $T_1$.
      \item Cualquier espacio infinito dotado de la topología cofinita es $T_1$ pero no Hausdorff.
      \item Cualquier espacio discreto es Hausdorff.
   \end{itemize}
\end{Ejm}
Ahora vamos a definir las funciones que nos permiten conectar espacios topológicos entre sí.
\begin{Def}
   Sean $(X,\tau_X), (Y,\tau_Y)$ espacios topológicos y $f:X\to Y$ una función.
   \begin{itemize}
      \item[1.] Decimos que $f$ es \textbf{continua} si para todo $V\in\tau_Y$ se tiene $f^{-1}(V)\in \tau_X$.
      \item[2.] Decimos que $f$ es \textbf{abierta} si para todo $U\in \tau_X$ se tiene $f(U)\in \tau_Y$.
      \item[3.] Decimos que $f$ es \textbf{cerrada} si para todo subconjunto cerrado $B$ de $(X,\tau_X)$ se tiene que $f(B)$ es un subconjunto cerrado de $(Y,\tau_Y)$.
   \end{itemize}
\end{Def}
De esta forma, solemos decir que una función continua es aquella que ``devuelve'' abiertos en abiertos por la imagen inversa, una función abierta es aquella que envía abiertos en abiertos por la imagen directa, y una función cerrada es aquella que envía cerrados en cerrados por la imagen directa.
\begin{Ejm}\label{Ejm:funciones_continuas}
   \phantom{o}
   \begin{itemize}
      \item Cualquier función que tenga como dominio un espacio discreto es continua.
      \item Cualquier función que tenga como codominio un espacio trivial es continua.
      \item Tomemos $X=\mathbb{N}$ dotado de la topología discreta y $Y=\mathbb{N}$ dotado de la topología trivial. Tomemos $f:X\to Y$ definida mediante $f(n)=n$ para todo $n\in X$. Por los ítems anteriores se tiene que $f$ es continua. Además $f$ tiene inversa $f^{-1}:Y\to X$ que satisface $f^{-1}(n)=n$ para todo $n\in Y$. Sin embargo, $f^{-1}$ no es continua, pues $\{0\}\subseteqab X$ pero $f^{-1}(\{0\})=\{0\}$ no es un subconjunto abierto de $Y$.   
      \item Sea $X = \{1, 2, 3\}$ y consideremos la función $f: X \to X$ dada por $f(1) = 1$, $f(2) = 1$ y $f(3) = 3$. Dotamos al dominio de la topología $\tau_1 = \{\emptyset, \{1, 2\}, \{3\}, X\}$ y al codominio de la topología $\tau_2 = \{\emptyset, \{1\}, \{3\}, \{1, 3\}, X\}$. Es fácil verificar con cada elemento de $\tau_1$ que $f$ es abierta, sin embargo la imagen del subconjunto cerrado $\{3\}$ de $(X,\tau_1)$ no es cerrada en $(X,\tau_2)$, de modo que $f$ no es cerrada.
      \item Cualquier función constante de $\mathbb{R}$ dotado de la topología usual en sí mismo es cerrada y no es abierta.
   \end{itemize}
\end{Ejm}
La siguiente proposición nos da algunas caracterizaciones de la continuidad de una función.
\begin{Prop}\label{Prop:Caracterizacion_continuidad}
   Sean $(X,\tau_X), (Y,\tau_Y)$ espacios topológicos y $f:X\to Y$ una función. Las siguientes proposiciones son equivalentes:
   \begin{itemize}
      \item[1.] $f$ es continua.
      \item[2.] Para todo subconjunto cerrado $B$ de $(Y,\tau_Y)$ se tiene que $f^{-1}(B)$ es un subconjunto cerrado de $(X,\tau_X)$.
      \item[3.] Para todo $A\subseteq X$ se tiene $f(\overline{A})\subseteq\overline{f(A)}$.
      \item[4.] Para cualquier $x\in X$ y cualquier vecindad $V$ de $f(x)$, existe una vecindad $U$ de $x$ tal que $f(U)\subseteq V$.
   \end{itemize}
\end{Prop}
La siguiente caracterización de las funciones cerradas será usada más adelante.
\begin{Prop}
   Sean $(X,\tau_X), (Y,\tau_Y)$ espacios topológicos y $f:X\to Y$ una función. Entonces $f$ es cerrada si, y solo si, $\overline{f(A)}\subseteq f(\overline{A})$ para todo $A\subseteq X$. 
\end{Prop}
En el \Cref{Ejm:funciones_continuas} vimos que puede tenerse una función $f$ que sea continua e invertible, tal que $f^{-1}$ no sea continua. Las funciones continuas que tienen inversa continua reciben un nombre especial.
\begin{Def}
   Sean $X$ y $Y$ espacios topológicos. Si $f:X\to Y$ es una función continua e invertible, tal que $f^{-1}:Y\to X$ es continua, decimos que $f$ es un \textbf{homeomorfismo}. Si existe un homeomorfismo de $X$ en $Y$, decimos que $X$ y $Y$ son \textbf{homeomorfos}. 
\end{Def}
En topología dos espacios homeomorfos son indistinguibles para la teoría; pueden considerarse como dos representaciones distintas del mismo espacio.
\begin{Ejm}\label{Ejm:homeomorfismos}
   \phantom{o}
   \begin{itemize}
      \item La función identidad es un homeomorfismo de cualquier espacio en sí mismo.
      \item Si $(X,\tau_X)$ y $(Y,\tau_Y)$ son espacios discretos con $|X|=|Y|$, entonces cualquier biyección de $X$ en $Y$ es un homeomorfismo.
      \item Si $(X,\tau_X)$ y $(Y,\tau_Y)$ son espacios triviales con $|X|=|Y|$, entonces cualquier biyección de $X$ en $Y$ es un homeomorfismo.
      \item Sean $X = \{1, 2, 3\}$ y $Y = \{a, b, c\}$. Dotamos a $X$ de la topología $\tau_X = \{\emptyset, \{1\},\allowbreak \{2, 3\},\allowbreak X\}$ y a $Y$ de la topología $\tau_Y = \{\emptyset, \{c\}, \{a, b\}, Y\}$. Puede verse que la función $f: X \to Y$ dada por $f(1) = c$, $f(2) = a$ y $f(3) = b$ es un homeomorfismo.
   \end{itemize}
\end{Ejm}
Una propiedad que se conserva bajo homeomorfismos se dice \textit{topológica}. Más formalmente:
\begin{Def}
   Sean $X$ y $Y$ espacios topológicos y $P$ una propiedad. Decimos que $P$ es una \textbf{propiedad topológica} si el hecho de que $X$ tenga la propiedad $P$ y $X$ sea homeomorfo a $Y$, impica que $Y$ tiene la propiedad $P$.
\end{Def}
Puede probarse que ser conexo, compacto, $T_0$, $T_1$ o Hausdorff son ejemplos de propiedades topológicas. Las propiedades topológicas son bastante útiles para demostrar que dos espacios no son homeomorfos. Por ejemplo, si un espacio $X$ es conexo y un espacio $Y$ no lo es, podemos afirmar que $X$ y $Y$ no son homeomorfos.

\subsection{Espacios de Alexandroff}\label{subsec:Alexandroff}
%Tensión discreto vs continuo y finito vs infinito
%Qué se gana con los espacios de Alexandroff frente a los espacios finitos?
Ahora procedemos a dar un pequeño vistazo a la teoría de \textit{espacios de Alexandroff}.
\begin{Def}
   Una topología $\tau$ sobre un conjunto $X$ se dice de \textbf{Alexandroff} si para cualquier colección $\left\{ U_i\right\}_{i\in I}$ de elementos de $\tau$ se tiene $\bigcap_{i\in I}U_i \in\tau$. En este caso decimos que $(X,\tau)$ es un \textbf{espacio de Alexandroff}.
\end{Def}
De esta manera, un espacio topológico es de \textit{Alexandroff} si la intersección arbitraria de abiertos es un abierto. En el \Cref{Ejm:Topologias} vimos un espacio topológico en el que no se tiene esta propiedad.\\\\
Pável Serguéyevich Alexandroff (Bogorodsk, 7 de mayo de 1896 - Moscú, 16 de noviembre de 1982) introduce en 1937, en \cite{alexandroff1937}, los espacios que ahora llevan su apellido, denominándolos originalmente como \textit{Diskrete Räume}, o traducido del alemán al español, \textit{espacios discretos} (no confundir con los espacios introducidos en el \Cref{Ejm:Topologias}). El objetivo original de Alexandroff se centró en construir una teoría de homología para sus \textit{Diskrete Räume}. Esta motivación se enmarcaba principalmente en el campo de la entonces llamada \textit{topología combinatoria}, precursora de lo que ahora conocemos como \textit{topología algebraica}. De acuerdo a \cite{arenas1999}, los espacios de Alexandroff cobraron interés en la década de 1980 debido a su uso dentro de la \textit{topología digital}, rama de la topología que trata principalmente con \textit{imágenes digitales}, entendidas como vectores discretos en dos o más dimensiones, y en la cual problemas alrededor del procesamiento de imágenes, tales como esqueletización, seguimiento de bordes, relleno de contornos o conteo de objetos, son de especial interés, conforme a \cite{kong1989}. Lo anterior desemboca, así mismo, en una amplia gama de aplicaciones. Por ejemplo, en \cite{khalimsky1990} se nombran el estudio de gráficas computarizadas, análisis de patrones, diseño robótico o la tomografía computarizada, de amplio uso en la medicina actual. La razón del éxito que tienen los espacios de Alexandroff en este campo, viene como fruto de su capacidad para abstraer y generalizar las propiedades, que tienen los espacios topológicos finitos, que son relevantes en los problemas nombrados. Siguiendo a \cite{rubiano2013} vemos que esto, a su vez, permite su uso en la teoría de las ciencias de la computación, donde limitaciones naturales, como las concernientes al uso de memoria, hacen que se tenga una prevalencia inicial por espacios topológicos finitos, y en general por aquellos que no son de Hausdorff, que históricamente han estado más ligados a la rama del análisis y al estudio del continuo.  Allende a su uso en las matemáticas aplicadas, el estudio de los espacios de Alexandroff también muestra influencia en ramas como las matemáticas discretas, donde permite uno de los primeros vínculos de esta área con la topología; la teoría de retículos, donde juegan un papel importante en el estudio del retículo de las topologías; o la teoría de números, donde, desde los espacios funcionales de Alexandroff, también conocidos como topologías primales, se presenta un nuevo punto de ataque a la célebre Conjetura de Collatz (ver \cite[Capítulo 5]{mancero2023}).\\\\
Como consecuencia de las leyes de De Morgan se tiene la siguiente equivalencia.
\begin{Prop}\label{Prop:Uniones_cerrados_Alexandroff}
   Un espacio topológico es de Alexandroff si, y solo si, la unión arbitraria de subconjuntos cerrados es un subconjunto cerrado.
\end{Prop}
Ahora damos ejemplos de espacios topológicos de Alexandroff.
\begin{Ejm}
   \phantom{o}
   \begin{itemize}
      \item Cualquier espacio finito es de Alexandroff.
      \item Cualquier espacio con topología finita es de Alexandroff, en particular los espacios triviales.
      \item Cualquier espacio discreto es de Alexandroff.
      \item $\mathbb{R}$ junto con la topología generada por la familia de intervalos $\left\lbrace [x,\infty)\mid x\in \mathbb{R}\right\rbrace$, llamada la topología de colas (a derecha) abiertas y cerradas, es un espacio de Alexandroff.
   \end{itemize}
\end{Ejm}
\begin{Def}
   Sean $X$ un espacio topológico y $x\in X$. Decimos que una vecindad $V$ de $x$ es una \textbf{vecindad mínima} de $x$ si para cualquier vecidad $U$ de $x$ se tiene $V\subseteq U$. Igualmente, decimos que una base $\mathcal{B}$ de $X$ es una \textbf{base mínima}, si para cualquier base $\mathcal{B}'$ de $X$ se tiene $\mathcal{B}\subseteq\mathcal{B}'$.
\end{Def}
En la literatura de referencia para espacios de Alexandroff, como por ejemplo \cite{arenas1999} y \cite{speer2007}, se suelen usar las palabras vecindad \textit{minimal} (\textit{minimal neighborhood}, en inglés), en lugar de vecindad \textit{mínima}, y base \textit{minimal} (\textit{minimal base}, en inglés), en lugar de base \textit{mínima}. Sin embargo, optamos por tomar las acepciones clásicas que en teoría de conjuntos ordenados se hacen de \textit{elemento mínimo} (\textit{least element}, en inglés) y \textit{elemento minimal} (\textit{minimal element}, en inglés). Por ejemplo, \cite{bloch2011} define:
\begin{Def}
   Sean $(A,\preccurlyeq)$ un conjunto parcialmente ordenado, y $a\in A$. Se dice que $a$ es un \textbf{elemento mínimo} de $A$ si $a\preccurlyeq x$ para todo $x\in A$. Se dice que $a$ es un \textbf{elemento minimal} de $A$ si no existe $x\in A$ tal que $x\preccurlyeq a$ y $a\ne x$. 
\end{Def}
Esta decisión se toma porque, a lo largo del presente trabajo, nos reservamos el término \textit{minimal} para usarlo en este sentido, como lo hacemos en las Definiciones \ref{Def:cortes} y \ref{Def:denso_minimal}, y más adelante, en las Proposiciones \ref{Prop:Conexo_Alexandroff} y \ref{Prop:cardinal_densos_minimales}.\\\\ 
Los principales órdenes parciales que acá tenemos en cuenta son aquellos que surgen de considerar conjuntos, espacios topológicos, o colecciones de grafos, ordenados por la contenencia. Nótese que en cualquier conjunto parcialmente ordenado $(A,\preccurlyeq)$, si $a$ es un mínimo de $A$, entonces $a$ es un elemento minimal de $A$, ya que para cualquier $x\in A$ con $x\preccurlyeq a$, por la minimalidad de $a$ se tiene que $a\preccurlyeq x$, y por la antisimetría de $\preccurlyeq$ se sigue que $x=a$. Sin embargo, ser elemento minimal no implica ser elemento mínimo. Por nombrar un caso, en la colección de subconjuntos no vacíos de $\mathbb{R}$, ordenada por la contenencia, se tiene que $\{0\}$ es un elemento minimal, pues ningún subconjunto no vacío de $\mathbb{R}$ está propiamente contenido en él. Sin embargo, $\{0\}$ no es un elemento mínimo, pues, por ejemplo, $\{0\}\nsubseteq \{1\}$. Puede verse que cualquier colección de conjuntos que contenga a $\emptyset$, oredenada por la contenencia, tiene a $\emptyset$ como elemento mínimo. A continuación mostramos algunas vecindades mínimas y bases mínimas en un par de espacios de Alexandroff.
\begin{Ejm}
   En cualquier espacio discreto $X$, se tiene que $\{x\}$ es una vecindad mínima de cualquier $x\in X$, y el conjunto $\{\{x\}\mid x\in X\}$ es una base mínima de $X$. En $\mathbb{R}$ dotado de la topología de colas (a derecha) abiertas y cerradas, $[x,\infty)$ es una vecindad mínima de $x$, y el conjunto $\{[x,\infty)\mid x\in \mathbb{R}\}$ es una base mínima de este espacio. 
\end{Ejm}
Una propiedad clave de los espacios de Alexandroff es que cada punto tiene una vecindad mínima, y cada espacio tiene una base mínima. Más precisamente:
\begin{Prop}\label{Prop:Alexandroff_discreto}
   Sea $X$ un espacio de Alexandroff. Para cada $x\in X$, la intersección de todas las vecindades de $x$ es una vecindad mínima de $x$, que denotamos por $U_x$. La colección $\mathcal{B}_{\min}=\left\lbrace U_x\mid x\in X\right\rbrace$ es una base mínima de $X$.
\end{Prop}
Usando las vecindades mínimas podemos caracterizar los espacios de Alexandroff discretos de la siguiente manera.
\begin{Prop}
   Sea $X$ un espacio de Alexandroff. Entonces $X$ es discreto si, y solo si, $U_x=\{x\}$ para todo $x\in X$. 
\end{Prop}
Contar con una subbase de un espacio topológico nos permite ``ahorrar trabajo'' para conocer si el espacio es de Alexandroff, pues para esto basta determinar si la intersección arbitraria de abiertos subbásicos es abierta. Más precisamente:
\begin{Prop}\label{Prop:Alexandroff_con_subbase}
   Sean $X$ un espacio topológico y $\mathcal{S}$ una subbase de $X$. Se tiene que $X$ es un espacio de Alexandroff si, y solo si, la intersección arbitraria de elementos de $\mathcal{S}$ es un abierto de $X$.
\end{Prop}
\begin{proof}
   $(\Rightarrow)$ Se sigue del hecho de que los elementos de $\mathcal{S}$ son abiertos de $X$.\\
   $(\Leftarrow)$ Supongamos que la intersección arbitraria de elementos de $\mathcal{S}$ es un abierto de $X$. Sea $\left\{ U_i\right\}_{i\in I}$ una colección de abiertos de $X$ y llamemos $U:=\bigcap_{i\in I}U_i$. Usamos la \Cref{Prop:Caracterizacion_abiertos} para ver que $U \subseteqab X$. Sea $x\in U$. Para cada $i\in I$ tenemos que $x\in U_i$ y $U_i=\bigcup_{B\in A_i}B$ para algún $A_i\subseteq \mathcal{B}_{\mathcal{S}}$. De este modo $x\in B_i$ para algún $B_i\in A_i$. Como en particular, $B_i\in \mathcal{B}_\mathcal{S}$ entonces existen $n_i\in \mathbb{Z}^+$ y $S_{i_j}\in \mathcal{S}$ para todo $j\in\left\lbrace 1,\dots ,n_i\right\rbrace$ tales que $B_i=\bigcap_{j=1}^{n_i}S_{i_j}$. Llamemos $V:=\bigcap_{i\in I}B_i=\bigcap_{i\in I}\left( \bigcap_{j=1}^{n_i}S_{i_j}\right)$. Como $V$ es una intersección de elementos de $\mathcal{S}$, se tiene que $V\subseteqab X$, por hipótesis. Como $x\in B_i$ para todo $i\in I$, entonces $x\in V$. Además, para todo $i\in I$, como $B_i\in A_i$, entonces $B_i\subseteq \bigcup_{B\in A_i}B=U_i$, luego $V=\bigcap_{i\in I}B_i\subseteq \bigcap_{i\in I}U_i=U$. Con esto, $U\subseteqab X$ y por tanto, $X$ es un espacio de Alexandroff.
\end{proof}
Como una base de un espacio tolpológico es en particular una subbase, la \Cref{Prop:Alexandroff_con_subbase} también puede aplicarse cuando conocemos una base del espacio en cuestión.\\\\
En un espacio de Alexandroff sí se tiene que la clausura ``respeta'' las uniones arbitrarias de subconjuntos, como lo muestra la siguiente proposición.
\begin{Prop}\label{Prop:Alexandroff_clausura}
   Sean $X$ un espacio de Alexandroff y $\left\{ A_i\right\}_{i\in I}$ una colección cualquiera de subconjuntos de $X$. Entonces $\overline{\bigcup_{i\in I}A_i}=\bigcup_{i\in I}\overline{A_i}$. 
\end{Prop}
\begin{proof}
   Por la \Cref{Prop:Clausura_uniones_finitas} se tiene $\bigcup_{i\in I}\overline{A_i}\subseteq \overline{\bigcup_{i\in I}A_i}$. Ahora, como $\overline{A_i}$ es cerrado para todo $i\in I$, y como $X$ es de Alexandroff, por la \Cref{Prop:Uniones_cerrados_Alexandroff} se tiene que $\bigcup_{i\in I}\overline{A_i}$ es cerrado. Además, para todo $i\in I$ se tiene $A_i\subseteq \overline{A_i}$, con lo cual $\bigcup_{i\in I}A_i\subseteq \bigcup_{i\in I}\overline{A_i}$. Por la definición de clausura de un conjunto se sigue que $\overline{\bigcup_{i\in I}A_i}\subseteq \bigcup_{i\in I}\overline{A_i}$, lo cual completa la prueba.
\end{proof}
Las vecindades y bases mínimas de un espacio de Alexandroff nos dan una condición suficiente para que el espacio sea disconexo, como lo muestra la siguiente proposición.
\begin{Prop}\label{Prop:Conexo_Alexandroff}
   Sea $X$ un espacio de Alexandroff. Si existe $x\in X$ tal que:
   \begin{itemize}
      \item $U_x$ es minimal en $\mathcal{B}_{\min}$, es decir, para todo $y\in X$, si $U_y\subseteq U_x$ entonces $U_x=U_y$, y,
      \item $U_x$ es maximal en $\mathcal{B}_{\min}$, es decir, para todo $y\in X$, si $U_x\subseteq U_y$ entonces $U_x=U_y$,
   \end{itemize}
   entonces $U_x$ es un clopen de $X$. Por lo tanto, si $U_x\ne X$, entonces $X$ es disconexo.
\end{Prop}
\begin{proof}
   Supongamos que $x\in X$ es tal que $U_x$ es tanto maximal como minimal en $\mathcal{B}_{\operatorname{min}}$. Ya sabemos que $U_x$ es abierto. Para ver que $U_x$ es cerrado, por la \Cref{Prop:Carac_abierto_cerrado} basta ver que $\overline{U_x}\subseteq U_x$. Sea $y\in\overline{U_x}$. Como, por la \Cref{Prop:Alexandroff_clausura}, tenemos $\overline{U_x}=\overline{\bigcup_{z\in U_x}\{z\}}=\bigcup_{z\in U_x}\overline{\{z\}}$, entonces $y\in\overline{\{z\}}$ para algún $z\in U_x$. Como $y\in U_y\subseteqab X$ entonces $U_y\cap\{z\}\ne\emptyset$, es decir, $z\in U_y$, luego $U_z\subseteq U_y$. Como $z\in U_x\subseteqab X$, se tiene $U_z\subseteq U_x$. Por la maximalidad de $U_x$ en $\mathcal{B}_{\operatorname{min}}$, se tiene $U_z=U_x$. Así, $U_x\subseteq U_y$. Por la minimalidad de $U_x$ en $\mathcal{B}_{\operatorname{min}}$, se tiene $U_x=U_y$. Como $y\in U_y$, entonces $y\in U_x$ y obtenemos $\overline{U_x}\subseteq U_x$. Por lo tanto, $U_x$ es un clopen de $X$. Como $x\in U_x$, entonces $U_x\ne\emptyset$. Así, si $U_x\ne X$, por la \Cref{Prop:Caracterizacion_conexo}, se sigue que $X$ es disconexo.
\end{proof}
{\color[HTML]{188539}
En \cite{jafarian2013}, el enunciado de la proposición (Proposición 6.4) no precisa la condición $U_x\ne X$ para concluir que $X$ es disconexo. Sin embargo esto es necesario. Por ejemplo, tomando $X=\{x\}$ con la topología trivial, tenemos $\mathcal{B}_{\operatorname{min}}=\left\lbrace U_x\right\rbrace=\left\lbrace\{x\}\right\rbrace$, de modo que $U_x$ es minimal y maximal en $\mathcal{B}_{\operatorname{min}}$. Sin embargo, $X$ es conexo por ser trivial.}\\\\
La siguiente proposición nos presenta algunas características de las propiedades de separación en los espacios de Alexandroff.
\begin{Prop}\label{Prop:Separacion_Alexandroff}
   Sea $X$ un espacio de Alexandroff.
   \begin{itemize}
      \item X es $T_0$ si, y solo si, $U_x=U_y$ implica $x=y$ para cualesquiera $x,y\in X$.
      \item X es $T_1$ si, y solo si, $X$ es discreto.
   \end{itemize}
\end{Prop}
El siguiente es un hecho bastante importante y útil de los espacios de Alexandroff.
\begin{Prop}
   Sean $X$ y $Y$ espacios topológicos. Si $X$ es un espacio de Alexandroff y $X$ y $Y$ son homeomorfos, entonces $Y$ es también un espacio de Alexandroff. 
\end{Prop}
En otras palabras, ser de Alexandroff es una propiedad topológica.\\\\
Por último, enunciamos y demostramos un resultado que será utilizado más adelante.
\begin{Prop}\label{Prop:cardinal_densos_minimales}
   Sea $X$ un espacio de Alexandroff. Tomemos 
   $$
   \mathcal{B}:=\left\lbrace U_x\mid x\in X \text{ y }U_x\text{ es minimal en }\mathcal{B}_{\min}\right\rbrace.
   $$
   Si $A$ es un subconjunto denso minimal de $X$ entonces existe una función sobreyectiva $f:\mathcal{B}\to A$ tal que $f\left( U_x\right)\in U_x$ para todo $U_x\in \mathcal{B}$. En particular,
         $$
         A\subseteq\left\lbrace x\in X\mid U_x \text{ es minimal en }\mathcal{B}_{\min}\right\rbrace.
         $$
   Además, si $A$ y $A'$ son subconjuntos densos minimales de $X$ entonces $|A|=|A'|$.
\end{Prop}
\begin{proof}
   Sea $A$ un subconjunto denso minimal de $X$. Sea $U_x\in\mathcal{B}$. Como $x\in U_x\subseteqab X$ y $x\in X=\overline{A}$, entonces $U_x\cap A\ne\emptyset$, de modo que existe $x^*\in U_x\cap A$. Sea $y\in U_x\cap A$. Como $A\setminus\{y\}\subsetneq A$, entonces $A\setminus\{y\}$ no es denso en $X$, con lo cual existen $z\in X$ y $V$ vecindad de $z$ con $V\cap\left( A\setminus\{y\}\right)=\emptyset$. Pero $z\in X=\overline{A}$, luego $V\cap A\ne\emptyset$, con lo cual $V\cap A=\{y\}$. Como $y\in U_x\subseteqab X$, entonces $U_y\subseteq U_x$; como $U_x$ es minimal en $\mathcal{B}_{\min}$, entonces $U_y=U_x$. Notemos que $y\in V\subseteqab X$, luego $U_x=U_y\subseteq V$. Así, $x^*\in U_x\cap A\subseteq V\cap A=\{y\}$, y por tanto $y=x^*$. Obtenemos así que $U_x\cap A=\{x^*\}$. Entonces definimos $f:\mathcal{B}\to A$ mediante $f(U_x)=x^*$ para cada $U_x\in\mathcal{B}$. Veamos que $f$ es sobreyectiva. Sea $a\in A$. Como $A\setminus\{a\}\subsetneq A$, se tiene que $A\setminus\{a\}$ no es denso en $X$, de modo que existen $z\in X$ y $V$ vecindad de $z$, con $V\cap\left( A\setminus\{a\}\right)=\emptyset$. Pero $z\in X=\overline{A}$, luego $V\cap A\ne\emptyset$, de modo que $V\cap A=\{a\}$. Como $a\in V$ entonces $U_a\subseteq V$. Sea $U_x\in\mathcal{B}_{\min}$ con $U_x\subseteq U_a$. Tenemos $x\in U_x\subseteqab X$ y $x\in X=\overline{A}$, luego $U_x\cap A\ne\emptyset$. Así, existe $w\in U_x\cap A\subseteq U_a\cap A\subseteq V\cap A=\{a\}$, es decir, $a=w\in U_x$, luego $U_a\subseteq U_x$ y $U_x=U_a$. Por tanto, $U_a$ es minimal en $\mathcal{B}_{\min}$, es decir, $U_a\in\mathcal{B}$. Como $a\in U_a\cap A=\{a^*\}$, entonces $a=a^*=f(U_a)$, de modo que $f$ es sobreyectiva. Tenemos además que, para cualquier $U_x\in\mathcal{B}$, $f(U_x)=x^*\in U_x\cap A$, y por tanto $f(U_x)\in U_x$. También, como vimos, si $a\in A$ entonces $U_a\in\mathcal{B}$, luego $U_a$ es minimal en $\mathcal{B}_{\min}$. Por tanto,
   $$
   A\subseteq\left\lbrace x\in X\mid U_x\text{ es minimal en }\mathcal{B}_{\min}\right\rbrace.
   $$
   Ahora, notemos que si $U_x,U_y\in\mathcal{B}$ con $f(U_x)=f(U_y)$, como $f(U_x)\in U_x$ y $f(U_y)\in U_y$, se tiene $U_x\cap U_y\ne\emptyset$. Así, existe $z\in U_x$ con $z\in U_y$. Como $U_x$ y $U_y$ son abiertos, se tiene $U_z\subseteq U_x$ y $U_z\subseteq U_y$. Por la minimalidad de $U_x$ y $U_y$ en $\mathcal{B}_{\min}$ se sigue que $U_x=U_z=U_y$. Por tanto $f$ es inyectiva. Así, $f$ es biyectiva y $|\mathcal{B}|=|A|$, con $A$ un subconjunto denso minimal de $X$ arbitrario. Por tanto, si $A$ y $A'$ son subconjuntos densos minimales de $X$ entonces $|A|=|\mathcal{B}|=|A'|$.
\end{proof}
\vspace{1em}
{\color[HTML]{188539}
   En \cite{jafarian2013} se presenta el anterior Teorema junto con otro ítem. Lo transcribo literalmente, junto con el párrafo que lo antecede:\\\\
The following theorem which we could not find any proof for it elsewhere, is true in every Alexandroff topological space and specially in every graphic topological space.\\\\
\textbf{Theorem 8.4.} \textit{Let $(V,\tau)$ be an Alexandroff topological space, $\mathcal{B}_\tau=\left\lbrace U_x\mid x\in V\right\rbrace$ and $\mathcal{B}=\left\lbrace U_x\mid x\in V,\, U_x \text{ is minimal in }\mathcal{B}_\tau\right\rbrace$.}
\begin{itemize}
   \item[1.] \textit{If $A\subseteq V$ is a minimal dense subset in $(V,\tau)$, then there exists a surjective function $f:\mathcal{B}\longrightarrow A$ such that $f(U_x)\in U_x$ for every $U_x\in\mathcal{B}$. In particular, $A\subseteq\left\lbrace x\in V\mid U_x \text{ is minimal in }\mathcal{B}_\tau\right\rbrace$.}
   \item[2.] \textit{Conversely, if $f:\mathcal{B}\longrightarrow V$ is a function such that $f(U_x)\in U_x$ for every $U_x\in\mathcal{B}$, then $f(\mathcal{B})$ is a minimal dense subset in $(V,\tau)$.}
\end{itemize}
\textit{Specially, if $A$ and $A'$ are minimal dense subsets in $(V,\tau)$, then we have $|A|=|A'|$.}\\\\
Sin embargo creo que el ítem 2. es falso. A continuación desarrollo lo que considero es un contraejemplo (algo similar a la topología de colas (a derecha) abiertas y cerradas, pero sobre $\mathbb{N}$ para hacerlo más sencillo, y agregando $\left\lbrace 0\right\rbrace$ como abierto para garantizar que $\mathcal{B}$ no es vacío, lo cual de por sí ya daría un contraejemplo, pero bastante trivial, tomando $f$ como la función vacía).
\begin{Prop}
   Para cada $n\in\mathbb{N}$ definimos $n_{\ge}:=\left\lbrace m\in\mathbb{N}\mid m\ge n\right\rbrace$. Tomamos $\mathcal{C}=\left\lbrace n_{\ge}\mid n\in\mathbb{N}\right\rbrace\cup\left\lbrace\{0\}\right\rbrace$. Entonces $\mathcal{C}$ es base para una topología sobre $\mathbb{N}$.
\end{Prop}
\begin{proof}
   Usamos el ítem 3. de la \Cref{Prop:Caracterizacion_base}. Tenemos que $\mathcal{C}$ es una colección de subconjuntos de $\mathbb{N}$. Notemos que $0_{\ge}=\mathbb{N}$ y $0_{\ge}\in\mathcal{C}$, luego $\mathbb{N}=\bigcup\mathcal{C}$. Además:
   \begin{itemize}
      \item $\{0\}\cap 0_{\ge}=\{0\}\in\mathcal{C}$;
      \item $\{0\}\cap n_{\ge}=\emptyset$ para todo $n\in\mathbb{N}\setminus\{0\}$;
      \item $m_{\ge}\cap n_{\ge}=\left(\max\{m,n\}\right)_{\ge}\in\mathcal{C}$ para cualesquiera $m,n\in\mathbb{N}$,
   \end{itemize}
   con lo cual la intersección de cualesquiera dos elementos de $\mathcal{C}$ es unión (posiblemente vacía) de elementos de $\mathcal{C}$. Obtenemos que $\mathcal{C}$ es base para una topología sobre $\mathbb{N}$.
\end{proof}
\begin{Prop}
   La topología $\tau_{\mathcal{C}}$ generada por $\mathcal{C}$ es una topología de Alexandroff sobre $\mathbb{N}$.
\end{Prop}
\begin{proof}
   Veamos que la intersección arbitraria de abiertos básicos es abierta. Sea $\left\lbrace C_i\right\rbrace_{i\in I}$ una colección de elementos de $\mathcal{C}$; sin pérdida de generalidad podemos suponer que sus elementos son distintos y que $|I|>1$.
   \begin{itemize}
      \item Supongamos que $I$ es finito. Si $\left\lbrace C_i\right\rbrace_{i\in I}=\left\lbrace \{0\},0_{\ge}\right\rbrace$ entonces $\bigcap_{i\in I}C_i=\{0\}\in\tau_{\mathcal{C}}$. Si para algún $n>0$ se tiene $\{0\},n_{\ge}\in\left\lbrace C_i\right\rbrace_{i\in I}$ entonces $\bigcap_{i\in I}C_i\subseteq\{0\}\cap n_{\ge}=\emptyset\in\tau_{\mathcal{C}}$. Si $\{0\}\notin\left\lbrace C_i\right\rbrace_{i\in I}$, tomando $n=\max\left\lbrace m\mid m_{\ge}\in\left\lbrace C_i\right\rbrace_{i\in I}\right\rbrace$, se tiene $\bigcap_{i\in I}C_i=n_{\ge}\in\tau_{\mathcal{C}}$.
      \item Supongamos que $I$ es infinito. Veamos que $\bigcap_{i\in I}C_i=\emptyset$. Razonando por contradicción supongamos que existe $z\in\bigcap_{i\in I}C_i$. Si no existiera $w>z$ con $w_{\ge}\in\left\lbrace C_i\right\rbrace_{i\in I}$ entonces $\left\lbrace C_i\right\rbrace_{i\in I}\subseteq\left\lbrace \{0\},0_{\ge},1_{\ge},\dots,z_{\ge}\right\rbrace$, contradiciendo que $I$ es infinito. Por tanto existe $w>z$ con $w_{\ge}\in\left\lbrace C_i\right\rbrace_{i\in I}$. Pero entonces $z\in\bigcap_{i\in I}C_i\subseteq w_{\ge}$, luego $z\in w_{\ge}$, es decir $z\ge w$, contradiciendo $w>z$. Por tanto $\bigcap_{i\in I}C_i=\emptyset\in\tau_{\mathcal{C}}$.
   \end{itemize}
   Obtenemos así que $\tau_{\mathcal{C}}$ es una topología de Alexandroff.
\end{proof}
\begin{Prop}
   En $(\mathbb{N},\tau_{\mathcal{C}})$ se tiene $U_0=\{0\}$ y $U_n=n_{\ge}$ para cualquier $n\in\mathbb{N}$ con $n>0$.
\end{Prop}
\begin{proof}
   Como $\{0\}\subseteqab\mathbb{N}$, se tiene $U_0=\{0\}$. Sea $n\in\mathbb{N}$ con $n>0$. Si $V$ es una vecindad de $n$, se tiene que $V=m_{\ge}$ para algún $m\in\mathbb{N}$ con $m>0$. Como $n\in m_{\ge}$ entonces $n\ge m$, luego para todo $l\ge n$ se tiene $l\ge m$, es decir $n_{\ge}=\left\lbrace l\mid l\ge n\right\rbrace\subseteq\left\lbrace l\mid l\ge m\right\rbrace=m_{\ge}=V$. Por tanto $n_{\ge}$ es la menor vecindad de $n$, es decir, $U_n=n_{\ge}$.
\end{proof}
\begin{Prop}
   En $(\mathbb{N},\tau_{\mathcal{C}})$, tomando la notación del Teorema 8.4 de \cite{jafarian2013}, se tiene $\mathcal{B}=\left\lbrace U_0\right\rbrace$.
\end{Prop}
\begin{proof}
   Claramente $U_0=\{0\}$ es minimal en $\mathcal{B}_\tau$, pues los elementos de $\mathcal{B}_\tau$ son no vacíos y por tanto no están propiamente contenidos en $\{0\}$. Así, $U_0\in\mathcal{B}$. Además, para cualquier $n\in\mathbb{N}$ con $n>0$ se tiene $U_{n+1}=(n+1)_{\ge}\subsetneq n_{\ge}=U_n$, luego $U_n$ no es minimal en $\mathcal{B}_\tau$, es decir, $U_n\notin\mathcal{B}$. Como $\mathcal{B}\subseteq\mathcal{B}_\tau=\left\lbrace U_n\mid n\in\mathbb{N}\right\rbrace$, se tiene $\mathcal{B}=\left\lbrace U_0\right\rbrace$.
\end{proof}
En $(\mathbb{N},\tau_{\mathcal{C}})$ tomemos entonces $f:\mathcal{B}\to\mathbb{N}:U_0\mapsto 0$. Como $f(U_0)=0\in\{0\}=U_0$, se cumple que $f(U_x)\in U_x$ para todo $U_x\in\mathcal{B}$. Tenemos $f(\mathcal{B})=\{0\}$ pero, por ejemplo, $\{0\}\cap 1_{\ge}=\{0\}\cap\left\lbrace l\mid l\ge 1\right\rbrace=\emptyset$, con $1_{\ge}$ abierto no vacío de $\mathbb{N}$, así que por el ítem 3. de la \Cref{Prop:Caracterizacion_denso} se tiene que $f(\mathcal{B})$ no es un subconjunto denso en $(\mathbb{N},\tau_{\mathcal{C}})$.\\\\
En la demostración que presenta \cite{jafarian2013} para el Teorema 8.4 se lee, sin ninguna prueba, ``For every $x\in V$, there exists an element $a\in V$ such that $U_a\subseteq U_x$ and $U_a\in\mathcal{B}$''. Sin embargo, en general esto es falso. En $(\mathbb{N},\tau_{\mathcal{C}})$, por ejemplo, la afirmación solo se cumple para $0$, pues como vimos, $\mathcal{B}=\left\lbrace U_0\right\rbrace$ con $U_0=\{0\}$, así que para $n\in\mathbb{N}$ con $n>0$, $U_n=n_{\ge}$ no contiene ningún elemento de $\mathcal{B}$.
}

\end{comment}

\chapter{Topología de adyacencia}\label{chap:topologia_de_adyacencia}
A lo largo de todo el capítulo supondremos que los grafos en cuestión son localmente finitos y sin vértices aislados.
\begin{Prop}
   Sea $G=(V,E)$ un grafo. El conjunto $\mathcal{S}_{\mathcal{A},G}:=\left\lbrace A_x\mid x\in V\right\rbrace$ es subbase para una topología sobre $V$.
\end{Prop}
\begin{proof}
   Por definición, $\mathcal{S}_{\mathcal{A},G}$ es una colección de subconjuntos de $V$. Si $V = \emptyset$, entonces $\mathcal{S}_{\mathcal{A},G} = \emptyset$, la cual es subbase para la topología $\{\emptyset\}$ sobre $V = \emptyset$. Supongamos $V \neq \emptyset$. Sea $v \in V$ un vértice arbitrario. Dado que $G$ no contiene vértices aislados, existe $u \in V$ tal que $v \in A_u$. En consecuencia, $v \in \bigcup_{x \in V} A_x$, de lo cual se deduce que $\bigcup_{x \in V} A_x = V$. Por lo tanto, $\mathcal{S}_{\mathcal{A},G}$ es una subbase para una topología sobre $V$.
\end{proof}
Como pudimos ver, la condición de no tener vértices aislados es clave para que $\mathcal{S}_{\mathcal{A},G}$ sea subbase para una topología sobre $V$, pues si $G$ tuviera algún vértice aislado $v\in V$, se tendría $v\notin\bigcup_{x\in V}A_x$, dado que $v\notin A_x$ para todo $x\in V$.
\begin{Def}
   Sea $G=(V,E)$ un grafo. La topología generada por la subbase $\mathcal{S}_{\mathcal{A},G}$, denotada por $\tau_{\mathcal{A},G}:=\tau_{\mathcal{S}_{\mathcal{A},G}}$, es llamada la \textbf{topología de adyacencia} de $G$. El espacio topológico $(V,\tau_{\mathcal{A},G})$ es llamado el \textbf{espacio de adyacencia} de $G$. 
\end{Def}
\begin{Prop}\label{Prop:Calculo_topologias_adyacencia}
   \phantom{o}
   \begin{itemize}
      \item[1.] Para todo $n\ge 2$, la topología de adyacencia de $K^n$ es discreta.
      \item[2.] Para todo $n\ge 5$, la topología de adyacencia de $C^n$ es discreta.
      \item[3.] Para todo $n\ge 3$, la topología de adyacencia de $P^n$ es no discreta.
      \item[4.] Para cualesquiera $n,m\ge 1$, la topología de adyacencia del grafo bipartito completo $K_{n,m}$ es $\{\emptyset, V, A, B\}$, donde $V=V\left( K_{n,m}\right)$ y $A$ y $B$ son las particiones de $V$.
   \end{itemize}
\end{Prop}
\begin{proof}
   \phantom{o}
   \begin{itemize}
      \item[1.] Sea $n\ge 2$. Denotemos $V\left( K^n\right)=\{v_1,\dots, v_n\}$. Sea $i\in\{1,\dots, n\}$. Basta ver que $\left\lbrace v_i\right\rbrace\in \tau_{\mathcal{A},K^n}$ Veamos que $\left\lbrace v_i\right\rbrace=\bigcap_{1\le j\le n, j\ne i}A_{v_j}$. Por la definición de $K^n$ se tiene $v_i\in A_{v_j}$ para todo $j$ tal que $1\le j\le n$ y $j\ne i$. Por tanto $v_i\in\bigcap_{1\le j\le n, j\ne i} A_{v_j}$ y $\left\lbrace v_i\right\rbrace\subseteq \bigcap_{1\le j\le n, j\ne i} A_{v_j}$. Ahora, dado $k$ con $1\le k\le n$ y $k\ne i$ se tiene $v_k\notin A_{v_k}$, luego $v_k\notin \bigcap_{1\le j\le n, j\ne i} A_{v_j}$. Con esto, $V\setminus \left\lbrace v_i\right\rbrace\subseteq V\setminus \bigcap_{1\le j\le n, j\ne i} A_{v_j}$, es decir $\bigcap_{1\le j\le n, j\ne i} A_{v_j} \subseteq \left\lbrace v_i\right\rbrace$. Obtenemos así $\left\lbrace v_i\right\rbrace=\bigcap_{1\le j\le n, j\ne i} A_{v_j}$. Como $\left\lbrace v_i\right\rbrace$ es una intersección finita de elementos de $\mathcal{S}_{\mathcal{A},K^n}$, se tiene $\left\lbrace v_i\right\rbrace\in \tau_{\mathcal{A},K^n}$.
      \item[2.] Sea $n\ge 5$. Denotemos $V\left( C^n\right)=\left\lbrace v_0,\dots,v_{n-1}\right\rbrace$. Basta ver que $\left\lbrace v_i\right\rbrace\in\tau_{\mathcal{A},C^n}$ para $i\in\{0,\dots,n-1\}$.
         \begin{figure}[H]
\centering
\resizebox{0.25\textwidth}{!}{
\begin{tikzpicture}[
  scale=0.7,
  vnode/.style={circle, fill=blue!20, draw=black, thick, inner sep=0pt, minimum size=0.9cm}
]

  % Vértices iniciales (lado derecho)
  \node[vnode] (v1) at (90:3) {$v_0$};
  \node[vnode] (v2) at (30:3) {$v_1$};
  \node[vnode] (v3) at (-30:3) {$v_2$};

  % Vértices finales (lado izquierdo)
  \node[vnode] (vn1) at (210:3) {$v_{n-2}$};
  \node[vnode] (vn) at (150:3) {$v_{n-1}$};

  % Puntos suspensivos y nodos auxiliares para continuar las aristas en la parte inferior
  \node[draw=none, fill=none] (dots) at (-90:3) {$\dots$};
  \node[draw=none, fill=none, minimum size=0pt] (v3_inv) at (-60:3) {};
  \node[draw=none, fill=none, minimum size=0pt] (vn1_inv) at (-120:3) {};

  % Aristas continuas visibles
  \draw[thick] (v1) -- (v2);
  \draw[thick] (v2) -- (v3);
  \draw[thick] (vn1) -- (vn);
  \draw[thick] (vn) -- (v1); % Arista que cierra el ciclo entre v_n y v_1

  % Aristas que se dirigen hacia la elipsis inferior
  \draw[thick] (v3) -- (v3_inv);
  \draw[thick] (vn1_inv) -- (vn1);

\end{tikzpicture}}
\caption{El grafo $C^n$}
\label{fig:ciclo_n_adyacencia}
\end{figure}

         Para todo $i\in\{0,\dots,n-1\}$ se tiene $A_{v_i}=\left\lbrace v_{(i-1)\bmod n},v_{(i+1)\bmod n}\right\rbrace$ y por tanto $A_{v_{(i-1)\bmod n}}=\left\lbrace v_{(i-2)\bmod n},v_i\right\rbrace$ y $A_{v_{(i+1)\bmod n}}=\left\lbrace v_i,v_{(i+2)\bmod n}\right\rbrace$. Como $n\ge 5$ se tiene $(i-2)\bmod n\ne (i+2)\bmod n$, de modo que $v_{(i-2)\bmod n}\ne v_{(i+2)\bmod n}$ y por tanto $\left\lbrace v_i\right\rbrace=A_{v_{(i-1)\bmod n}}\cap A_{v_{(i+1)\bmod n}}$. Como $\left\lbrace v_i\right\rbrace$ es una intersección finita de elementos de $\mathcal{S}_{\mathcal{A},C^n}$ entonces $\left\lbrace v_i\right\rbrace\in\tau_{\mathcal{A},C^n}$.
      \item[3.] Sea $n\ge 3$.
         \begin{figure}[H]
\centering
\resizebox{0.45\textwidth}{!}{
\begin{tikzpicture}[
  scale=1.2,
  vnode/.style={circle, fill=blue!20, draw=black, thick, inner sep=0pt, minimum size=0.9cm}
]

  % Posicionamiento de vértices horizontales
  \node[vnode] (v1) at (0, 0) {$v_1$};
  \node[vnode] (v2) at (2, 0) {$v_2$};
  
  % Puntos suspensivos y nodos auxiliares para conectar las aristas
  \node[draw=none, fill=none] (dots) at (4, 0) {$\dots$};
  \node[draw=none, fill=none, inner sep=0pt, minimum size=0pt] (v2_out) at (3.1, 0) {};
  \node[draw=none, fill=none, inner sep=0pt, minimum size=0pt] (vn_in) at (4.9, 0) {};

  % Vértice final
  \node[vnode] (vn) at (6, 0) {$v_n$};

  % Arista continua entre v1 y v2
  \draw[thick] (v1) -- (v2);

  % Aristas de entrada/salida hacia la elipsis
  \draw[thick] (v2) -- (v2_out);
  \draw[thick] (vn_in) -- (vn);

\end{tikzpicture}}
\caption{El grafo $P^n$}
\label{fig:camino_n}
\end{figure}

         Veamos que $\left\lbrace v_1\right\rbrace\notin\tau_{\mathcal{A},P^n}$. Sea $U\in\tau_{\mathcal{A},P^n}$ con $v_1\in U$. Basta probar que $v_3\in U$. Tenemos $A_{v_1}=\left\lbrace v_2\right\rbrace$, $A_{v_i}=\left\lbrace v_{i-1},v_{i+1}\right\rbrace$ para cualquier $i$ con $2\le i\le n-1$ y $A_{v_n}=\left\lbrace v_{n-1}\right\rbrace$. Existe $B\subseteq\mathcal{B}_{\mathcal{S}_{\mathcal{A},P^n}}$ tal que $U=\bigcup B$. Entonces $v_1\in S$ para algún $S\in B$. Como en particular $S\in \mathcal{B}_{\mathcal{S}_{\mathcal{A},P^n}}$, luego $S=\bigcap_{j=1}^{m}A_{v_{i_j}}$ para algún $m\in \mathbb{Z}^+$ y con $i_j\in\{1,\dots,n\}$ para todo $j\in\{1,\dots,m\}$. Tenemos que $v_1\in A_{v_{i_j}}$ para todo $j\in\{1,\dots,m\}$. Como $v_1\in A_{v_2}$ y $v_1\notin A_{v_k}$ para todo $k\in\{3,\dots,n\}$ entonces $A_{v_{i_j}}=A_{v_2}$ para todo $j\in\{1,\dots,m\}$, de modo que $A_{v_2}=S\subseteq U$ y por tanto $v_3\in U$.
      \item[4.] Sean $n,m\ge 1$. Dado $a\in A$ se tiene $A_a=B$, y dado $b\in B$ se tiene $A_b=A$. Por tanto $\mathcal{S}_{\mathcal{A},K_{n,m}}=\left\lbrace A,B\right\rbrace$. Como $A\cap B=\emptyset$, tenemos $\mathcal{B}_{\mathcal{S}_{\mathcal{A},K_{n,m}}}=\left\lbrace A,B\right\rbrace$, y como $A\cup B=V$ entonces $\tau_{\mathcal{A},k_{n,m}}=\left\lbrace \emptyset,V,A,B\right\rbrace$. 
   \end{itemize}
\end{proof}
\begin{Obs}
   Se tiene que $K^1=P^1$ está formado por un solo vértice (que por tanto es aislado), de modo que no tiene definida topología de adyacencia. Además $C^3=K^3$, y por tanto $C^3$ tiene topología de adyacencia discreta. En la \Cref{sec:computacional_adyacencia} calcularemos la topología de adyacencia de $C^4$ y veremos que no es discreta. 
\end{Obs}
Hasta el momento no hemos necesitado la propiedad de que los grafos sean localmente finitos; esta la exigimos para que la topología de adyacencia sea de Alexandroff, como muestra la siguiente proposición.
\begin{Prop}\label{Prop:Adyacencia_Alexandroff}
   Sea $G=(V,E)$ un grafo. Se tiene que $(V,\tau_{\mathcal{A},G})$ es un espacio de Alexandroff.
\end{Prop}
\begin{proof}
   Sea $U$ una intersección arbitraria de elementos de $\mathcal{S}_{\mathcal{A},G}$. Por la \Cref{Prop:Alexandroff_con_subbase} basta ver que $U\subseteqab X$ para concluir que $X$ es de Alexandroff. Si $U=\emptyset$ entonces $U\subseteqab X$. Supongamos que $U\ne \emptyset$. Sea $y\in U$. Por la definición de $\mathcal{S}_{\mathcal{A},G}$, existe $V'\subseteq V$ tal que $U=\bigcap_{x\in V'}A_x$. Para todo $x\in V'$ tenemos $y\in A_x$ y por tanto $x\in A_y$, de modo que $V'\subseteq A_y$. Como $G$ es localmente finito, $A_y$ es finito, y por tanto $V'$ es finito. Con esto, $U$ es una intersección finita de conjuntos abiertos, con lo que obtenemos que $U\subseteqab X$.
\end{proof}
%\textcolor{red}{Mostrar un grafo que no sea localmente finito y no tenga topología de adyacencia de Alexandroff / ¿todo grafo que no es localmente finito tiene topología de adyacencia que no es de Alexandroff?}\\\\
Como la topología de adyacencia de un grafo es de Alexandroff, podemos hablar de vecindades mínimas y bases mínimas en la forma que precisamos en \Cref{sec:Topologias_Alexandroff}.
\begin{Prop}\label{Prop:U_x_2}
   Sea $G=(V,E)$ un grafo. Para todo $x\in V$ tenemos $U_x=\bigcap_{y\in A_x}A_y$ y por tanto $U_x$ es finito.
\end{Prop}
\begin{proof}
   Sea $x\in V$. Para todo $y\in A_x$ tenemos $x\in A_y$, luego $x\in\bigcap_{y\in A_x}A_y$. Además $\bigcap_{y\in A_x}A_y\subseteqab X$. Como $U_x$ es la vecindad mínima de $x$, se sigue $U_x\subseteq \bigcap_{y\in A_x}A_y$. Ahora, como $x\in U_x$ y $U_x\subseteqab X$, existe $B\in\mathcal{B}_{\mathcal{S}_{\mathcal{A},G}}$ tal que $x\in B \subseteq U_x$. Por definición de $\mathcal{B}_{\mathcal{S}_{\mathcal{A},G}}$, existe algún subconjunto finito $V'$ de $V$ tal que $B=\bigcap_{y\in V'}A_y$. Para todo $y\in V'$ se tiene $x\in A_y$ y por tanto $y\in A_x$, de donde $V'\subseteq A_x$. Con esto, $\bigcap_{y\in A_x}A_y\subseteq\bigcap_{y\in V'}A_y \subseteq U_x$. Obtenemos así $U_x=\bigcap_{y\in A_x}A_y$. Como $G$ es localmente finito, $A_y$ es finito para todo $y\in A_x$ y por tanto $U_x$ es finito.
\end{proof}
\begin{Obs}\label{Obs:Igualdad_adyacencias_vecindades}
   De la \Cref{Prop:U_x_2} se tiene que si $x,y\in V$ son tales que $A_x=A_y$ entonces $U_x=\bigcap_{z\in A_x}A_z=\bigcap_{z\in A_y}A_z=U_y$.
\end{Obs}
\begin{Cor}\label{Cor:U_x_contenencia}
    Sea $G=(V,E)$ un grafo. Para cualesquiera $x,z\in V$ se tiene $z\in U_x$ si, y solo si, $A_x\subseteq A_z$.
\end{Cor}
\begin{proof}
   Sean $x,z\in V$.\\
   ($\Rightarrow$) Supongamos $z\in U_x$. Como $U_x=\bigcap_{y\in A_x}A_y$ entonces $z\in \bigcap_{y\in A_x}A_y$. Para cualquier $y\in A_x$ se tiene $z\in A_y$ y por tanto $y\in A_z$. Luego $A_x\subseteq A_z$.\\
   ($\Leftarrow$) Supongamos $A_x\subseteq A_z$. Para cualquier $y\in A_x$ se tiene $y\in A_z$ y por tanto $z\in A_y$. Así, $z\in \bigcap_{y\in A_x}A_y=U_x$.
\end{proof}
\begin{Obs}\label{Obs:U_x}
   Por lo anterior, para cualquier vértice $x$ de un grafo $G=(V,E)$ se tiene $U_x=\left\lbrace z\in V\mid A_x\subseteq A_z\right\rbrace$.
\end{Obs}
El siguiente corolario nos da una condición necesaria y suficiente para que un grafo tenga topología de adyacencia discreta.
\begin{Cor}\label{Cor:Discreto_Alexandroff}
   Sea $G=(V,E)$ un grafo. El espacio $(V,\tau_{\mathcal{A},G})$ es discreto si, y solo si $A_x \nsubseteq A_y$ para cualesquiera $x,y\in V$ con $x\ne y$.
\end{Cor}
\begin{proof}
   ($\Rightarrow$) Supongamos que $(V,\tau_{\mathcal{A},G})$ es un espacio discreto. Sean $x,y\in V$ con $x\ne y$. Como $(V, \tau_{\mathcal{A},G})$ es un espacio de Alexandroff discreto, se tiene que $U_x=\{x\}$. Si $A_x\subseteq A_y$ entonces por la \Cref{Obs:U_x} se tiene $y\in U_x$, contradiciendo $U_x=\{x\}$. Por tanto $A_x\nsubseteq A_y$.\\
   ($\Leftarrow$) Supongamos $A_x\nsubseteq A_y$ para cualesquiera $x,y\in V$ con $x\ne y$. Sea $x\in V$. Sea $y\in V$ con $x\ne y$. Si $y\in U_x$ entonces por la \Cref{Obs:U_x} se tiene $A_x\subseteq A_y$, de modo que $y\notin U_x$. Con esto, $\{x\}=U_x\subseteqab V$. Se sigue que $(V,\tau_{\mathcal{A},G})$ es un espacio discreto.  
\end{proof}
\begin{Cor}
   Sea $G=(V,E)$ un grafo. Entonces $(V,\tau_{\mathcal{A},G})$ es $T_1$ si, y solo si, $A_x\nsubseteq A_y$ para cualesquiera $x,y\in V$ con $x\ne y$.
\end{Cor}
\begin{proof}
   Se sigue directamente de la \Cref{Prop:Separacion_Alexandroff} y el \Cref{Cor:Discreto_Alexandroff}.
\end{proof}
\begin{Cor}\label{Cor:T_0_Adyacencia}
   Sea $G=(V,E)$ un grafo. Entonces $(V,\tau_{\mathcal{A},G})$ es $T_0$ si, y solo si $A_x\ne A_y$ para cualesquiera $x,y\in V$ con $x\ne y$.
\end{Cor}
\begin{proof}
   ($\Rightarrow$) Supongamos que $(V,\tau_{\mathcal{A},G})$ es $T_0$. Sean $x,y\in V$ con $x\ne y$. Sin pérdida de generalidad podemos afirmar que existe $U\subseteqab V$ con $x\in U$ y $y\notin U$. Como $U$ es una vecindad de $x$, se tiene $U_x\subseteq U$. Si $A_x=A_y$ entonces, por la \Cref{Obs:Igualdad_adyacencias_vecindades}, se tiene $U_x=U_y$ y por tanto $U_y\subseteq U$, con lo cual en particular $y\in U$. Por lo tanto $A_x\ne A_y$.\\
   ($\Leftarrow$) Supongamos que $A_x\ne A_y$ para $x,y\in V$ con $x\ne y$. Sean $x,y\in V$ con $x\ne y$. Como $A_x\ne A_y$ tenemos $A_x\nsubseteq A_y$ o $A_y\nsubseteq A_x$. Sin pérdidad de generalidad supongamos $A_x\nsubseteq A_y$. Entonces existe $z\in A_x$ con $z\notin A_y$, de donde $x\in A_z$ y $y\notin A_z$. Como $A_z\subseteqab V$, se sigue que $(V,\tau_{\mathcal{A},G})$ es $T_0$.
\end{proof}
Si $G$ es un árbol, basta verificar la anterior condición para hojas distintas, como lo muestra el siguiente corolario.
\begin{Cor}
   Sea $T=(V,E)$ un árbol. Entonces $(V,\tau_{\mathcal{A},T})$ es $T_0$ si, y solo si, $A_x\ne A_y$ para cualesquiera $x,y\in V$ con $x\ne y$ y $\operatorname{deg}(x)=\operatorname{deg}(y)=1$.
\end{Cor}
\begin{proof}
   ($\Rightarrow$) Es inmediato por el \Cref{Cor:T_0_Adyacencia}.\\
   ($\Leftarrow$) Supongamos que $A_x\ne A_y$ para $x,y\in V$ con $x\ne y$ y $\operatorname{deg}(x)=\operatorname{deg}(y)=1$. Veamos que esto implica $A_x\ne A_y$ para $x,y\in V$ con $x\ne y$. Razonando por contradicción supongamos que existen $x,y\in V$ con $x\ne y$ tales que $A_x=A_y$. Por hipótesis, $x$ y $y$ no son hojas. Por tanto $|A_x|=|A_y|\ge 2$. Existen $u,v\in A_x = A_y$ con $u\ne v$; además $u\ne x$, $u\ne y$, $v\ne x$, $v\ne y$ (pues, por ejemplo, $u=x$ implica $x\in A_x$, lo cual es una contradicción de nuestra definición de grafo simple; similarmente con las demás afirmaciones). Así, $xuyvx$ es un ciclo contenido en $T$, lo cual contradice que $T$ sea un árbol. Se sigue que $A_x\ne A_y$ para cualesquiera $x,y\in V$ con $x\ne y$. Por el \Cref{Cor:T_0_Adyacencia} se obtiene que $(V,\tau_{\mathcal{A},T})$ es $T_0$.
\end{proof}
\begin{Cor}\label{Cor:U_x_3}
    Sea $G=(V,E)$ un grafo. Entonces:
    \begin{itemize}
       \item[1.] Para todo $x\in V$ se tiene $U_x\subseteq \{x\}\cup \left\lbrace y\in V\mid d(x,y)=2\right\rbrace$.
       \item[2.] Para todo $x\in V$ se tiene $U_x\cap A_x=\emptyset$.
       \item[3.] Para cualesquiera $x,y\in V$, si $x$ y $y$ son adyacentes entonces $U_x\cap U_y=\emptyset$.
      \item[4.] Para cualquier $x\in V$ se tiene $\overline{\{x\}}\subseteq \overline{U_x}\subseteq A_x^c$ y $\overline{A_x}\subseteq U_x^c$.
    \end{itemize}
\end{Cor}
\begin{proof}
   \phantom{o}
   \begin{itemize}
      \item[1.] Sea $x\in V$. Sea $z\in U_x$. Si $z=x$ entonces $z\in \{x\}\cup \left\lbrace y\in V\mid d(x,y)=2\right\rbrace$. Supongamos $z\ne x$. Sea $w\in A_x$. Por la \Cref{Prop:U_x_2} tenemos $U_x=\bigcap_{y\in A_x}A_y$, de modo que $z\in A_w$. Así, $zwx$ es un camino en $G$ que conecta $z$ y $x$. Si $z\in A_x$, como $z\in U_x$ se tendría $z\in A_z$, contradiciendo la definición de grafo simple. Por tanto $z\notin A_x$, de lo que se sigue que $d(x,z)=2$. Por tanto $z\in \{x\}\cup \left\lbrace y\in V\mid d(x,y)=2\right\rbrace$ y obtenemos $U_x\subseteq \{x\}\cup \left\lbrace y\in V\mid d(x,y)=2\right\rbrace$.
      \item[2.] Sea $x\in V$. Si $z\in A_x$ entonces $d(x,z)=1$, y por el ítem $1.$ se sigue que $z\notin U_x$. Por tanto $U_x\cap A_x=\emptyset$.
      \item[3.] Sean $x,y\in V$ y supongamos que $x$ y $y$ son adyacentes. Tenemos $y\in A_x$ y $A_x\subseteqab V$, de modo que $A_x$ es una vecindad de $y$, con lo cual $U_y\subseteq A_x$. Si existiera $z\in U_x\cap U_y$ entonces se tendría $z\in U_x\cap A_x$, contradiciendo el ítem $2.$ Por tanto $U_x\cap U_y=\emptyset$.
      \item[4.] Sea $x\in V$. Sabemos que $\{x\}\subseteq U_x$, y por el ítem $2.$ se tiene $U_x\subseteq A_x^c$ y $A_x\subseteq U_x^c$. Por la \Cref{Prop:Contenencia_Clausura} se sigue que $\overline{\{x\}}\subseteq \overline{U_x}\subseteq \overline{A_x^c}$ y $\overline{A_x}\subseteq \overline{U_x^c}$. Como $A_x$ y $U_x$ son abiertos, entonces $A_x^c$ y $U_x^c$ son cerrados, de modo que $\overline{A_x^c}=A_x^c$ y $\overline{U_x^c}=U_x^c$ por la \Cref{Prop:Carac_abierto_cerrado}, con lo cual el resultado se obtiene.
   \end{itemize}
\end{proof}
\textcolor{red}{Qué es la distancia entre dos vértices??????\\Ejemplo de que en 1. la otra contenencia no se tiene.}
\begin{Cor}\label{Cor:Clausura_unitario}
   Sea $G=(V,E)$ un grafo. Para cualesquiera $x,y\in V$ se tiene $y\in\overline{\{x\}}$ si, y solo si, $A_y\subseteq A_x$. 
\end{Cor}
\begin{proof}
   Sean $x,y\in V$.\\
   ($\Rightarrow$) Supongamos que $y\in\overline{\{x\}}$. Por la \Cref{Prop:Caracterizacion_punto_adherencia} se tiene que cualquier vecindad de $y$ contiene a $x$; en particular $x\in U_y$. Por el \Cref{Cor:U_x_contenencia} se sigue que $A_y\subseteq A_x$.\\
   ($\Leftarrow$) Supongamos que $A_y\subseteq A_x$. Por el \Cref{Cor:U_x_contenencia} se tiene $x\in U_y$. Sea $V$ una vecindad de $y$. Se tiene $U_y\subseteq V$ y por tanto $x\in V$, es decir, $V$ interseca a $\{x\}$. Por la \Cref{Prop:Caracterizacion_punto_adherencia} se sigue que $y\in\overline{\{x\}}$.
\end{proof}
\begin{Lema}\label{Lema:lema_homeomorfismo}
   Sean $G_1=(V_1,E_1)$ y $G_2=(V_2,E_2)$ os grafos. Si $\varphi:V_1\to V_2$ es un isomorfismo de $G_1$ en $G_2$ entonces:\begin{itemize}
      \item[1.] $\varphi (A_x) = A_{\varphi (x)}$ para todo $x\in V_1$.
      \item[2.] $\varphi (U_x) = U_{\varphi (x)}$ para todo $x\in V_1$.
   \end{itemize}
\end{Lema}
\begin{proof}
   \phantom{o}
   \begin{itemize}
      \item[1.] Sea $x\in V_1$.
         \begin{itemize}
            \item[$\subseteq$:] Sea $w\in\varphi(A_x)$. Existe $x_w\in A_x$ tal que $w=\varphi(x_w)$. Entonces $x_wx\in E_1$, de donde $\varphi(x_w)\varphi(x)\in E_2$, pues $\varphi$ es un isomorfismo, con lo cual $w=\varphi(w_x)\in A_{\varphi(x)}$ y por tanto $\varphi(A_x)\subseteq A_{\varphi(x)}$.
            \item[$\supseteq$:] Sea $w\in A_{\varphi(x)}$. Como $\varphi$ es en particular sobreyectiva, existe $x_w\in V_1$ tal que $\varphi(x_w)=w$. Tenemps $\varphi(x_w)\varphi(x)=w\varphi(x)\in E_2$, de donde $x_wx\in E_1$, pues $\varphi$ es un isomorfismo, con lo cual $x_w\in A_x$ y por tanto $w=\varphi(x_w)\in\varphi(A_x)$. Con esto, $A_{\varphi(x)}\subseteq\varphi(A_x)$.
         \end{itemize}
      \item[2.] Sea $x\in V_1$. Por la \Cref{Prop:U_x_2} tenemos $U_x=\bigcap_{y\in A_x}A_y$ y $U_{\varphi(x)}=\bigcap_{w\in A_{\varphi(x)}}A_w$. Además tenemos 
         $$
         \begin{aligned}
            \varphi(U_x)&=\varphi\left( \bigcap_{y\in A_x}A_y\right)\\
                        &=\bigcap_{y\in A_x}\varphi(A_y)\\
                        &\phantom{=}\text{(}\varphi\text{ es biyectiva)}\\
                        &=\bigcap_{y\in A_x}A_{\varphi(y)}.\\
                        &\phantom{=}\text{(ítem 1.)}
         \end{aligned}
         $$
         Para cualquier $z\in \bigcap_{y\in A_x}A_{\varphi(y)}$ existe $y\in A_x$ tal que $z\in A_{\varphi(y)}$. Como $\varphi$ es un isomorfismo se tiene $\varphi(y)\in A_{\varphi(x)}$, y por tanto $z\in\bigcap_{w\in A_{\varphi(x)}}A_w$. Con esto, $\bigcap_{y\in A_x}A_{\varphi(y)}\subseteq\bigcap_{w\in A_{\varphi(x)}}A_w$. Ahora, para cualquier $z\in\bigcap_{w\in A_{\varphi(x)}}A_w$ existe $w\in A_{\varphi(x)}$ tal que $z\in A_w$. Como $\varphi$ es en particular sobreyectiva, exite $y\in V_1$ tal que $w=\varphi(y)$, luego $\varphi(y)\in A_{\varphi(x)}$. Como $\varphi$ es un isomorfismo, lo anterior implica $y\in A_x$, y $z\in A_w=A_{\varphi(y)}$, con lo cual $z\in \bigcap_{y\in A_x}A_{\varphi(y)}$ y tenemos $\bigcap_{w\in A_{\varphi(x)}}A_w \subseteq \bigcap_{y\in A_x}A_{\varphi(y)}$. Así,
         $$
         \varphi(U_x)=\bigcap_{y\in A_x}A_{\varphi(y)}=\bigcap_{w\in A_{\varphi(x)}}A_w=U_{\varphi(x)}.
         $$
   \end{itemize}
\end{proof}
\begin{Prop}\label{Prop:Isomorfismo_continua}
   Sean $G_1=(V_1,E_1)$ y $G_2=(V_2,E_2)$ dos grafos y $\varphi:V_1\to V_2$ un isomorfismo entre $G_1$ y $G_2$. Entonce $\varphi$ es una función continua entre los espacios topológicos $(V_1,\tau_{\mathcal{A},G_1})$ y $(V_2, \tau_{\mathcal{A},G_2})$.
\end{Prop}
\begin{proof}
   Usamos la caracterización dada en el ítem $4.$ de la \Cref{Prop:Caracterizacion_continuidad}. Sean $x\in V_1$ y $A$ una vecindad de $\varphi(x)$. Tenemos que $U_x$ es una vecindad de $x$, y $\varphi(U_x)=U_{\varphi(x)}$ por el \Cref{Lema:lema_homeomorfismo}, es decir, $\varphi(U_x)$ es la vecindad mínima de $\varphi(x)$, de modo que $\varphi(U_x)\subseteq A$. Por tanto $\varphi$ es continua.
\end{proof}
Como consecuencia inmediata de lo anterior tenemos que si dos grafos son isomorfos, entonces sus espacios de adyacencia son homeomorfos. Más precisamente:
\begin{Cor}
   Sean $G_1=(V_1,E_1)$ y $G_2=(V_2,E_2)$ dos grafos y $\varphi:V_1\to V_2$ un isomorfismo entre $G_1$ y $G_2$. Entonces $\varphi$ es también un homeomorfismo entre los espacios topológicos $(V_1,\tau_{\mathcal{A},G_1})$ y $(V_2,\tau_{\mathcal{A},G_2})$. 
\end{Cor}
\begin{proof}
   Por el \Cref{Cor:Inversa_isomorfismo} se tiene que $\varphi^{-1}$ es un isomorfismo de $G_2$ en $G_1$. El resultado se sigue de la \Cref{Prop:Isomorfismo_continua} aplicada a $\varphi$ y $\varphi^{-1}$.
\end{proof}
\begin{Obs}
   No se tiene en cambio la implicación de que grafos con espacios de adyacencia homeomorfos sean isomorfos. Por ejemplo, por la \Cref{Prop:Calculo_topologias_adyacencia}, las topologías de adyacencia de $K^5$ y $C^5$ son ambas discretas; como además $|V\left( K^5\right)|=5=|V\left( C^5\right)|$ entonces, por el \Cref{Ejm:homeomorfismos}, se tiene que $\left( V\left( K^5\right),\tau_{\mathcal{A},K^5}\right)$ y $\left( V\left( C^5\right),\tau_{\mathcal{A},C^5}\right)$ son homeomorfos. Sin embargo, por el \Cref{Cor:Isomorfismo_grado_iguales}, como $K^5$ tiene un vértice de grado $5$, mientras que $C^5$ no lo tiene, sabemos que $K^5$ y $C^5$ no son isomorfos. También podemos construir dos grafos infinitos no isomorfos con espacios de adyacencia isomorfos de la siguiente manera: sea $P$ el grafo sobre $V:=\{v_i\mid i\in\mathbb{N}\}$ dado por $E(P)=\{v_iv_{i+1}\mid i\in \mathbb{N}\}$.
   \begin{figure}[H]
     \centering
       \resizebox{0.45\textwidth}{!}{\begin{tikzpicture}
  [scale=.7,auto=left,every node/.style={circle,fill=blue!20,draw=black,thick}]
  \node (n1) at (0,0) {$v_0$};
  \node (n2) at (2,0)  {$v_1$};
  \node (n3) at (4,0)  {$v_2$};
  \node (n4) at (6,0) {$v_3$};
  \node (n5) at (8,0)  {$v_4$};

  \foreach \from/\to in {n1/n2,n2/n3, n3/n4, n4/n5}
    \draw (\from) -- (\to);

   %Flecha de (n5) a un nodo invisible en (10,0)   
   \draw (n5) -- (10,0);
   
   %Etiqueta de puntos suspensivos
   \node[draw=none, fill=none, scale=1.5] at (11,0) {$\cdots$};
   % Etiqueta G_1 debajo del grafo
%   \node[draw=none, fill=none, scale=1.5] at (6,-1.5) {$G_5$};
\end{tikzpicture}
} 
     \caption{El grafo $P$}
   \end{figure}
   Ahora, sea $V^*=\{v_0,y_1,y_2,\dots,y_{n-1}\}$ con $n\ge 5$, y tomemos $K^n$ y $C^n$ como el grafo completo sobre $V^*$ y el ciclo sobre $V^*$, respectivamente. Definimos entonces $G_1=\left( V\cup V^*, E(P)\cup E\left( K^n\right)\right)$ y $G_2=\left( V\cup V^*, E(P)\cup E\left( C^n\right) \right)$.
   \begin{figure}[H]
\centering
\resizebox{0.6\textwidth}{!}
{
\begin{tikzpicture}[
  scale=0.8,
  vnode/.style={circle, fill=blue!20, draw=black, thick, inner sep=2pt, minimum size=7mm},
  edge/.style={thick, draw=black}
]

  % --- Vértice superior v_1 (origen del camino y cúspide del K_5) ---
  \node[vnode] (v1) at (0, 0) {$v_0$};

  % --- Vértices restantes del K_5 (disposición en pentágono regular simétrico abajo de v_1) ---
  % Centro del pentágono en (0, -1.9), radio 1.8
  \node[vnode] (y2) at (-1.71, -1.24) {$y_{m}$};
  \node[vnode] (y3) at ( 1.71, -1.24) {$y_1$};
  \node[vnode] (y4) at ( 1.06, -3.51) {$y_2$};
  \node[vnode] (y5) at (-1.06, -3.51) {$y_3$};

  % --- Aristas del K_5 completo ---
  % Conexiones desde v_1 a todos los y_i
  \draw[edge] (v1) -- (y2);
  \draw[edge] (v1) -- (y3);
  \draw[edge] (v1) -- (y4);
  \draw[edge] (v1) -- (y5);
  
  % Conexiones entre los vértices y_i
  \draw[edge] (y2) -- (y3);
  \draw[edge] (y2) -- (y4);
  
  % Arista interrumpida entre y_n (y2) y y_4 (y5) con puntos suspensivos horizontales
  \draw[edge] (y2) -- (-1.46, -2.1);
  \draw[edge] (-1.31, -2.65) -- (y5);
  \node[draw=none, fill=none] at (-1.385, -2.375) {$\cdots$};

  \draw[edge] (y3) -- (y4);
  \draw[edge] (y3) -- (y5);
  \draw[edge] (y4) -- (y5);

  % --- Camino horizontal superior (v_2 a v_5 y puntos suspensivos) ---
  \node[vnode] (v2) at (2.0, 0) {$v_1$};
  \node[vnode] (v3) at (4.0, 0) {$v_2$};
  \node[vnode] (v4) at (6.0, 0) {$v_3$};
  \node[vnode] (v5) at (8.0, 0) {$v_4$};

  \draw[edge] (v1) -- (v2);
  \draw[edge] (v2) -- (v3);
  \draw[edge] (v3) -- (v4);
  \draw[edge] (v4) -- (v5);
   
  % Continuación con puntos suspensivos
  \draw[edge, line width=1.5pt] (v5) -- (9.5, 0);
  \node[draw=none, fill=none, scale=1.5] at (10.2, 0) {$\cdots$};

\end{tikzpicture}
}
\caption{El grafo $G_1$, con $m=n-1$}
\end{figure}
 
   \begin{figure}[H]
\centering
\resizebox{0.6\textwidth}{!}{
\begin{tikzpicture}[
  scale=0.8,
  vnode/.style={circle, fill=blue!20, draw=black, thick, inner sep=2pt, minimum size=7mm},
  edge/.style={thick, draw=black}
]

  % --- Vértice superior v_1 (origen del camino y del ciclo C^5) ---
  \node[vnode] (v1) at (0, 0) {$v_0$};

  % --- Vértices restantes del C^5 (disposición en pentágono regular simétrico abajo de v_1) ---
  \node[vnode] (y2) at (-1.71, -1.24) {$y_{m}$};
  \node[vnode] (y3) at ( 1.71, -1.24) {$y_1$};
  \node[vnode] (y4) at ( 1.06, -3.51) {$y_2$};
  \node[vnode] (y5) at (-1.06, -3.51) {$y_3$};

  % --- Aristas del ciclo C^5 con corte y puntos suspensivos horizontales ---
  \draw[edge] (v1) -- (y2);
  
  % Arista interrumpida entre y_n (y2) y y_4 (y5)
  \draw[edge] (y2) -- (-1.46, -2.1);
  \draw[edge] (-1.31, -2.65) -- (y5);
  \node[draw=none, fill=none] at (-1.385, -2.375) {$\cdots$};

  \draw[edge] (y5) -- (y4);
  \draw[edge] (y4) -- (y3);
  \draw[edge] (y3) -- (v1);

  % --- Camino horizontal superior (v_2 a v_5 y puntos suspensivos) ---
  \node[vnode] (v2) at (2.0, 0) {$v_2$};
  \node[vnode] (v3) at (4.0, 0) {$v_3$};
  \node[vnode] (v4) at (6.0, 0) {$v_4$};
  \node[vnode] (v5) at (8.0, 0) {$v_5$};

  \draw[edge] (v1) -- (v2);
  \draw[edge] (v2) -- (v3);
  \draw[edge] (v3) -- (v4);
  \draw[edge] (v4) -- (v5);
   
  % Continuación con puntos suspensivos
  \draw[edge, line width=1.5pt] (v5) -- (9.5, 0);
  \node[draw=none, fill=none, scale=1.5] at (10.2, 0) {$\cdots$};

\end{tikzpicture}}
\caption{El grafo $G_2$, con $m=n-1$}
\end{figure}

   Veamos que tanto $G_1$ como $G_2$ tienen topología de adyacencia discreta.
   \begin{itemize}
      \item[$G_1$:] Sea $i\in\{1,2,\dots,n-1\}$. Tenemos
         $$
         \begin{aligned}
            U_{y_i} &= \bigcap_{w\in A_{y_i}}A_w \\
                    &\phantom{=} \text{(\Cref{Prop:U_x_2})} \\
                    &= A_{v_0}\cap \left( \bigcap_{\substack{1\le j\le n-1 \\ j\ne i}}A_{y_j}\right)\\
                    &= A_{v_0}\cap \left( \bigcap_{\substack{1\le j\le n-1 \\ j\ne i}}\left(V^* \setminus \{y_j\}\right)\right)\\
                    &= A_{v_0} \cap \left( V^*\setminus \left( \bigcup_{\substack{1\le j\le n-1 \\ j\ne i}}\{y_j\}\right)\right)\\
                    &= A_{v_0}\cap \{v_0,y_i\}\\
                    &= \{v_1,y_1,y_2,\dots ,y_{n-1}\}\cap\{v_0,y_i\}\\
                    &= \{y_i\}.
         \end{aligned}
         $$
         Para cualquier $i\in\{2,3,\dots\}$ tenemos
         $$
         \begin{aligned}
            U_{v_i} &= \bigcap_{w\in A_{v_i}}A_w \\
                    &= A_{v_{i-1}}\cap A_{v_{i+1}} \\
                    &= \{v_{i-2},v_i\}\cap\{v_i,v_{i+2}\}\\
                    &= \{v_i\}.
         \end{aligned}
         $$
         Además
         $$
         \begin{aligned}
            U_{v_1} &= \bigcap_{w\in A_{v_1}}A_w \\
                    &= A_{v_0}\cap A_{v_2}\\
                    &= \{v_1,y_1,y_2,\dots ,y_{n-1}\}\cap\{v_1,v_3\}\\
                    &= \{v_1\},
         \end{aligned}
         $$
         y
         $$
         \begin{aligned}
            U_{v_0} &= \bigcap_{w\in A_{v_0}}A_w \\
                    &= \left( \bigcap_{i=1}^{n-1} A_{y_i}\right)\cap A_{v_1} \\
                    &= \left( \bigcap_{i=1}^{n-1} \left( V^*\setminus \{y_i\}\right)\right) \cap A_{v_1} \\
                    &= \left( V^*\setminus\bigcup_{i=1}^{n-1} \{y_i\}\right)\cap A_{v_1} \\
                    &= \{v_0\}\cap A_{v_1} \\
                    &= \{v_0\}\cap \{v_0,v_2\} \\
                    &= \{v_0\}.
         \end{aligned}
         $$
         Por lo tanto, de la \Cref{Prop:U_x=ximplicadiscreto} se tiene que $(V(G_1),\tau_{\mathcal{A},G_1})$ es un espacio discreto.
      \item[$G_2$:] De forma análoga a como se hizo para $G_1$ se tiene que $U_{v_i}=\{v_i\}$ para todo $i\in\{2,3,\dots\}$. Sea $i\in\{2,3,\dots ,n-2\}$. Denotemos $y_0=v_0$. Entonces:
         $$
         \begin{aligned}
            U_{y_i} &= \bigcap_{w\in A_{y_i}}A_w \\
                    &= A_{y_{i-1}}\cap A_{y_{i+1}} \\
                    &= \{y_{i-2},y_i\}\cap \{y_i, y_{(i+2)\bmod n}\} \\
                    &= \{y_i\},
         \end{aligned}
         $$
         pues como $n\ge 5$ y $2\le i\le n-2$ entonces $i-2\ne (i+2)\bmod n$ y $y_{i-2}\ne y_{(i+2)\bmod n}$. Además,
         $$
         \begin{aligned}
            U_{y_1} &= \bigcap_{w\in A_{y_1}}A_w \\
                    &= A_{v_0}\cap A_{y_2} \\
                    &= \{v_1,y_1,y_{n-1}\} \cap \{y_1,y_3\} \\
                    &= \{y_1\},
         \end{aligned}
         $$
         pues $n-1\ge 4$, luego $y_{n-1}\ne y_3$. También,
         $$
         \begin{aligned}
            u_{y_{n-1}} &= \bigcap_{w\in A_{y_{n-1}}}A_w \\
                        &= A_{v_0}\cap A_{y_{n-2}} \\
                        &= \{v_1,y_1,y_{n-1}\}\cap \{y_{n-1},y_{n-3}\} \\
                        &= \{y_{n-1}\},
         \end{aligned}
         $$
         pues $n-3\ge 2$, luego $y_{n-3}\ne y_1$. Por último,
         $$
         \begin{aligned}
            U_{v_0} &= \bigcap_{w\in A_{v_0}}A_w \\
                    &= A_{v_1}\cap A_{y_1}\cap A_{y_{n-1}} \\
                    &= \{v_0, v_2\}\cap \{v_0,y_2\}\cap \{v_0,y_{n-1}\} \\
                    &= \{v_0\},
         \end{aligned}
         $$
         y
         $$
         \begin{aligned}
            U_{v_1} &= \bigcap_{w\in A_{v_1}}A_w \\
                    &= A_{v_0}\cap A_{v_2} \\
                    &= \{v_1,y_1,y_{n-1}\}\cap \{v_1,v_3\} \\
                    &= \{v_1\}.
         \end{aligned}
         $$
         Por lo tanto, de la \Cref{Prop:Alexandroff_discreto} se tiene que $(V(G_2),\tau_{\mathcal{A},G_2})$ es un espacio discreto.
   \end{itemize}
   Como $(V(G_1),\tau_{\mathcal{A},G_1})$ y $(V(G_2),\tau_{\mathcal{A},G_2})$ son espacios discretos y $|V(G_1)|=|V(G_2)|$, entonces por el \Cref{Ejm:homeomorfismos} se tiene que  $(V(G_1),\tau_{\mathcal{A},G_1})$ y $(V(G_2),\tau_{\mathcal{A},G_2})$ son homeomorfos. Sin embargo $G_1$ tiene un vértice de grado $n$, mientras que $G_2$ no lo tiene, así que por el \Cref{Cor:Isomorfismo_grados_iguales} se tiene que $G_1$ y $G_2$ no son isomorfos.
\end{Obs}
\begin{Prop}
   Sea $G=(V,E)$ un grafo. Entonces $(V,\tau_{\mathcal{A},G})$ es compacto si, y solo si, $V$ es finito. 
\end{Prop}
\begin{proof}
   ($\Rightarrow$) Probamos la contrarrecíproca. Supongamos que $V$ es infinito. Tenemos que $\mathcal{B}_{\operatorname{min}}=\left\lbrace U_x\mid x\in V\right\rbrace$ es un cubrimiento abierto de $V$. Para cualquier subconjunto finito $\left( U_{x_i}\right)_{i=1}^{n}$ de $\mathcal{B}_{\operatorname {min}}$, como, por la \Cref{Prop:U_x_2}, cada $U_{x_i}$ es finito, entonces $\bigcup_{i=1}^{n}U_{x_i}$ es finito y por tanto $\bigcup_{i=1}^{n}U_{x_i}\ne V$. Con esto, $\mathcal{B}_{\operatorname{min}}$ no tiene ningún subcubrimiento finito, y por tanto $(V,\tau_{\mathcal{A},G})$ no es compacto.\\
   ($\Leftarrow$) Sabemos que cualquier espacio topológico sobre un conjunto finito es compacto.
\end{proof}
\begin{Prop}
   Sean $G=(V,E)$ un grafo y $M:=\{x\in V\mid d(v)=\Delta (G)\}$. Entonces $M\in\tau_{\mathcal{A},G}$. 
\end{Prop}
\begin{proof}
   Usamos la \Cref{Prop:Caracterizacion_abiertos}. Sea $x\in M$. Veamos que $U_x\subseteq M$. Sea $y\in U_x$. Por la \Cref{Prop:U_x_2} se tiene que $y\in A_w$ para todo $w\in A_x$. De esta forma, para cualquier $w\in A_x$ se tiene $y\in A_w$ y por tanto $w\in A_y$, es decir, $A_x\subseteq A_y$, con lo cual $\Delta(G)=d(x)=|A_x|\le |A_y|=d(y)\le \Delta(G)$ y por tanto $d(y)=\Delta (G)$. Con esto, $y\in M$ y $U_x\subseteq M$. Obtenemos que $U_x$ es una vecindad de $x$ contenida en $M$. Por tanto $M\in\tau_{\mathcal{A},G}$. 
\end{proof}
\begin{Prop}
   Sean $G=(V,E)$ un grafo y $L:=\{x\in V\mid d(v)=\delta(G)\}$. Entonces $L$ es cerrado en $(V,\tau_{\mathcal{A},G})$. 
\end{Prop}
\begin{proof}
   Por la \Cref{Prop:Alexandroff_clausura} se tiene $\overline{L}=\overline{\bigcup_{x\in L}\{x\}}=\bigcup_{x\in L}\overline{\{x\}}$. Sea $y\in \overline{L}$. Entonces existe $x\in L$ tal que $y\in\overline{\{x\}}$. Por el \Cref{Cor:Clausura_unitario} se tiene $A_y\subseteq A_x$. Con esto, $\delta (G)\le d(y)=|A_y|\le |A_x|=\delta (G)$ y por tanto $d(y)=\delta(G)$, es decir, $y\in L$. Obtenemos que $\overline{L}\subseteq L$. Como además $L\subseteq\overline{L}$, entonces $L=\overline{L}$, y por la \Cref{Prop:Caracterizacion_abierto_cerrado} se sigue que $L$ es cerrado en $(V,\tau_{\mathcal{A},G})$.
\end{proof}
\begin{Lema}\label{Lema:x_x}
   Sean $G$ un grafo y $x\in V(G)$ un vértice de corte de $G$. Entonces existen $y_1,y_2\in A_x$ y $D_1,D_2$ componentes distintas de $G-\{x\}$ tales que $y_1\in V(D_1)$ y $y_2\in V(D_2)$.
\end{Lema}
\begin{proof}
   Llamemos $\mathcal{C}$ a la colección de componentes de $G$ y $\mathcal{D}$ a la colección de componentes de $G-\{x\}$. Como $G$ no tiene vértices asilados, existe $y_1\in A_x$. Por la definición de grafo simple tenemos $y_1\ne x$, es decir, $y_1\in V(G)\setminus \{x\}=V(G-\{x\})$. Por la \Cref{Prop:Componentes_particion}, existe $D_1\in\mathcal{D}$ tal que $y_1\in V(D_1)$. Razonando por contradicción supongamos que $A_x\subseteq V(D_1)$.\\\\
   Veamos que $\mathcal{D}=\{C\in \mathcal{C}\mid x\notin V(C)\}\cup \{D_1\}$.
   \begin{itemize}
      \item[$\subseteq$:] Sea $D\in\mathcal{D}$. Consideremos dos casos.
         \begin{itemize}
            \item[\textbf{I.}] Supongamos que $A_x\cap V(D) \ne\emptyset$. Sea $w\in A_x\cap V(D)$. Como $A_x\subseteq V(D_1)$ entonces $w\in V(D_1)$. Además $w\in V(D)$. De esta forma, $V(D)\cap V(D_1)\ne\emptyset$. Como $D,D_1\in\mathcal{D}$, por la \Cref{Prop:Componentes_particion} se sigue que $D=D_1$.   
            \item[\textbf{II.}] Supongamos que $A_x\cap V(D) = \emptyset$. Como $D$ es una componente de $G-\{x\}$, se tiene que $x\notin V(D)$ y que $D$ es conexo. Como $D\subseteq G-\{x\} \subseteq G$, tenemos que $D$ es un subgrafo de $G$. Sea $H$ un subgrafo conexo de $G$ tal que $D\subseteq H$. Veamos que $D=H$. Razonando por contradicción supongamos $V(H)\nsubseteq V(D)$. Tomemos $w\in V(H)\setminus V(D)$. Como las componentes de un grafo son no vacías, podemos tomar $d\in V(D)$. Tenemos que $d,w\in V(H)$ y $H$ es conexo, luego existe un camino $p_0\dots p_n$ contenido en $H$ con $n\ge 1$, $p_0=d$ y $p_n=w$. Sea $j$ el menor índice en $\{1,\dots , n\}$ tal que $p_j\notin V(D)$. Como $p_{j-1}\in V(D)$, por la \Cref{Prop:Conexo+1} se tiene que $H':=(V(D)\cup \{p_j\},E(D)\cup \{p_{j-1}p_j\})$ es un grafo conexo con $D\subsetneq H'$. Notemos además que si $p_j=x$ entonces $p_{j-1}\in A_x\cap V(D)$, contradiciendo $A_x\cap V(D)=\emptyset$. Por tanto $p_j\ne x$ y se sigue que $H'\subseteq G-\{x\}$. Pero entonces se contradice que $D$ sea un subgrafo conexo-maximal de $G-\{x\}$. Por lo tanto $V(H)\subseteq V(D)$ y $V(D)=V(H)$. Ahora, sea $e\in E(H)$. Entonces $e=st$ para algunos $s,t\in V(H)$. Como $V(H)=V(D)\subseteq V(G-\{x\})=V(G)\setminus\{x\}$ entonces $s,t\in V(G)\setminus\{x\}$. Además $H\subseteq G$, así que $E(H)\subseteq E(G)$ y $st\in E(G)$. De esta forma $e\in\{uv\in E(G)\mid u,v\in V(G)\setminus \{x\}\}=E(G-\{x\})$. Con esto, $E(H)\subseteq E(G-\{x\})$ y se tiene $H\subseteq G-\{x\}$, con $H$ conexo y $D\subseteq H$. Como $D$ es un subgrafo conexo-maximal de $G-\{x\}$ entonces $D=H$. Por lo tanto $D$ es un subgrafo conexo-maximal de $G$, es decir $D\in\mathcal{C}$, y además tenemos $x\notin V(D)$.
         \end{itemize}
         En cualquier caso, $D\in \{C\in \mathcal{C}\mid x\notin V(C)\}\cup \{D_1\}$ y por tanto $\mathcal{D}\subseteq \{C\in \mathcal{C}\mid x\notin V(C)\}\cup \{D_1\}$.
      \item[$\supseteq$:] Sabemos que $D_1\in \mathcal{D}$. Sea $C\in\mathcal{C}$ con $x\notin V(C)$. Como $C\subseteq G$ tenemos $V(C)\subseteq V(G)$; además $x\notin V(C)$, luego $V(C)\subseteq V(G)\setminus \{x\}=V(G-\{x\})$. Sea $e\in E(C)$. Entonces $e=st$ para algunos $s,t\in V(C)$. En particular $s,t\in V(G)\setminus \{x\}$. Como $C\subseteq G$ entonces $E(C)\subseteq E(G)$ y $st\in E(G)$. De esta forma $e\in \{uv\in E(G)\mid u,v\in V(G)\setminus \{x\}\}=E(G-\{x\})$ y $E(C)\subseteq E(G-\{x\})$. Por tanto $C\subseteq G-\{x\}$. Además sabemos que $C$ es conexo. Supongamos que $H$ es un subgrafo conexo de $G-\{x\}$ tal que $C\subseteq H$. Como $G-\{x\}\subseteq G$, se tiene $H\subseteq G$. Como $C$ es un subgrafo conexo-maximal de $G$ se tiene $C=H$. Por lo tanto $C$ es un subgrafo conexo-maximal de $G-\{x\}$, es decir $C\in\mathcal{D}$. Obtenemos $\{C\in \mathcal{C}\mid x\notin V(C)\}\cup \{D_1\}\subseteq\mathcal{D}$.
   \end{itemize}
   De lo anterior, $\mathcal{D}=\{C\in \mathcal{C}\mid x\notin V(C)\}\cup \{D_1\}$. Definimos entonces
   \begin{equation*}
     \hspace*{6em}
     \begin{array}{r c c c l}
       f \colon
         & \mathcal{C} & \longrightarrow & \mathcal{D} & \\[0.5em]
       \hphantom{{}=\begin{cases} C, & \text{si } x \notin C, \\ C_1, & \text{si } x \in C. \end{cases}}
         & C & \longmapsto & f(C) &
       {}=\begin{cases}
         C,   & \text{si } x \notin V(C), \\
         D_1, & \text{si } x \in V(C).
       \end{cases}
     \end{array}
   \end{equation*}
   Veamos que $f$ es inyectiva. Sean $C_1,C_2\in\mathcal{C}$ cualesquiera tales que $f(C_1)=f(C_2)$. Si $x\notin V(C_1)$ y $x\notin V(C_2)$ entonces $f(C_1)=C_1$ y $f(C_2)=C_2$, por tanto $C_1=C_2$. Si $x\in V(C_1)$ y $x\notin V(C_2)$ entonces $C_1\ne C_2$ y además $f(C_1)=D_1$ y $f(C_2)=C_2$, por tanto $D_1=C_2$. Como $y_1\in V(D_1)$ entonces $y_1\in V(C_2)$. Pero $y_1\in A_x$ con $x\in V(C_1)$, contradiciendo la \Cref{Obs:Componentes_no_adyacentes}. De forma similar, si $x\notin V(C_1)$ y $x\in V(C_2)$ se llega a una contradicción. Si $x\in V(C_1)$ y $x\in V(C_2)$ entonces por la \Cref{Prop:Componentes_particion} se tiene $C_1=C_2$. De los casos anteriores se sigue que $C_1=C_2$ y por tanto $f$ es inyectiva.\\\\
   Ahora, para cualquier $C\in\mathcal{C}$ con $x\notin V(C)$ tenemos $f(C)=C$, y por la \Cref{Prop:Componentes_particion} existe $C_x\in \mathcal{C}$ con $x\in V(C_x)$, de modo que $f(C_x)=D_1$. Por tanto $f$ es sobreyectiva.\\\\
   De lo anterior, tenemos que $f$ es una biyección entre $\mathcal{C}$ y $\mathcal{D}$, luego $|\mathcal{C}|=|\mathcal{D}|$, es decir que $G$ tiene el mismo número de componentes de $G-\{x\}$, con lo cual se contradice que $x$ sea un vértice de corte de $G$. Se sigue por tanto que $A_x\nsubseteq V(D_1)$. Sea $y_2\in A_x\setminus V(D_1)$. Como $y_2\in A_x$, por la definición de grafo simple se tiene $y_2\ne x$, es decir $y_2\in V(G)\setminus\{x\}=V(G-\{x\})$. Por la \Cref{Prop:Componentes_particion}, existe $D_2\in\mathcal{D}$ tal que $y_2\in V(D_2)$. Como $y_2\notin V(D_1)$ se sigue que $D_1\ne D_2$, lo cual completa la prueba.
\end{proof}
\begin{Prop}
   Sean $G=(V,E)$ un grafo y $x\in V$. Si $x$ es un vértice de corte de $G$ entonces $\{x\}\in\tau_{\mathcal{A},G}$.
\end{Prop}
\begin{proof}
   Supongamos que $x$ es un vértice de corte de $G$. Por el \Cref{Lema:x_x}, existen $y_1,y_2\in A_x$ y $D_1,D_2$ componentes distintas de $G-\{x\}$ tales que $y_1\in V(D_1)$ y $y_2\in V(D_2)$. Sea $z\in A_{y_1}\cap A_{y_2}$. Si $z\ne x$ entonces $z\in V(G)\setminus \{x\}=V(G-\{x\})$, de modo que $y_1zy_2$ es un camino en $G-\{x\}$ que conecta $y_1$ y $y_2$, contradiciendo la \Cref{Prop:Componentes_no_conectan}. Por tanto $z=x$ y $A_{y_1}\cap A_{y_2}\subseteq \{x\}$. De esta manera, usando la \Cref{Prop:U_x_2} tenemos
   $$
   \{x\}\subseteq U_x=\bigcap_{w\in A_x}A_w\subseteq A_{y_1}\cap A_{y_2}\subseteq \{x\},
   $$
   con lo cual $\{x\}=U_x\in \tau_{\mathcal{A},G}$.
\end{proof}

\chapter{Topología de incidencia}\label{chap:topologia_de_incidencia}
En este capítulo no exigimos que los grafos en cuestión sean localmente finitos. Sin embargo mantenemos la condición de que los grafos no tengan vértices aislados.
\begin{Prop}
   Sea $G=(V,E)$ un grafo. El conjunto $E$ es subbase para una topología sobre $V$. 
\end{Prop}
En \cite{kilicman2018} se define primero, para cada arista del grafo $G$, el conjunto $I_e$ de todos los vértices ``incidentes'' con $e$, para luego tomar como subbase el conjunto $S_{IG}=\{I_e\mid e\in E\}$. Sin embargo, con las definiciones que desarrollamos en el \Cref{chap:Preliminares}, cada conjunto $I_e$ concide exactamente con $e$, y por tanto el conjunto $S_{IG}$ es igual al conjunto de aristas del grafo, $E$.
\begin{Def}
   Sea $G=(V,E)$ un grafo. La topología generada por el conjunto $E$, denotada por $\tau_{\mathcal{I},G}$, es llamada la \textbf{topología de incidencia} de $G$. El espacio topológico $(V,\tau_{\mathcal{I},G})$ es llamado el \textbf{espacio de incidencia} de $G$. 
\end{Def}
\begin{Prop}
   Sea $G=(V,E)$ un grafo. Se tiene que $(V,\tau_{\mathcal{I},G}))$ es un espacio de Alexandroff. 
\end{Prop}
\begin{Prop}
   Sea $G=(V,E)$ un grafo. Para toda $x\in V$ tenemos $U_x=\bigcap_{v\in e}e$.
\end{Prop}
\begin{Prop}
   Sean $G=(V,E)$ un grafo y $x\in V$. Si $\operatorname{deg}(v)\ge 2$ entonces $U_x=\{x\}$ 
\end{Prop}

\chapter{Implementación computacional}\label{chap:implementacion_computacional}
En este capítulo describimos el desarrollo de diversos códigos, escritos en el lenguaje de programación \Cpp{}, para facilitar el cálculo explícito de la topología de adyacencia y la topología de incidencia de un grafo finito, calculando, en cada caso, los abiertos subbásicos, los abiertos básicos y la colección de todos los abiertos de la topología en cuestión, recibiendo como entrada el número de vértices y la matriz de adyacencia del grafo que se desea estudiar. 
\section{Funciones preliminares}\label{sec:funciones_preliminares}
En esta sección describimos las funciones que se usarán conjuntamente para el cálculo de la topología de adyacencia y de la topología de incidencia. El objetivo grande que se desea con el grueso de funciones es lograr realizar intersecciones y uniones de conjuntos. Para esto nos apoyaremos fuertemente en el uso de matrices y vectores, y de las funcionalidades básicas que ofrece \Cpp{} para la manipulación de éstas.
Para la implementación de las funciones usamos cuatro cabeceras de la biblioteca estándar de \Cpp{}.
\begin{lstlisting}[caption={Encabezado del código}, label={cod:encabezado}]
#include <iostream>
#include <fstream>
#include <vector>
#include <cmath>
\end{lstlisting}
La cabecera \texttt{<iostream>} no es necesaria para imprimir texto en la terminal. La cabecera \texttt{<fstream>} la incluimos para poder leer archivos del tipo \texttt{.txt}, donde guardaremos el número de vértices y la matriz de adyacencia del grafo en cuestión. La cabecera \texttt{<vector>} nos permite la creación y manipulación de vectores y matrices. Por último, la cabecera \texttt{<cmath>} la implementamos para usar la función \texttt{pow}, con la cual realizamos la operación potencia de números enteros.\\\\
Para el funcionamiento del código usaremos archivos \texttt{.txt} para almacenar la información necesaria del grafo al que se le calcularán las topologías da adyacencia y de incidencia. Esta se conforma únicamente por un número entero que representa el número de vértices del grafo en cuestión, y un arreglo rectangular de ceros y unos, separados por \texttt{tabs} y renglones, que representa la matriz de adyacencia del mismo. Mostramos a continuación un par de ejemplos.
\begin{figure}[H]
   \centering
   % Fila 1: el grafo y el archivo, centrados verticalmente
   \begin{minipage}[c]{0.40\textwidth}
      \centering
      \resizebox{0.78\linewidth}{!}{\input{Implementacion_computacional/Diagramas/Diag1.tex}}
   \end{minipage}%
   \hspace{1.5em}%
   \begin{minipage}[c]{0.42\textwidth}
\begin{lstlisting}[framexrightmargin=-4.4pt]
5
0  1  1  1  0
1  0  1  1  0
1  1  0  1  1
1  1  1  0  0
0  0  1  0  0
\end{lstlisting}
   \end{minipage}

   % Fila 2: los pies, alineados arriba y centrados
   \captionsetup{format=plain, justification=centering}
   \begin{minipage}[t]{0.40\textwidth}
      \captionof{figure}{Grafo $G_1$}
      \label{fig:G_1}
   \end{minipage}%
   \hspace{1.5em}%
   \begin{minipage}[t]{0.42\textwidth}
      \captionof{lstlisting}{Archivo \texttt{1\_Mat\_Ady.txt} para guardar la información de $G_1$}
      \label{cod:1_Mat_Ady}
   \end{minipage}
\end{figure}

\begin{figure}[H]
   \centering
   % Fila 1: el grafo y el archivo, centrados verticalmente
   \begin{minipage}[c]{0.40\textwidth}
      \centering
      \resizebox{0.78\linewidth}{!}{\input{Implementacion_computacional/Diagramas/Diag2.tex}}
   \end{minipage}%
   \hspace{1.5em}%
   \begin{minipage}[c]{0.42\textwidth}
\begin{lstlisting}[framexrightmargin=-4.4pt]
8 (*@\label{lin:r1_2_Mat_Ady}@*)
0 0 0 1 1 1 1 1 (*@\label{lin:inicio_2_Mat_Ady}@*)
0 0 0 1 1 1 1 1
0 0 0 1 1 1 1 1
1 1 1 0 0 0 0 0
1 1 1 0 0 0 0 0
1 1 1 0 0 0 0 0
1 1 1 0 0 0 0 0
1 1 1 0 0 0 0 0 (*@\label{lin:fin_2_Mat_Ady}@*)
\end{lstlisting}
\end{minipage}

   % Fila 2: los pies, alineados arriba y centrados
   \captionsetup{format=plain, justification=centering}
   \begin{minipage}[t]{0.40\textwidth}
      \captionof{figure}{Grafo $K_{3,5}$}
      \label{fig:K_3,5}
   \end{minipage}%
   \hspace{1.5em}%
   \begin{minipage}[t]{0.42\textwidth}
      \captionof{lstlisting}{Archivo \texttt{2\_Mat\_Ady.txt} para guardar la información de $K_{3,5}$}
      \label{cod:2_Mat_Ady}
   \end{minipage}
\end{figure}

Así, por ejemplo, en la \Cref{lin:r1_1_Mat_Ady} del \Cref{cod:1_Mat_Ady}, el número $5$ hace referencia al número de vectores del grafo $G_1$, mientras que la \Cref{lin:r1_2_Mat_Ady} del \Cref{cod:2_Mat_Ady} hace referencia al número de vértices del grafo $K_{3,5}$. Igualmente, en las \Crefrange{lin:inicio_1_Mat_Ady}{lin:fin_1_Mat_Ady} del \Cref{cod:1_Mat_Ady}, se representa la matriz de adyacencia
\begin{ceqn}
\[
\begin{pmatrix}
   0 & 1 & 1 & 1 & 0 \\
   1 & 0 & 1 & 1 & 0 \\
   1 & 1 & 0 & 1 & 1 \\
   1 & 1 & 1 & 0 & 0 \\
   0 & 0 & 1 & 0 & 0
\end{pmatrix}
\]
\end{ceqn}
del grafo $G_1$. Análogamente, en las líneas \Crefrange{lin:inicio_2_Mat_Ady}{lin:fin_2_Mat_Ady} del \Cref{cod:2_Mat_Ady}, se representa la matriz de adyacencia
\begin{ceqn}
\[
\begin{pmatrix}
   0 & 0 & 0 & 1 & 1 & 1 & 1 & 1 \\
   0 & 0 & 0 & 1 & 1 & 1 & 1 & 1 \\
   0 & 0 & 0 & 1 & 1 & 1 & 1 & 1 \\
   1 & 1 & 1 & 0 & 0 & 0 & 0 & 0 \\
   1 & 1 & 1 & 0 & 0 & 0 & 0 & 0 \\
   1 & 1 & 1 & 0 & 0 & 0 & 0 & 0 \\
   1 & 1 & 1 & 0 & 0 & 0 & 0 & 0 \\
   1 & 1 & 1 & 0 & 0 & 0 & 0 & 0
\end{pmatrix}
\]
\end{ceqn}
del grafo $K_{3,5}$.
\begin{lstlisting}[caption={Función \lstinline|leerMatriz|: lectura de la matriz de adyacencia desde un archivo}, label={cod:leerMatriz}]
std::vector<std::vector<bool>> leerMatriz(std::ifstream& file){
   int n;
   file >> n;
   std::vector<std::vector<bool>> Mat_Ady(n, std::vector<bool>(n));
   for(int i=0; i<=n-1; i++){
      for(int j=0; j<=n-1; j++){
         bool valor;
         file >> valor;
         Mat_Ady[i][j] = valor;
      }
   }
   return Mat_Ady;
}\end{lstlisting}

\section{Topología de adyacencia}\label{sec:adyacencia_computacional}
\section{Topología de incidencia}\label{sec:incidencia_computacional}

% ==============================================================================
% BIBLIOGRAFÍA
% ==============================================================================
\clearpage
\phantomsection
\addcontentsline{toc}{chapter}{\bibname}
\nocite{*}
\bibliographystyle{apalike}
\bibliography{refs.bib}

\end{document}
